% Lesson 13 — Speaking: Exchanges information about buying and selling articles/foodstuff at a nearby store/market — Anglais 1ère A — Cours signé OnBuch+
% Programme officiel MINESEC. Compiler avec : tectonic ENA14-speaking-exchanges-information-about-buying-and-selling.tex
\newcommand{\DOCMATIERE}{Anglais}
\newcommand{\DOCNIVEAU}{1ère A}
\newcommand{\DOCMODULE}{Chapter 2 : Economic Life and Occupations}
\newcommand{\DOCLECON}{ENA14}
\newcommand{\DOCTITRE}{Lesson 13 — Speaking: Exchanges information about buying and selling articles/foodstuff at a nearby store/market}
% OnBuch+ — préambule « cours signés » ENRICHI (compilé avec tectonic / XeLaTeX)
% Système visuel « Pop » d'OnBuch+ : fond crème (jamais blanc), bordures encre
% épaisses, ombres « dures » décalées, titres Archivo Black en capitales.
% Version MATHÉMATIQUES : graphes (pgfplots/tikz), boîtes pédagogiques
% (définition, propriété, méthode, attention, le savais-tu, exercice résolu,
% exercice, TP, bilan), page de garde et signature OnBuch+ ancrée en bas.
%
% Macros attendues dans main.tex AVANT \input{preamble} :
%   \DOCMATIERE  \DOCNIVEAU  \DOCMODULE  \DOCLECON (numéro)  \DOCTITRE
\documentclass[11pt]{article}

\usepackage{fontspec}
\usepackage[english,french]{babel} % français = langue principale ; l'anglais via \eng{..} / textebox
\usepackage[a4paper,margin=2cm,top=3.4cm,bottom=2.8cm,headheight=1.4cm,footskip=1.3cm]{geometry}
\usepackage{amsmath,amssymb,amsfonts}
\usepackage{mathtools} % \xmapsto, \coloneqq... souvent utilisés par les modèles
\usepackage{cancel} % simplifications barrées (bilans de forces)
% \abs{z} (module, valeur absolue) et \norm{u} (norme), utilisés par les modèles
\providecommand{\abs}{}\let\abs\relax\DeclarePairedDelimiter{\abs}{\lvert}{\rvert}
\providecommand{\norm}{}\let\norm\relax\DeclarePairedDelimiter{\norm}{\lVert}{\rVert}
\usepackage[table]{xcolor} % [table] : fournit aussi \rowcolors (alternance de lignes)
\usepackage{pagecolor}
\usepackage{tikz}
\usetikzlibrary{arrows.meta,positioning,calc,shapes.geometric,shapes.misc,shapes.arrows,decorations.pathmorphing,decorations.pathreplacing,decorations.markings,patterns,shadows,mindmap,trees,fit,backgrounds,matrix,intersections,angles,quotes}
\usepackage{pgfplots}
\usepackage[version=4]{mhchem} % équations nucléaires \ce{^{235}_{92}U}
\usepackage[european,siunitx]{circuitikz} % schémas électriques
\pgfplotsset{compat=1.18}
\usepgfplotslibrary{fillbetween,groupplots} % aires entre courbes (calcul intégral)
\usepackage{enumitem}
\usepackage{fancyhdr}
% [most] charge toutes les bibliothèques tcolorbox (listings, minted, raster,
% documentation...) et gonfle inutilement la mémoire TeX sur un document long
% (~300 boîtes) : on ne charge que ce qui est réellement utilisé.
\usepackage[skins,breakable]{tcolorbox}
\usepackage{graphicx}
\usepackage{colortbl}
\usepackage{booktabs}
\usepackage{tabularx}
\usepackage{diagbox} % case d'en-tête diagonale des tableaux à double entrée
% Colonnes extensibles centrée / gauche / droite (C, L, R), souvent utilisées par les modèles
\newcolumntype{C}{>{\centering\arraybackslash}X}
\newcolumntype{L}{>{\raggedright\arraybackslash}X}
\newcolumntype{R}{>{\raggedleft\arraybackslash}X}
\usepackage{array}
\usepackage{multicol}
\usepackage{multirow}
\usepackage{siunitx}
% parse-numbers=false : accepte aussi \SI{2,5 \times 10^{-3}}{...}
\sisetup{locale=FR,output-decimal-marker={,},group-minimum-digits=4,parse-numbers=false}
\DeclareSIUnit\ohm{\ensuremath{\symup{\Omega}}} % U+2126 absent de la police
\DeclareSIUnit\division{div} % réglages d'oscilloscope (ms/div, V/div)
\DeclareSIUnit\tr{tr} % tours (vitesse de rotation, ex. \SI{1500}{\tr\per\minute})
\DeclareSIUnit\molar{M} % concentration molaire (ex. \SI{3}{\molar}, \SI{1}{\micro\molar})
\DeclareSIUnit\an{an} % année (le modèle confond parfois avec \year, un compteur TeX) — ex. \SI{5}{\kilo\meter\per\an}
\DeclareSIUnit\FCFA{FCFA} % franc CFA (ex. \SI{500}{\FCFA\per\kilo\gram})
\DeclareSIUnit\degreeBrix{\textdegree Brix} % teneur en sucre d'un sirop/jus (réfractométrie)
\DeclareSIUnit\month{mois} % le modèle utilise parfois \month, un compteur TeX, comme unité de durée
\usepackage{qrcode}
% \degree : commande fréquemment utilisée par les modèles pour un angle ou une
% température en dehors de \SI{}{} (non fournie par les paquets chargés).
\providecommand{\degree}{^{\circ}}
% \qed (fin de démonstration), utilisé par les modèles sans amsthm
\providecommand{\qed}{\hfill$\square$}
% \barre : double barre des valeurs interdites dans un tableau de variations (tkz-tab)
\providecommand{\barre}{\ensuremath{\|}}
% environnement proof (amsthm non chargé) utilisé par les modèles
\ifdefined\proof\else\newenvironment{proof}[1][Démonstration]{\par\noindent\textit{#1.}\ }{\qed\par}\fi
\usepackage{lastpage}
\usepackage{hyperref}
\hypersetup{hidelinks}

% ============================================================
% Palette « Pop » OnBuch+ — mêmes hex que la classe P (plus_ui.dart)
% ============================================================
\definecolor{poporange}{HTML}{F59321}
\definecolor{popdark}{HTML}{E07A0C}
\definecolor{popcream}{HTML}{FFF6E8}
\definecolor{popink}{HTML}{1C1714}
\definecolor{poppurple}{HTML}{7A5AE0}
\definecolor{popgreen}{HTML}{2E9E6A}
\definecolor{poppink}{HTML}{E0457B}
\definecolor{popblue}{HTML}{2F7DE1}
\definecolor{popgold}{HTML}{C98A12}
\definecolor{popmuted}{HTML}{8A7F76}
\definecolor{popline}{HTML}{EADFD2}
\definecolor{popgray}{HTML}{8A7F76}
\colorlet{popgrey}{popgray}
% Alias tolérants (le modèle nomme parfois les couleurs différemment)
\colorlet{popwhite}{white}
\colorlet{popblack}{popink}
\colorlet{popred}{poppink}
\colorlet{popyellow}{popgold}
% Teintes claires (fonds de tableaux/encadrés) — le modèle nomme parfois
% les couleurs avec un suffixe L (ex. popblueL) sans les avoir définies.
\colorlet{poporangeL}{poporange!20}
\colorlet{popdarkL}{popdark!20}
\colorlet{poppurpleL}{poppurple!20}
\colorlet{popgreenL}{popgreen!20}
\colorlet{poppinkL}{poppink!20}
\colorlet{popblueL}{popblue!20}
\colorlet{popgoldL}{popgold!20}
\colorlet{popredL}{popred!20}
\colorlet{popgrayL}{popgray!20}
% Teintes claires pour fonds de tableaux / zones
\colorlet{popblueL}{popblue!10!white}
\colorlet{poporangeL}{poporange!14!white}
\colorlet{popgreenL}{popgreen!12!white}
\colorlet{poppurpleL}{poppurple!12!white}
\colorlet{poppinkL}{poppink!12!white}
\colorlet{popgoldL}{popgold!14!white}

\pagecolor{popcream}

% ============================================================
% Typographie
% ============================================================
\defaultfontfeatures{Ligatures=TeX}
\setmainfont{PlusJakartaSans-Medium.ttf}[Path=../../../fonts/,BoldFont=PlusJakartaSans-Bold.ttf]
% Symboles, API et emojis absents de la police principale : repli sur DejaVu Sans (caractères actifs)
\newfontface\symfont{DejaVuSans.ttf}[Path=../../../fonts/]
\makeatletter
\newcommand\symactive[1]{\catcode#1=13 \begingroup\lccode`\~=#1 \lowercase{\endgroup\def~}{{\symfont\char#1}}}
\newcommand\symblank[1]{\catcode#1=13 \begingroup\lccode`\~=#1 \lowercase{\endgroup\def~}{}}
\newcommand\symhyph[1]{\catcode#1=13 \begingroup\lccode`\~=#1 \lowercase{\endgroup\def~}{\mbox{-}}}
\newcount\symc
\newcommand\symrange[2]{\symc=#1 \@whilenum\symc<#2 \do{\expandafter\symactive\expandafter{\the\symc}\advance\symc by 1 }}
\newcommand\symblankrange[2]{\symc=#1 \@whilenum\symc<#2 \do{\expandafter\symblank\expandafter{\the\symc}\advance\symc by 1 }}
\makeatother
\symrange{"0250}{"02FF}\symactive{"2713}\symactive{"2714}\symactive{"2717}\symactive{"2718}\symactive{"2610}\symactive{"2611}\symactive{"2612}
\symhyph{"2011}\symhyph{"2E1A}\symblankrange{"1F300}{"1FAFF}\symblank{"FE0F}
% Maths en sans-serif (Fira Math) avec chiffres et lettres droites de Plus
% Jakarta, pour que les formules se fondent dans le texte.
\usepackage[math-style=ISO,bold-style=ISO]{unicode-math}
\setmathfont{FiraMath-Regular.otf}[Path=../../../fonts/]
\setmathfont{PlusJakartaSans-Medium.ttf}[Path=../../../fonts/,range={up/num,up/{Latin,latin},\mathup}]
\usepackage{icomma} % virgule décimale sans espace en mode math
% Caractères mathématiques Unicode absents des polices de texte (Plus Jakarta,
% Archivo Black) : rendus par leur équivalent mathématique, en texte comme en
% formule — sinon ils s'affichent en carré vide (ex. « Arithmétique dans ℤ »).
\usepackage{newunicodechar}
\newunicodechar{ℝ}{\ensuremath{\mathbb{R}}}
\newunicodechar{ℤ}{\ensuremath{\mathbb{Z}}}
\newunicodechar{ℂ}{\ensuremath{\mathbb{C}}}
\newunicodechar{ℕ}{\ensuremath{\mathbb{N}}}
\newunicodechar{ℚ}{\ensuremath{\mathbb{Q}}}
\newunicodechar{𝔻}{\ensuremath{\mathbb{D}}}
\newunicodechar{α}{\ensuremath{\alpha}}
\newunicodechar{β}{\ensuremath{\beta}}
\newunicodechar{γ}{\ensuremath{\gamma}}
\newunicodechar{δ}{\ensuremath{\delta}}
\newunicodechar{ε}{\ensuremath{\varepsilon}}
\newunicodechar{θ}{\ensuremath{\theta}}
\newunicodechar{λ}{\ensuremath{\lambda}}
\newunicodechar{μ}{\ensuremath{\mu}}
\newunicodechar{σ}{\ensuremath{\sigma}}
\newunicodechar{τ}{\ensuremath{\tau}}
\newunicodechar{φ}{\ensuremath{\varphi}}
\newunicodechar{ψ}{\ensuremath{\psi}}
\newunicodechar{ω}{\ensuremath{\omega}}
\newunicodechar{Σ}{\ensuremath{\Sigma}}
% majuscules produites par \MakeUppercase dans les titres (α-aminés -> Α)
\newunicodechar{Α}{\ensuremath{\alpha}}
\newunicodechar{Β}{\ensuremath{\beta}}
\newunicodechar{Γ}{\ensuremath{\Gamma}}
\newunicodechar{Δ}{\ensuremath{\Delta}}
\newunicodechar{Ε}{\ensuremath{\varepsilon}}
\newunicodechar{Θ}{\ensuremath{\Theta}}
\newunicodechar{Λ}{\ensuremath{\Lambda}}
\newunicodechar{Μ}{\ensuremath{\mu}}
\newunicodechar{Τ}{\ensuremath{\tau}}
\newunicodechar{Φ}{\ensuremath{\Phi}}
\newunicodechar{Ψ}{\ensuremath{\Psi}}
\newunicodechar{Ω}{\ensuremath{\Omega}}
\newunicodechar{↦}{\ensuremath{\mapsto}}
\newunicodechar{∈}{\ensuremath{\in}}
\newunicodechar{∉}{\ensuremath{\notin}}
\newunicodechar{∧}{\ensuremath{\wedge}}
\newunicodechar{⇔}{\ensuremath{\Leftrightarrow}}
\newunicodechar{⟺}{\ensuremath{\Longleftrightarrow}}
\newunicodechar{⇒}{\ensuremath{\Rightarrow}}
\newunicodechar{⟹}{\ensuremath{\Longrightarrow}}
\newunicodechar{∘}{\ensuremath{\circ}}
\newunicodechar{∪}{\ensuremath{\cup}}
\newunicodechar{∩}{\ensuremath{\cap}}
\newunicodechar{≡}{\ensuremath{\equiv}}
\newunicodechar{⊂}{\ensuremath{\subset}}
\newunicodechar{⊥}{\ensuremath{\perp}}
\newunicodechar{∃}{\ensuremath{\exists}}
\newunicodechar{∀}{\ensuremath{\forall}}
\newunicodechar{∛}{\ensuremath{\sqrt[3]{\,}}}
\newunicodechar{∜}{\ensuremath{\sqrt[4]{\,}}}
\newunicodechar{⃗}{\ensuremath{{}^{\to}}}
\newunicodechar{ⁿ}{\ifmmode^{n}\else\textsuperscript{\ensuremath{n}}\fi}
\newunicodechar{ˣ}{\ifmmode^{x}\else\textsuperscript{\ensuremath{x}}\fi}
\newunicodechar{ᵇ}{\ifmmode^{b}\else\textsuperscript{\ensuremath{b}}\fi}
\newunicodechar{ʳ}{\ifmmode^{r}\else\textsuperscript{\ensuremath{r}}\fi}
\newunicodechar{ᵘ}{\ifmmode^{u}\else\textsuperscript{\ensuremath{u}}\fi}
\newunicodechar{ʸ}{\ifmmode^{y}\else\textsuperscript{\ensuremath{y}}\fi}
\newunicodechar{ᵃ}{\ifmmode^{a}\else\textsuperscript{\ensuremath{a}}\fi}
\newunicodechar{ᵉ}{\ifmmode^{e}\else\textsuperscript{\ensuremath{e}}\fi}
\newunicodechar{ᵏ}{\ifmmode^{k}\else\textsuperscript{\ensuremath{k}}\fi}
\newunicodechar{ᵖ}{\ifmmode^{p}\else\textsuperscript{\ensuremath{p}}\fi}
\newunicodechar{ᴺ}{\ifmmode^{N}\else\textsuperscript{\ensuremath{N}}\fi}
\newunicodechar{ᶜ}{\ifmmode^{c}\else\textsuperscript{\ensuremath{c}}\fi}
\newunicodechar{ʰ}{\ifmmode^{h}\else\textsuperscript{\ensuremath{h}}\fi}
\newunicodechar{ᵠ}{\ifmmode^{q}\else\textsuperscript{\ensuremath{q}}\fi}
\newunicodechar{ₙ}{\ifmmode_{n}\else\textsubscript{\ensuremath{n}}\fi}
\newunicodechar{ₐ}{\ifmmode_{a}\else\textsubscript{\ensuremath{a}}\fi}
\newunicodechar{ᵢ}{\ifmmode_{i}\else\textsubscript{\ensuremath{i}}\fi}
\newunicodechar{ᵣ}{\ifmmode_{r}\else\textsubscript{\ensuremath{r}}\fi}
\newunicodechar{ₖ}{\ifmmode_{k}\else\textsubscript{\ensuremath{k}}\fi}
\newunicodechar{ᵦ}{\ifmmode_{\beta}\else\textsubscript{\ensuremath{\beta}}\fi}
\newunicodechar{ₓ}{\ifmmode_{x}\else\textsubscript{\ensuremath{x}}\fi}
\newfontfamily\popdisplay{ArchivoBlack-Regular.ttf}[Path=../../../fonts/]
\newfontfamily\poplogo{SpaceGrotesk-Bold.ttf}[Path=../../../fonts/]
\newfontfamily\popsemi{PlusJakartaSans-SemiBold.ttf}[Path=../../../fonts/]
\newfontfamily\popxbold{PlusJakartaSans-ExtraBold.ttf}[Path=../../../fonts/]
\newfontfamily\popxboldls{PlusJakartaSans-ExtraBold.ttf}[Path=../../../fonts/,LetterSpace=3.0]

\setlist[enumerate]{leftmargin=1.7em,itemsep=5pt,topsep=4pt}
\setlist[itemize]{leftmargin=1.7em,itemsep=4pt,topsep=4pt,label={\color{poporange}\textbullet}}
\setlength{\parindent}{0pt}
\setlength{\parskip}{5pt}
\renewcommand{\arraystretch}{1.25}

% pgfplots : style Pop par défaut
\pgfplotsset{
  every axis/.append style={axis line style={popink,line width=1pt},tick style={popink},
    grid style={popline,line width=0.5pt},label style={font=\small},tick label style={font=\footnotesize}},
  popcurve/.style={poporange,line width=1.8pt,no markers,smooth},
}
% Même style disponible dans un \draw TikZ simple (hors pgfplots)
\tikzset{popcurve/.style={poporange,line width=1.8pt}}
\tikzset{popgrid/.style={popline,very thin}}
\colorlet{popcurve}{poporange} % le nom du style est parfois utilisé comme couleur

% ============================================================
% Pill / squiggle
% ============================================================
\newcommand{\poppill}[3]{%
  \tikz[baseline=-0.55ex]\node[draw=popink,line width=1.2pt,rounded corners=2.6pt,%
    fill=#2,inner xsep=6.5pt,inner ysep=3pt]{%
    \popxboldls\fontsize{7}{7}\selectfont\color{#3}\MakeUppercase{#1}};%
}
\newcommand{\popsquiggle}[1]{%
  \tikz[baseline=-0.4ex]{
    \draw[#1,line width=2.6pt,line cap=round]
      plot [smooth] coordinates {(0,0.09) (0.34,0.01) (0.68,0.21) (1.02,0.01) (1.36,0.10)};
  }%
}
% Mot-clé surligné (marqueur orange sous le texte)
\newcommand{\cle}[1]{{\popxbold\color{popdark}#1}}

% ============================================================
% En-tête / pied
% ============================================================
\pagestyle{fancy}
\fancyhf{}
\renewcommand{\headrulewidth}{0pt}
\renewcommand{\footrulewidth}{0pt}
\lhead{\raisebox{-0.15ex}{\poplogo\fontsize{13}{13}\selectfont\color{popink}OnBuch\color{poporange}+}}
\rhead{\poppill{\DOCMATIERE\ \textbullet\ \DOCNIVEAU\ \textbullet\ Leçon \DOCLECON}{white}{popink}}
\lfoot{\popxbold\fontsize{7.6}{7.6}\selectfont\color{popmuted}\MakeUppercase{Cours signé OnBuch+}\ \textbullet\ {\popsemi\DOCTITRE}}
\rfoot{\tikz[baseline=-0.6ex]\node[rounded corners=8.5pt,fill=popink,minimum height=17pt,minimum width=17pt,inner xsep=5pt,inner sep=0]{\popxbold\fontsize{8}{8}\selectfont\color{popcream}\thepage\,/\,\pageref*{LastPage}};}

% ============================================================
% Titres
% ============================================================
\newcommand{\chaptitre}[2]{%
  \begin{tcolorbox}[enhanced,colback=poporange,colframe=popink,boxrule=2.6pt,arc=11pt,
      drop shadow=popink,left=15pt,right=15pt,top=12pt,bottom=11pt,before skip=3pt,after skip=18pt]
    \poppill{#2}{popink}{popcream}\par\vspace{9pt}
    {\raggedright\hyphenpenalty=10000\exhyphenpenalty=10000\popdisplay\fontsize{22}{24}\selectfont\color{popcream}\MakeUppercase{#1}\par}
  \end{tcolorbox}
}
% Section numérotée : pastille orange avec le numéro + titre Archivo
\newcounter{popsec}
\newcommand{\coursec}[1]{%
  \stepcounter{popsec}\setcounter{popsub}{0}%
  \par\vspace{16pt}\noindent
  \begin{minipage}{\linewidth}
  \tikz[baseline=-0.7ex]\node[draw=popink,line width=1.6pt,fill=poporange,rounded corners=4pt,minimum size=22pt,inner sep=0]{\popdisplay\fontsize{12}{12}\selectfont\color{popcream}\thepopsec};%
  \hspace{8pt}{\popdisplay\fontsize{15}{17}\selectfont\color{popink}\MakeUppercase{#1}}\par
  \vspace{4pt}\noindent{\color{poporange}\rule{36pt}{3.6pt}}
  \end{minipage}\par\nopagebreak\vspace{8pt}
}
\newcounter{popsub}
\newcommand{\courssub}[1]{%
  \stepcounter{popsub}%
  \par\vspace{9pt}\noindent{\popxbold\fontsize{11.5}{13}\selectfont\color{popink}{\color{poporange}\thepopsec.\thepopsub}\hspace{6pt}#1}\par\nopagebreak\vspace{4pt}
}

% ============================================================
% Boîtes pédagogiques — même recette : fond blanc, bordure épaisse colorée,
% ombre dure encre, badge plein. Titre optionnel [#1].
% ============================================================
\tcbset{popbase/.style={breakable,enhanced,colback=white,boxrule=2.3pt,arc=9pt,
  shadow={3pt}{-3pt}{0pt}{popink},left=14pt,right=14pt,top=11pt,bottom=11pt,before skip=11pt,after skip=15pt,
  fontupper=\popsemi\fontsize{10.6}{14.5}\selectfont\color{popink},
  fontlower=\popsemi\fontsize{10.6}{14.5}\selectfont\color{popink}}}
% Coche dessinée (le glyphe ✓ n'existe pas dans les polices du document)
\newcommand{\popcheck}[1]{\tikz[baseline=-0.5ex]\draw[#1,line width=1.6pt,line cap=round,line join=round](-0.13,0.0)--(-0.04,-0.09)--(0.13,0.1);}
\providecommand{\coche}{\popcheck{popgreen}}
\providecommand{\resultat}[1]{\par\smallskip\noindent\fcolorbox{popgreen}{popgreenL}{\parbox{\dimexpr\linewidth-2\fboxsep-2\fboxrule\relax}{\textbf{Résultat.} #1}}\par\smallskip}
\newcommand{\popboxhead}[3]{\poppill{#1}{#2}{white}\ifx\relax#3\relax\else\hspace{6pt}{\popxbold\fontsize{10.6}{12}\selectfont\color{popink}#3}\fi\par\vspace{7pt}}

\newtcolorbox{aretenir}[1][]{popbase,colframe=popgreen,before upper={\popboxhead{À retenir}{popgreen}{#1}}}
% Texte en anglais (document de lecture, dialogue, extrait) : cadre
% d'étude. Un titre optionnel (source, auteur) s'affiche dans l'en-tête.
\newtcolorbox{textebox}[1][]{popbase,colframe=popblue,before upper={\popboxhead{Text}{popblue}{#1}\begin{otherlanguage}{english}},after upper={\end{otherlanguage}}}
% Marques ✓ / ✗ (forme correcte / fautive) souvent utilisées par les modèles
\providecommand{\cross}{\textcolor{poppink}{\ensuremath{\times}}}
\providecommand{\xmark}{\cross}\providecommand{\crossmark}{\cross}\providecommand{\cmark}{\ensuremath{\checkmark}}
% Ligne de réponse / justification à compléter (inventée par les modèles)
\providecommand{\justification}[1]{\par\noindent\textit{Justification~:} \dotfill\par}
% \textipa (package tipa non chargé) : la police ne couvre pas l'API -> simple passage
\providecommand{\textipa}[1]{#1}
% Soulignement (ulem non chargé) : \uline, \ul
\providecommand{\uline}[1]{\underline{#1}}\providecommand{\ul}[1]{\underline{#1}}
% \glossaire{\item ...} inventée par les modèles : liste de glossaire
\providecommand{\glossaire}[1]{\begin{itemize}#1\end{itemize}}
% \alert (beamer) inventée par les modèles : texte mis en évidence
\providecommand{\alert}[1]{\textbf{\textcolor{poppink}{#1}}}
% variantes ASCII d'ordinaux écrites en macro
\providecommand{\iere}{\textsuperscript{re}}\providecommand{\ieme}{\textsuperscript{e}}\providecommand{\ere}{\textsuperscript{re}}\providecommand{\eme}{\textsuperscript{e}}\providecommand{\er}{\textsuperscript{er}}\providecommand{\ie}{\textsuperscript{e}}
% Ordinaux français écrits en macro par les modèles : 1\ière, 3\ième
\newcommand{\ière}{\textsuperscript{re}}\newcommand{\ième}{\textsuperscript{e}}
% \exemple (inventée par les modèles) : étiquette « Exemple : »
\providecommand{\exemple}{\textbf{Exemple~:}\ }
% \square utilisable aussi hors mode math (cases à cocher)
\AtBeginDocument{\let\squareMath\square\renewcommand{\square}{\ensuremath{\squareMath}}}
% Mot ou phrase en anglais à l'intérieur d'un paragraphe français
\newcommand{\eng}[1]{\ifmmode\text{\textit{#1}}\else\begingroup\selectlanguage{english}\itshape #1\endgroup\fi}
\let\esp\eng % alias toléré
% flèches d'intonation inventées par les modèles
\providecommand{\textsearrow}{}\renewcommand{\textsearrow}{\ensuremath{\searrow}}\providecommand{\textnearrow}{}\renewcommand{\textnearrow}{\ensuremath{\nearrow}}
% \fr{..} : mot français (inventé par les modèles)
\providecommand{\fr}[1]{\textit{#1}}
\newtcolorbox{exemplebox}[1][]{popbase,colframe=poporange,before upper={\popboxhead{Exemple}{poporange}{#1}}}
\newtcolorbox{definition}[1][]{popbase,colframe=popblue,before upper={\popboxhead{Définition}{popblue}{#1}}}
\newtcolorbox{propriete}[1][]{popbase,colframe=poppurple,before upper={\popboxhead{Propriété}{poppurple}{#1}}}
\newtcolorbox{methode}[1][]{popbase,colframe=popgold,before upper={\popboxhead{Méthode}{popgold}{#1}}}
\newtcolorbox{attention}[1][]{popbase,colframe=poppink,before upper={\popboxhead{Attention}{poppink}{#1}}}
\newtcolorbox{experience}[1][]{popbase,colframe=popblue,colback=popblueL,before upper={\popboxhead{Expérience}{popblue}{#1}}}
% « Le savais-tu ? » : carte violette pleine (contexte, culture, Cameroun)
\newtcolorbox{savaistu}[1][]{popbase,colback=poppurple,colframe=popink,
  fontupper=\popsemi\fontsize{10.4}{14.2}\selectfont\color{white},
  before upper={\poppill{Le savais-tu\,?}{white}{poppurple}\ifx\relax#1\relax\else\hspace{6pt}{\popxbold\color{white}#1}\fi\par\vspace{7pt}}}
% Exercice résolu : énoncé en haut, \tcblower puis la solution sur fond crème
\newcounter{popexo}\newcounter{popexr}
\newtcolorbox{exoresolu}[1][]{popbase,colframe=poporange,colbacklower=popcream,
  segmentation style={popink,line width=1.2pt,dashed},
  before upper={\stepcounter{popexr}\popboxhead{Exercice résolu \thepopexr}{poporange}{#1}},
  before lower={\poppill{Solution}{popink}{popcream}\par\vspace{6pt}}}
% Exercice (non résolu en place) — la correction est dans la partie Corrigés
\newtcolorbox{exercice}[1][]{popbase,colframe=popink,
  before upper={\stepcounter{popexo}\popboxhead{Exercice \thepopexo}{popink}{#1}}}
% Corrigé
\newtcolorbox{corrige}[1][]{popbase,colframe=popgreen,colback=popgreenL,
  before upper={\popboxhead{Corrigé}{popgreen}{#1}}}
% Objectifs / prérequis / bilan
\newtcolorbox{objectifs}{popbase,colframe=popink,colback=poporangeL,
  before upper={\popboxhead{Objectifs de la leçon}{popink}{}}}
\newtcolorbox{prerequis}{popbase,colframe=popink,colback=popblueL,
  before upper={\popboxhead{Prérequis}{popblue}{}}}
\newtcolorbox{bilan}{popbase,colframe=popink,colback=white,boxrule=2.8pt,
  before upper={\popboxhead{Fiche bilan}{poporange}{L'essentiel à connaître pour l'examen}}}
% Figure encadrée avec légende
\newtcolorbox{popfigure}[1][]{enhanced,breakable=false,colback=white,colframe=popline,boxrule=1.4pt,arc=8pt,
  left=10pt,right=10pt,top=10pt,bottom=8pt,before skip=10pt,after skip=12pt,halign=center,
  lower separated=false,halign lower=center,
  fontlower=\popsemi\fontsize{9}{11.5}\selectfont\color{popmuted},
  #1}
\newcommand{\legende}[1]{\par\vspace{6pt}{\popsemi\fontsize{9}{11.5}\selectfont\color{popmuted}#1}}

% ============================================================
% Signature OnBuch+
% ============================================================
\newcommand{\poplogoinline}[1]{{\poplogo\fontsize{#1}{#1}\selectfont\color{popink}OnBuch\color{poporange}+}}
\newcommand{\signatureonbuch}{%
  \begin{tcolorbox}[enhanced,colback=popink,colframe=popink,boxrule=2.2pt,arc=10pt,
      drop shadow=poporange,left=14pt,right=14pt,top=10pt,bottom=10pt,before skip=0pt,after skip=0pt]
    \tikz[baseline=-0.7ex]\node[circle,fill=popgreen,draw=popcream,line width=1.3pt,minimum size=22pt,inner sep=0]{\popcheck{white}};%
    \hspace{9pt}\begin{minipage}[c]{0.62\linewidth}
      {\popxbold\fontsize{12}{13}\selectfont\color{popcream}Signé\ \textbullet\ {\poplogo OnBuch\color{poporange}+}}\par\vspace{2pt}
      {\raggedright\popsemi\fontsize{8}{10.5}\selectfont\color{popline}\MakeUppercase{Équipe pédagogique OnBuch+}\\\MakeUppercase{Conforme au programme officiel MINESEC}\par}
    \end{minipage}\hfill
    {\poplogo\fontsize{22}{22}\selectfont\color{popcream}OnBuch\color{poporange}+}
  \end{tcolorbox}%
}

% ============================================================
% Page de garde
% #1 titre · #2 matière · #3 niveau · #4 module · #5 n° leçon · #6 durée ·
% #7 description / objectifs (texte ou itemize)
% ============================================================
\newcommand{\pagegarde}[7]{%
  \thispagestyle{empty}
  \newgeometry{a4paper,margin=1.7cm,top=1.6cm,bottom=1.4cm}%
  \begin{tikzpicture}[remember picture,overlay]
    \fill[poporange!18] ([xshift=-3.2cm,yshift=-3.6cm]current page.north east) circle (4.6cm);
    \fill[poppurple!14] ([xshift=2.4cm,yshift=7.6cm]current page.south west) circle (3.8cm);
  \end{tikzpicture}%
  {\poplogo\fontsize{30}{30}\selectfont\color{popink}OnBuch\color{poporange}+}\hfill
  \raisebox{0.6ex}{\poppill{Cours signé \textbullet\ Premium}{popink}{popcream}}\par
  \vspace{1pt}\popsquiggle{poppurple}\par
  \vspace{8pt}
  \begin{tcolorbox}[enhanced,colback=poporange,colframe=popink,boxrule=2.8pt,arc=13pt,
      drop shadow=popink,left=18pt,right=18pt,top=12pt,bottom=13pt,after skip=10pt]
    \poppill{#2 \textbullet\ #3}{popink}{popcream}\hspace{5pt}\poppill{#4}{white}{popink}\par\vspace{8pt}
    {\popdisplay\fontsize{14}{14}\selectfont\color{popink}LEÇON #5}\par\vspace{4pt}
    {\raggedright\hyphenpenalty=10000\exhyphenpenalty=10000\popdisplay\fontsize{25}{27}\selectfont\color{popcream}\MakeUppercase{#1}\par}
    \vspace{8pt}
    {\popxbold\fontsize{9.5}{9.5}\selectfont\color{popink}\MakeUppercase{Durée conseillée : #6}}
  \end{tcolorbox}
  \begin{tcolorbox}[enhanced,colback=white,colframe=popink,boxrule=2.2pt,arc=10pt,
      drop shadow=popink,left=16pt,right=16pt,top=10pt,bottom=10pt,after skip=8pt]
    \poppill{Dans cette leçon}{poppurple}{white}\par\vspace{5pt}
    \popsemi\fontsize{9.3}{12.6}\selectfont\color{popink}#7
  \end{tcolorbox}
  \noindent\begin{minipage}[t]{0.487\linewidth}
    \begin{tcolorbox}[enhanced,colback=popgreen,colframe=popink,boxrule=2.2pt,arc=10pt,
        drop shadow=popink,left=11pt,right=11pt,top=10pt,bottom=10pt,height=3.1cm,valign=center]
      \tcbox[colback=white,colframe=popink,boxrule=1.3pt,arc=5pt,boxsep=0pt,nobeforeafter,
        left=3pt,right=3pt,top=3pt,bottom=3pt,valign=center]{\qrcode[height=1.9cm]{https://play.google.com/store/apps/details?id=com.luvvix.onbuch&hl=fr}}%
      \hspace{8pt}\begin{minipage}[c]{0.5\linewidth}
        {\popdisplay\fontsize{11}{12.5}\selectfont\color{white}\MakeUppercase{Télécharge OnBuch}}\par\vspace{4pt}
        {\raggedright\popsemi\fontsize{8.4}{11}\selectfont\color{white}Gratuit \textbullet\ Google Play\\Tuteur IA, annales, concours, résultats\par}
      \end{minipage}
    \end{tcolorbox}
  \end{minipage}\hfill
  \begin{minipage}[t]{0.487\linewidth}
    \begin{tcolorbox}[enhanced,colback=poppurple,colframe=popink,boxrule=2.2pt,arc=10pt,
        drop shadow=popink,left=11pt,right=11pt,top=10pt,bottom=10pt,height=3.1cm,valign=center]
      \tikz[baseline=-0.7ex]\node[circle,fill=white,draw=popink,line width=1.3pt,minimum size=1.6cm,inner sep=0]{\popdisplay\fontsize{24}{24}\selectfont\color{poppurple}+};%
      \hspace{8pt}\begin{minipage}[c]{0.6\linewidth}
        {\popdisplay\fontsize{11}{12.5}\selectfont\color{white}\MakeUppercase{Passe à OnBuch+}}\par\vspace{4pt}
        {\raggedright\popsemi\fontsize{8.4}{11}\selectfont\color{white}Tous les cours signés, Tuteur IA illimité, fiches PDF — onglet OnBuch+ de l'app\par}
      \end{minipage}
    \end{tcolorbox}
  \end{minipage}
  \vspace{8pt}
  \popcontenu
  \vfill
  \signatureonbuch
  \restoregeometry
  \newpage
}

% Rangée « Ce cours contient » (page de garde)
\newcommand{\poptile}[3]{% #1 couleur #2 gros texte #3 légende
  \begin{tcolorbox}[enhanced,colback=#1,colframe=popink,boxrule=2pt,arc=9pt,shadow={2.5pt}{-2.5pt}{0pt}{popink},
      left=6pt,right=6pt,top=9pt,bottom=9pt,halign=center,valign=center,height=1.75cm,width=\linewidth,nobeforeafter]
    {\popdisplay\fontsize{12.5}{13}\selectfont\color{white}\MakeUppercase{#2}}\par\vspace{3pt}
    {\centering\popsemi\fontsize{7.8}{9.5}\selectfont\color{white}#3\par}
  \end{tcolorbox}}
\newcommand{\popcontenu}{%
  \par\noindent{\popxboldls\fontsize{7.5}{7.5}\selectfont\color{popmuted}CE COURS SIGNÉ CONTIENT}\par\vspace{7pt}
  \noindent\begin{minipage}[t]{0.235\linewidth}\poptile{popblue}{Cours}{complet, illustré, pas à pas}\end{minipage}\hfill
  \begin{minipage}[t]{0.235\linewidth}\poptile{poporange}{Méthodes}{et exercices résolus}\end{minipage}\hfill
  \begin{minipage}[t]{0.235\linewidth}\poptile{poppink}{Exercices}{type examen + corrigés}\end{minipage}\hfill
  \begin{minipage}[t]{0.235\linewidth}\poptile{popgreen}{Fiche bilan}{l'essentiel à retenir}\end{minipage}\par}

% Sommaire
\newtcolorbox{sommaire}{enhanced,colback=white,colframe=popink,boxrule=2.2pt,arc=10pt,
  drop shadow=popink,left=18pt,right=18pt,top=13pt,bottom=13pt,before skip=6pt,after skip=20pt,
  before upper={\poppill{Sommaire}{poppurple}{white}\par\vspace{9pt}\popsemi\fontsize{10.6}{14.5}\selectfont\color{popink}}}

% Signature de fin de cours — poussée tout en bas de la dernière page
\newcommand{\findecours}{%
  \par\vfill\nopagebreak
  \begin{center}\popsquiggle{poporange}\end{center}\vspace{4pt}
  \signatureonbuch
}
\AtBeginDocument{\def\llbracket{\mathopen{[\mkern-3.5mu[}}\def\rrbracket{\mathclose{]\mkern-3.5mu]}}} % intervalles d'entiers (glyphe ⟦ absent de la police)
% Glyphes absents de Fira Math : redéfinis à partir de glyphes disponibles
\AtBeginDocument{%
  \def\setminus{\mathbin{\backslash}}\let\smallsetminus\setminus
  \def\vdots{\mathord{\vbox{\baselineskip3.2pt\lineskiplimit0pt\kern4pt\hbox{.}\hbox{.}\hbox{.}}}}%
  \def\ddots{\mathinner{\mkern1mu\raise6.5pt\hbox{.}\mkern2mu\raise3.6pt\hbox{.}\mkern2mu\raise0.7pt\hbox{.}\mkern1mu}}%
  \def\triangle{\mathord{\scalebox{0.9}{$\symup{\Delta}$}}}%
  \def\nabla{\mathord{\raisebox{\depth}{\scalebox{1}[-1]{$\symup{\Delta}$}}}}%
  \def\checkmark{\popcheck{popdark}}%
  \def\star{\mathbin{\ast}}\def\bigstar{\mathord{\ast}}%
  \def\diamond{\mathbin{\scalebox{0.6}{$\Diamond$}}}%
  \def\bigcup{\mathop{\mathchoice{\raisebox{-0.3em}{\scalebox{1.7}{$\cup$}}}{\scalebox{1.25}{$\cup$}}{\cup}{\cup}}\displaylimits}%
  \def\bigcap{\mathop{\mathchoice{\raisebox{-0.3em}{\scalebox{1.7}{$\cap$}}}{\scalebox{1.25}{$\cap$}}{\cap}{\cap}}\displaylimits}%
  \def\lesseqqgtr{\mathrel{\lessgtr}}\def\bigoplus{\mathop{\oplus}\displaylimits}%
}
\providecommand{\ssi}{si et seulement si} % abréviation fréquente
\AtBeginDocument{\providecommand{\wideparen}{\overparen}} % arc de cercle

%% FIGFIX-BEGIN v3 — correctifs globaux des figures (docs/onbuch-generation/tools/figfix)
\usepackage{adjustbox}\usepackage{etoolbox}
% toute figure TikZ/pgfplots plus large que la ligne est réduite à la largeur disponible
\BeforeBeginEnvironment{tikzpicture}{\begin{adjustbox}{max width=\linewidth}}
\AfterEndEnvironment{tikzpicture}{\end{adjustbox}}
% légendes pgfplots lisibles par-dessus les courbes
\pgfplotsset{every axis/.append style={legend style={fill=white,fill opacity=0.92,text opacity=1,draw=black!15,inner sep=3pt}}}
% étiquettes d'arêtes (syntaxe edge["..."]) avec fond blanc
\tikzset{every edge quotes/.append style={fill=white,inner sep=1.2pt}}
% bulles de cartes mentales : texte à la ligne dans la bulle, bulle un peu plus grande
\tikzset{every concept/.append style={text width=2.0cm,align=center,minimum size=2.3cm,inner sep=2pt}}
%% FIGFIX-END
\begin{document}

\pagegarde{Lesson 13 — Speaking: Exchanges information about buying and selling articles/foodstuff at a nearby store/market}{Anglais}{1ère A}{Chapter 2 : Economic Life and Occupations}{ENA14}{2 h}{Cette leçon d'expression orale apprend aux élèves de 1ère A à mener un dialogue complet d'achat et de vente dans un magasin ou au marché : saluer, demander des informations sur les prix et les quantités, négocier, conclure la transaction. À travers des dialogues modèles, des formules fonctionnelles et des jeux de rôle guidés, les élèves développent leur aisance, leur prononciation et leur intonation en anglais.
\vspace{4pt}
\begin{multicols}{2}\small
\begin{itemize}
\item À la fin de cette leçon, je sais saluer et prendre contact poliment avec un vendeur ou un client en anglais.
\item À la fin de cette leçon, je sais demander le prix, la quantité et la disponibilité d'un article ou d'un produit alimentaire.
\item À la fin de cette leçon, je sais donner une information précise (prix, poids, mesure, qualité) à un client.
\item À la fin de cette leçon, je sais relancer la conversation et négocier un prix de façon polie.
\item À la fin de cette leçon, je sais conclure une transaction : accepter, refuser, payer et remercier.
\item À la fin de cette leçon, je sais utiliser le vocabulaire des articles et des denrées alimentaires du marché camerounais.
\end{itemize}\end{multicols}}

\begin{sommaire}
\begin{multicols}{2}
\begin{enumerate}
\item Setting the scene: at the store and at the market
\item Functional language: asking for and giving information
\item Bargaining, relaunching and concluding the exchange
\item Model dialogues
\item Pronunciation, stress and intonation
\item Guided role-plays: let's speak!
\item Activité d'intégration
\item Méthodes et exercices résolus
\item Exercices
\item Corrigés des exercices
\item Fiche bilan
\end{enumerate}
\end{multicols}
\end{sommaire}

\begin{objectifs}
\begin{itemize}
\item À la fin de cette leçon, je sais saluer et prendre contact poliment avec un vendeur ou un client en anglais.
\item À la fin de cette leçon, je sais demander le prix, la quantité et la disponibilité d'un article ou d'un produit alimentaire.
\item À la fin de cette leçon, je sais donner une information précise (prix, poids, mesure, qualité) à un client.
\item À la fin de cette leçon, je sais relancer la conversation et négocier un prix de façon polie.
\item À la fin de cette leçon, je sais conclure une transaction : accepter, refuser, payer et remercier.
\item À la fin de cette leçon, je sais utiliser le vocabulaire des articles et des denrées alimentaires du marché camerounais.
\item À la fin de cette leçon, je sais prononcer correctement les nombres, les prix et les formules de politesse avec la bonne intonation.
\item À la fin de cette leçon, je sais jouer un court dialogue vendeur/client devant la classe avec un partenaire.
\end{itemize}
\end{objectifs}

\begin{prerequis}
\begin{itemize}
\item Les nombres cardinaux en anglais (Seconde) : one, twenty, one hundred, one thousand...
\item Les salutations et formules de politesse de base : Good morning, please, thank you, excuse me.
\item Le vocabulaire élémentaire de la nourriture (Seconde) : rice, bread, fish, oil, tomatoes...
\item Les questions en Wh- et en Yes/No (present simple) : How much...? Do you have...?
\item Les quantifieurs de base : some, any, much, many, a kilo of, a bottle of.
\end{itemize}
\end{prerequis}

\newpage

\coursec{Setting the scene: at the store and at the market}

It is Saturday morning in Mokolo Market, Yaoundé. Amina wants to buy two kilos of tomatoes, a bottle of groundnut oil and some bread for the weekend. The trader greets her in English: \eng{“Good morning, Madam! What can I do for you today?”} Amina hesitates… How should she ask for prices, bargain politely and pay? In the North West and South West regions of Cameroon, as in Nigeria or Ghana, these exchanges happen in English every day. This lesson will give you the exact words and confidence to buy and sell like a professional!

\courssub{Places where we buy and sell}

Before we speak, we must know \cle{where} the exchange takes place. In Cameroon, you will hear these common places:

\begin{textebox}[Places of purchase]
A \textbf{nearby store} is a small shop in your neighbourhood. A \textbf{provision store} sells mainly foodstuff and household items. A \textbf{market stall} is a table or covered space inside a big market like Mokolo or Mfoundi. A \textbf{supermarket} is a large self‑service shop (e.g. Mahima, Casino). A \textbf{roadside kiosk} is a tiny wooden cabin along the road, often selling phone credit, sweets and cold drinks.
\end{textebox}

\begin{center}
\begin{tabularx}{\linewidth}{|X|X|X|}
\hline
\rowcolor{popblueL} \textbf{English} & \textbf{Français} & \textbf{Remarque} \\
\hline
a nearby store & un petit magasin de quartier & \emph{store} (US) = \emph{shop} (UK) \\
a provision store & un magasin de provisions & \emph{provisions} = vivres, denrées \\
a market stall & un étal de marché & \emph{stall} = emplacement loué \\
a supermarket & un supermarché & libre‑service, grands rayons \\
a roadside kiosk & un kiosque au bord de la route & souvent en bois, toit en tôle \\
\hline
\end{tabularx}
\end{center}

\begin{popfigure}
\begin{tikzpicture}[
  mindmap,
  concept color=popblue!60,
  level 1/.style={level distance=3.2cm,sibling angle=72},
  level 2/.style={level distance=2.2cm,sibling angle=36},
  every node/.style={font=\small, text=popdark, align=center}
]
\node[concept] {BUYING AND SELLING}
  child[concept color=poporange!70] { node[concept] {Places}
    child { node[concept] {nearby store} }
    child { node[concept] {provision store} }
    child { node[concept] {market stall} }
    child { node[concept] {supermarket} }
    child { node[concept] {roadside kiosk} }
  }
  child[concept color=popgreen!70] { node[concept] {People}
    child { node[concept] {customer / buyer} }
    child { node[concept] {seller / trader} }
    child { node[concept] {shopkeeper} }
    child { node[concept] {stallholder} }
  }
  child[concept color=poppurple!70] { node[concept] {Articles \& Foodstuff}
    child { node[concept] {soap, matches, candles} }
    child { node[concept] {rice, beans, gari} }
    child { node[concept] {plantains, tomatoes} }
    child { node[concept] {oil, dried fish} }
  }
  child[concept color=poppink!70] { node[concept] {Money \& Measures}
    child { node[concept] {price, cost, discount} }
    child { node[concept] {change, FCFA} }
    child { node[concept] {bag, kilo, litre} }
    child { node[concept] {bottle, tin, packet} }
    child { node[concept] {bunch, cup} }
  };
\end{tikzpicture}
\legende{Carte mentale lexicale : les quatre piliers de l'acte d'achat-vente.}
\end{popfigure}

\courssub{People in the transaction}

Two main roles exist. The person who buys is the \cle{customer} or the \cle{buyer}. The person who sells can be called \cle{seller}, \cle{trader}, \cle{shopkeeper} (in a shop) or \cle{stallholder} (at a market stall).

\begin{definition}[customer / buyer]
\eng{customer} (countable) = le client, celui qui achète. \eng{buyer} = l'acheteur (terme plus formel, souvent utilisé dans les contrats).
\end{definition}

\begin{definition}[seller / trader / shopkeeper / stallholder]
\eng{seller} = le vendeur (terme général). \eng{trader} = le commerçant (qui achète et revend). \eng{shopkeeper} = le propriétaire ou gérant d'une boutique. \eng{stallholder} = le titulaire d'un étal sur un marché.
\end{definition}

\begin{center}
\begin{tabularx}{\linewidth}{|X|X|X|}
\hline
\rowcolor{popblueL} \textbf{English} & \textbf{Français} & \textbf{Contexte typique} \\
\hline
the customer & le client & supermarché, boutique \\
the buyer & l'acheteur & transaction formelle, marché \\
the seller & le vendeur & terme neutre, universel \\
the trader & le commerçant & marché, gros volume \\
the shopkeeper & le boutiquier & provision store, nearby store \\
the stallholder & le titulaire d'étal & market stall (Mokolo, Mfoundi) \\
\hline
\end{tabularx}
\end{center}

\courssub{Everyday articles and foodstuff}

We distinguish \cle{articles} (non‑food items) from \cle{foodstuff} (denrées alimentaires). Both are essential in a Cameroonian shopping basket.

\begin{definition}[foodstuff]
\eng{foodstuff} (uncountable) = les denrées alimentaires, c'est‑à‑dire tout ce qu'on achète pour cuisiner et manger (rice, oil, fish, etc.). On ne dit pas “a foodstuff” mais “an item of foodstuff” ou “a food item”.
\end{definition}

\begin{center}
\begin{tabularx}{\linewidth}{|X|X|}
\hline
\rowcolor{popblueL} \textbf{Articles (non‑food)} & \textbf{Foodstuff (denrées)} \\
\hline
soap (savon) & rice (riz) \\
matches (allumettes) & beans (haricots) \\
candles (bougies) & gari (gari / farine de manioc) \\
batteries (piles) & plantains (plantains) \\
notebooks (cahiers) & tomatoes (tomates) \\
pens (stylos) & onions (oignons) \\
sugar (sucre) & groundnut oil (huile d'arachide) \\
salt (sel) & palm oil (huile de palme) \\
\hline
\end{tabularx}
\end{center}

\begin{exemplebox}[Vocabulaire courant — traduction]
\eng{soap} = savon ; \eng{matches} = allumettes ; \eng{candles} = bougies ; \eng{batteries} = piles ; \eng{notebooks} = cahiers ; \eng{pens} = stylos ; \eng{sugar} = sucre ; \eng{salt} = sel. \\
\eng{rice} = riz ; \eng{beans} = haricots ; \eng{gari} = gari ; \eng{plantains} = plantains ; \eng{tomatoes} = tomates ; \eng{onions} = oignons ; \eng{groundnut oil} = huile d'arachide ; \eng{palm oil} = huile de palme ; \eng{dried fish} = poisson séché ; \eng{meat} = viande ; \eng{eggs} = œufs ; \eng{bread} = pain ; \eng{flour} = farine.
\end{exemplebox}

\courssub{Containers and measures}

In a market, goods are rarely sold “by the piece” only. We use specific containers and measures. Mastering them avoids confusion at the stall.

\begin{exemplebox}[Containers and measures — exemples traduits]
\eng{a bag of rice} = un sac de riz ; \eng{a kilo of meat} = un kilo de viande ; \eng{a litre of oil} = un litre d'huile ; \eng{a bottle of palm wine} = une bouteille de vin de palme ; \eng{a tin of tomatoes} = une boîte de tomates ; \eng{a packet of sugar} = un paquet de sucre ; \eng{a bunch of plantains} = un régime de plantains ; \eng{a cup of gari} = une tasse de gari.
\end{exemplebox}

\begin{center}
\begin{tabularx}{\linewidth}{|X|X|X|}
\hline
\rowcolor{popblueL} \textbf{English measure} & \textbf{Français} & \textbf{Produit typique} \\
\hline
a bag of … & un sac de … & rice, beans, gari, flour \\
a kilo of … & un kilo de … & meat, fish, tomatoes, onions \\
a litre of … & un litre de … & groundnut oil, palm oil, petrol \\
a bottle of … & une bouteille de … & palm wine, soda, water \\
a tin of … & une boîte (métallique) de … & tomatoes, sardines, milk \\
a packet of … & un paquet / sachet de … & sugar, salt, biscuits, spaghetti \\
a bunch of … & un régime / une botte de … & plantains, bananas, vegetables \\
a cup of … & une tasse / un bol de … & gari, rice (mesure locale), beans \\
\hline
\end{tabularx}
\end{center}

\begin{popfigure}
\begin{tikzpicture}[
  every node/.style={font=\small, text=popdark},
  label/.style={rectangle, draw=popblue, fill=popblueL, rounded corners=2pt, inner sep=2pt},
  stall/.style={rectangle, draw=poporange, thick, minimum width=8cm, minimum height=4cm}
]
\node[stall] (stall) {};
% Labels inside stall
\node[label, anchor=north west] at ([xshift=0.5cm,yshift=-0.5cm]stall.north west) {tomatoes};
\node[label, anchor=north east] at ([xshift=-0.5cm,yshift=-0.5cm]stall.north east) {onions};
\node[label, anchor=south west] at ([xshift=0.5cm,yshift=0.5cm]stall.south west) {groundnut oil};
\node[label, anchor=south east] at ([xshift=-0.5cm,yshift=0.5cm]stall.south east) {rice (bags)};
% Seller and customer
\node[draw=popgreen, fill=popgreenL, circle, inner sep=2pt] (seller) at ([yshift=-2.5cm]stall.south) {};
\node[below=2mm of seller] {seller};
\node[draw=poppink, fill=poppinkL, circle, inner sep=2pt] (customer) at ([yshift=-4cm]stall.south) {};
\node[below=2mm of customer] {customer};
% Arrow
\draw[->, thick, poporange] (seller) -- (customer) node[midway, right] {exchange};
\end{tikzpicture}
\legende{Schéma simplifié d'un étal de marché (stall) avec quatre produits clés.}
\end{popfigure}

\courssub{Money and prices}

Talking about money requires precise vocabulary. In Cameroon, prices are in \cle{FCFA} (Franc CFA). The key words are \cle{price}, \cle{cost}, \cle{discount}, \cle{change}.

\begin{center}
\begin{tabularx}{\linewidth}{|X|X|X|}
\hline
\rowcolor{popblueL} \textbf{English} & \textbf{Français} & \textbf{Exemple} \\
\hline
the price & le prix & \eng{The price is 2,000 FCFA.} \\
the cost & le coût & \eng{The cost of a bag of rice is 25,000 FCFA.} \\
a discount & une remise / une réduction & \eng{Can you give me a discount?} \\
change & la monnaie (rendue) & \eng{Here is your change: 500 FCFA.} \\
500 FCFA & cinq cents francs CFA & \eng{It costs 500 FCFA.} \\
2,000 FCFA & deux mille francs CFA & \eng{That makes 2,000 FCFA altogether.} \\
\hline
\end{tabularx}
\end{center}

\begin{attention}[How much or How many? — Piège classique]
Ne dites jamais \eng{“How many is it?”} pour demander le prix. \cle{Price} est indénombrable (uncountable), on utilise \cle{much} :
\eng{“How much is it?”} = Combien ça coûte ? \\
\eng{“How much does it cost?”} = Combien ça coûte ? \\
\eng{“How much for the tomatoes?”} = Combien pour les tomates ? \\
\eng{How many} ne s'utilise qu'avec des noms dénombrables au pluriel : \eng{“How many kilos do you want?”} (Combien de kilos voulez‑vous ?).
\end{attention}

\begin{savaistu}[“How much last?” — Cameroun / Nigeria]
Au Cameroun anglophone et au Nigeria, on entend souvent \eng{“How much last?”} pour demander le dernier prix (prix final après négociation). En anglais standard, on dit : \eng{“What is your last price?”} ou \eng{“Is that your final price?”}. Apprenez la forme standard pour les examens et les contextes formels.
\end{savaistu}

\courssub{Authentic situation: At Mokolo Market}

The following dialogue puts everything together. Read it aloud, notice the containers, the polite forms and the money words.

\begin{textebox}[At Mokolo Market]
Customer: Good morning, Madam!
Seller: Good morning, my dear. What can I do for you today?
Customer: I want to buy some foodstuff for Sunday lunch. How much is a bag of rice?
Seller: A bag of rice is 25,000 FCFA. But I can give you a discount: 24,000 FCFA.
Customer: That's a good price. I'll take one bag. And how much for a litre of groundnut oil?
Seller: 1,500 FCFA per litre.
Customer: Give me two litres, please. Also, I need a bunch of plantains and a tin of tomatoes.
Seller: The plantains are 1,000 FCFA the bunch. The tin of tomatoes is 600 FCFA.
Customer: Okay. Here is 30,000 FCFA.
Seller: Thank you. Let me count… Here is your change: 2,300 FCFA.
Customer: Thank you, Madam. Have a nice day!
Seller: You too, my dear. Come again!
\end{textebox}

\begin{center}
\begin{tabularx}{\linewidth}{|X|X|}
\hline
\rowcolor{popblueL} \textbf{Mot difficile} & \textbf{Traduction / Explication} \\
\hline
foodstuff & denrées alimentaires (indénombrable) \\
discount & remise, réduction sur le prix \\
per litre & par litre (prix unitaire) \\
change & monnaie rendue au client \\
count & compter (l'argent) \\
Come again! & Revenez ! (formule de politesse commerciale) \\
\hline
\end{tabularx}
\end{center}

\courssub{Consolidation: asking for the price}

\begin{exoresolu}[Asking for the price correctly]
\textbf{Énoncé :} Complétez les questions par \eng{How much} ou \eng{How many} selon le nom qui suit. Traduisez chaque phrase.
\begin{enumerate}
\item \_\_\_\_\_\_ is a kilo of meat?
\item \_\_\_\_\_\_ kilos of rice do you need?
\item \_\_\_\_\_\_ for a bottle of palm wine?
\item \_\_\_\_\_\_ tins of tomatoes are on the shelf?
\item \_\_\_\_\_\_ does a bunch of plantains cost?
\end{enumerate}

\textbf{Solution détaillée :}
\begin{enumerate}
\item \eng{How much is a kilo of meat?} — \eng{kilo} (singulier) + \eng{meat} (indénombrable) $\rightarrow$ \cle{How much}. Traduction : « Combien coûte un kilo de viande ? »
\item \eng{How many kilos of rice do you need?} — \eng{kilos} (pluriel dénombrable) $\rightarrow$ \cle{How many}. Traduction : « Combien de kilos de riz te faut‑il ? »
\item \eng{How much for a bottle of palm wine?} — \eng{bottle} (singulier) + \eng{palm wine} (indénombrable) $\rightarrow$ \cle{How much}. Traduction : « Combien pour une bouteille de vin de palme ? »
\item \eng{How many tins of tomatoes are on the shelf?} — \eng{tins} (pluriel dénombrable) $\rightarrow$ \cle{How many}. Traduction : « Combien de boîtes de tomates sont sur l'étagère ? »
\item \eng{How much does a bunch of plantains cost?} — \eng{bunch} (singulier) + \eng{plantains} (pluriel mais \eng{a bunch of} = un régime, quantité unique) $\rightarrow$ \cle{How much}. Traduction : « Combien coûte un régime de plantains ? »
\end{enumerate}
\end{exoresolu}

\begin{aretenir}[Rappel essentiel — Section 1]
\begin{itemize}
\item Lieux : \eng{nearby store, provision store, market stall, supermarket, roadside kiosk}.
\item Acteurs : \eng{customer / buyer} vs \eng{seller / trader / shopkeeper / stallholder}.
\item Articles : \eng{soap, matches, candles, batteries, notebooks, pens, sugar, salt}.
\item Foodstuff : \eng{rice, beans, gari, plantains, tomatoes, onions, groundnut oil, palm oil, dried fish, meat, eggs, bread, flour}.
\item Mesures : \eng{a bag of, a kilo of, a litre of, a bottle of, a tin of, a packet of, a bunch of, a cup of}.
\item Argent : \eng{price, cost, discount, change, FCFA}.
\item Question prix : \eng{How much is it?} (jamais \eng{How many is it?}).
\end{itemize}
\end{aretenir}

\coursec{Functional language: asking for and giving information}

\courssub{Asking for information (Customer)}

To buy efficiently, the customer uses specific question patterns. Politeness markers (\eng{please}, \eng{could you}, \eng{would you}) are essential in Cameroonian anglophone culture.

\begin{propriete}[Structures clés pour demander]
\begin{itemize}
\item \textbf{Prix :} \eng{How much is / are …?} | \eng{How much does … cost?} | \eng{What is the price of …?}
\item \textbf{Quantité / Disponibilité :} \eng{Do you have …?} | \eng{Have you got …?} | \eng{Is there any … left?}
\item \textbf{Qualité / Choix :} \eng{Which one is better?} | \eng{Can I see …?} | \eng{Are these fresh?}
\item \textbf{Préférence :} \eng{I'd like …} (I would like) | \eng{I want …} (plus direct) | \eng{Can I have …?}
\end{itemize}
\end{propriete}

\begin{exemplebox}[Exemples traduits — Client]
\eng{How much is a kilo of onions?} « Combien coûte le kilo d'oignons ? » \\
\eng{Do you have any fresh tomatoes?} « Avez-vous des tomates fraîches ? » \\
\eng{I'd like a bag of rice, please.} « Je voudrais un sac de riz, s'il vous plaît. » \\
\eng{Can I have two litres of groundnut oil?} « Puis-je avoir deux litres d'huile d'arachide ? » \\
\eng{Are the plantains ripe?} « Les plantains sont-ils mûrs ? »
\end{exemplebox}

\courssub{Giving information (Seller)}

The seller answers, proposes, advises and closes the sale. The tone must be welcoming and professional.

\begin{propriete}[Structures clés pour informer / vendre]
\begin{itemize}
\item \textbf{Annoncer le prix :} \eng{It is / They are … FCFA.} | \eng{It costs …} | \eng{That comes to …}
\item \textbf{Proposer / Conseiller :} \eng{I recommend …} | \eng{This one is better.} | \eng{We have a special offer on …}
\item \textbf{Quantifier :} \eng{We sell by the kilo / bag / tin.} | \eng{One bag contains 50 kilos.}
\item \textbf{Politesse commerciale :} \eng{Here you are.} | \eng{Anything else?} | \eng{Thank you, come again.}
\end{itemize}
\end{propriete}

\begin{exemplebox}[Exemples traduits — Vendeur]
\eng{A bag of beans is 20,000 FCFA.} « Un sac de haricots coûte 20 000 FCFA. » \\
\eng{These plantains are very ripe and sweet.} « Ces plantains sont bien mûrs et sucrés. » \\
\eng{We sell the oil by the litre.} « On vend l'huile au litre. » \\
\eng{Here is your change. Anything else?} « Voici votre monnaie. Autre chose ? » \\
\eng{Thank you, Madam. Have a nice day!} « Merci, Madame. Bonne journée ! »
\end{exemplebox}

\begin{attention}[Erreurs fréquentes — Francophones]
\begin{itemize}
\item \textbf{“I want…” seul} : Trop brutal. Préférez \eng{“I'd like…”} ou \eng{“Can I have…?”}.
\item \textbf{“How much cost?”} : Incorrect. Dire \eng{“How much does it cost?”} (auxiliaire \eng{does}).
\item \textbf{“Give me…” sans “please”} : Impoli. Dire \eng{“Give me…, please”} ou mieux \eng{“Could I have…, please?”}.
\item \textbf{Confusion “price / prize”} : \eng{Price} = prix ; \eng{Prize} = prix (récompense). Ne dites pas \eng{“What is the prize?”}.
\end{itemize}
\end{attention}

\courssub{Practice: Matching questions and answers}

\begin{exercice}[Matching]
Associez chaque question du client (1--6) à la réponse appropriée du vendeur (A--F).
\begin{enumerate}
\item How much is a tin of tomatoes?
\item Do you have any dried fish?
\item I'd like a bunch of plantains.
\item Are these onions fresh?
\item Can I have a bag of rice?
\item What is your last price for the oil?
\end{enumerate}
\begin{itemize}
\item[A] Yes, they were harvested this morning.
\item[B] It is 600 FCFA.
\item[C] Here you are. That will be 25,000 FCFA.
\item[D] Yes, we have smoked fish at 2,000 FCFA the kilo.
\item[E] My final price is 1,400 FCFA per litre.
\item[F] Certainly. This bunch is 1,000 FCFA.
\end{itemize}
\end{exercice}

\begin{corrige}[Exercice Matching]
1--B ; 2--D ; 3--F ; 4--A ; 5--C ; 6--E.
\end{corrige}

\coursec{Bargaining, relaunching and concluding the exchange}

\courssub{Bargaining (Marchander)}

Bargaining is standard in Cameroonian markets (Mokolo, Mfoundi, Bamenda Main Market) but \textbf{not} in supermarkets or modern provision stores where prices are fixed (\eng{fixed prices}).

\begin{propriete}[Étapes et formules de la négociation]
\begin{enumerate}
\item \textbf{Exprimer que c'est cher :} \eng{That's too expensive.} | \eng{It's a bit pricey.} | \eng{That's more than I expected.}
\item \textbf{Proposer un prix / Demander une remise :} \eng{Can you give me a discount?} | \eng{Will you take … FCFA?} | \eng{What about … FCFA?} | \eng{Can you reduce the price?}
\item \textbf{Argumenter :} \eng{I'm a regular customer.} | \eng{I'm buying in bulk.} | \eng{The other stall sells it cheaper.}
\item \textbf{Accepter / Refuser (Vendeur) :} \eng{Okay, for you: … FCFA.} (Accepter) | \eng{I'm afraid I can't go lower.} (Refuser poliment) | \eng{That's my last price.} (Prix final)
\end{enumerate}
\end{propriete}

\begin{exemplebox}[Mini-dialogue de négociation]
\eng{Customer:} How much for this bunch of plantains? \\
\eng{Seller:} 1,500 FCFA. \\
\eng{Customer:} \textbf{That's too expensive. Will you take 1,200 FCFA?} \\
\eng{Seller:} \textbf{I can do 1,300 FCFA. That's my last price.} \\
\eng{Customer:} \textbf{Okay, agreed. I'll take it.}
\end{exemplebox}

\courssub{Relaunching (Relancer la conversation)}

To keep the interaction natural or add items, both speakers use relaunching phrases.

\begin{center}
\begin{tabularx}{\linewidth}{|X|X|}
\hline
\rowcolor{popblueL} \textbf{Fonction} & \textbf{Formules utiles} \\
\hline
Client : ajouter un article & \eng{Also, I need…} | \eng{And what about…?} | \eng{Do you also sell…?} \\
Vendeur : proposer plus & \eng{Would you like some… with that?} | \eng{We have fresh… today.} | \eng{Anything else?} \\
Vérifier / Confirmer & \eng{So that's one bag of rice and two litres of oil?} | \eng{Is that everything?} \\
\hline
\end{tabularx}
\end{center}

\courssub{Concluding the exchange (Conclure la transaction)}

The final phase involves payment, giving change and leave-taking formulas.

\begin{propriete}[Séquence de conclusion type]
\begin{enumerate}
\item \textbf{Récapitulatif :} \eng{That makes 5,000 FCFA altogether.} / \eng{The total comes to 5,000 FCFA.}
\item \textbf{Paiement :} \eng{Here is 10,000 FCFA.} (Client tend l'argent)
\item \textbf{Rendre la monnaie :} \eng{Here is your change: 5,000 FCFA.} / \eng{Here are 5,000 FCFA change.}
\item \textbf{Formules de séparation :} \eng{Thank you. / Thanks a lot.} | \eng{Have a nice day!} | \eng{Goodbye! / Bye!} | \eng{Come again!} (Vendeur)
\end{enumerate}
\end{propriete}

\begin{attention}[Vocabulaire argent : “Change” vs “Exchange”]
\begin{itemize}
\item \cle{Change} (noun, uncountable) = la monnaie qu'on vous rend. \eng{“Here is your change.”}
\item \cle{Exchange} (noun/verb) = le change (bureau de change) / échanger. \eng{“The exchange rate is 650 FCFA for 1 USD.”}
\item Ne dites pas \eng{“Here is your exchange”} pour la monnaie.
\end{itemize}
\end{attention}

\begin{savaistu}[Culture du “Dash” / Pourboire]
Dans les marchés camerounais, le vendeur offre parfois un petit plus (un oignon, une tomate, un morceau de sucre) après l'achat : c'est le \eng{dash} (argot nigérian/camerounais) ou \eng{a little something extra}. Le client répond : \eng{“Thank you, Madam! God bless you!”} Cela scelle la relation commerciale amicale.
\end{savaistu}

\coursec{Model dialogues}

Three typical situations. Read them in pairs, then swap roles.

\courssub{Dialogue 1: At a provision store (Fixed prices, polite)}

\begin{textebox}[Dialogue 1 — Provision Store, Buea]
Shopkeeper: Good afternoon. How can I help you?
Customer: Good afternoon. I'd like to buy some provisions for the week.
Shopkeeper: Certainly. What do you need?
Customer: A bag of rice, two packets of spaghetti and a tin of milk.
Shopkeeper: The rice is 25,000 FCFA, spaghetti is 600 FCFA a packet, milk is 800 FCFA a tin.
Customer: Okay. That's 25,000 plus 1,200 plus 800… 27,000 FCFA altogether?
Shopkeeper: That's correct. 27,000 FCFA.
Customer: Here is 30,000 FCFA.
Shopkeeper: Thank you. Here is your change: 3,000 FCFA. Here are your goods.
Customer: Thank you very much. Goodbye.
Shopkeeper: Goodbye, Sir. Come again.
\end{textebox}

\courssub{Dialogue 2: At a market stall (Bargaining, quantities)}

\begin{textebox}[Dialogue 2 — Mokolo Market Stall]
Stallholder: Welcome, my sister! Buy your tomatoes, very fresh today!
Customer: Good morning. Are these tomatoes really fresh?
Stallholder: Yes, brought from the farm this morning. How many kilos?
Customer: I want three kilos. How much per kilo?
Stallholder: 800 FCFA per kilo.
Customer: \textbf{800? That's too much. The other stall sells at 700.}
Stallholder: \textbf{My sister, mine are better quality. Let's say 750 FCFA, last price.}
Customer: \textbf{Okay, 750 FCFA. Give me three kilos.}
Stallholder: Here you are. That makes 2,250 FCFA.
Customer: Here is 3,000 FCFA.
Stallholder: Your change: 750 FCFA. Thank you! Come again!
Customer: Thank you. Bye.
\end{textebox}

\courssub{Dialogue 3: At a roadside kiosk (Quick, small items)}

\begin{textebox}[Dialogue 3 — Roadside Kiosk, Douala]
Customer: (Arrives on motorbike) Good evening. Do you have cold Malta?
Kiosk keeper: Yes, small bottle or large?
Customer: Large, please. And a recharge card for MTN, 500 FCFA.
Kiosk keeper: Here is the Malta. And the card. That's 800 FCFA altogether.
Customer: (Hands 1,000 FCFA) Here.
Kiosk keeper: Change: 200 FCFA. Enjoy your drink!
Customer: Thanks. Bye.
Kiosk keeper: Bye bye.
\end{textebox}

\courssub{Analyse fonctionnelle des dialogues}

\begin{center}
\begin{tabularx}{\linewidth}{|X|X|X|}
\hline
\rowcolor{popblueL} \textbf{Fonction} & \textbf{Dialogue 1 (Boutique)} & \textbf{Dialogue 2 (Marché)} \\
\hline
Salutation & Good afternoon / How can I help? & Welcome, my sister! \\
Demande prix & Implicite (liste donnée) & How much per kilo? \\
Négociation & Aucune (prix fixes) & That's too much / Last price \\
Récapitulatif prix & That's 27,000 FCFA altogether & That makes 2,250 FCFA \\
Paiement / Monnaie & Here is 30,000 / Here is your change & Here is 3,000 / Your change 750 \\
Congé & Goodbye / Come again & Thank you / Bye / Come again \\
\hline
\end{tabularx}
\end{center}

\coursec{Pronunciation, stress and intonation}

\courssub{Word stress: Two-syllable nouns vs verbs}

Many shopping words change stress depending on grammatical category. \textbf{Noun/Adjective = stress on 1st syllable} | \textbf{Verb = stress on 2nd syllable}.

\begin{center}
\begin{tabularx}{\linewidth}{|X|X|X|X|}
\hline
\rowcolor{popblueL} \textbf{Mot} & \textbf{Nom / Adjectif} & \textbf{Verbe} & \textbf{Prononciation (approx. FR)} \\
\hline
record & \eng{RE-cord} (un disque) & \eng{re-CORD} (enregistrer) & « RÉ-kord » / « ri-KORD » \\
present & \eng{PRE-sent} (cadeau) & \eng{pre-SENT} (présenter) & « PRÉ-zent » / « pri-ZENT » \\
discount & \eng{DIS-count} (remise) & \eng{dis-COUNT} (réduire) & « DIS-kaount » / « dis-KAOUNT » \\
produce & \eng{PRO-duce} (produits) & \eng{pro-DUCE} (produire) & « PRO-dyouce » / « pro-DYOUCE » \\
refund & \eng{RE-fund} (remboursement) & \eng{re-FUND} (rembourser) & « RÉ-fonde » / « ri-FONDE » \\
\hline
\end{tabularx}
\end{center}

\begin{exemplebox}[Entraînement]
\eng{The \textbf{DIS}count is 10\%.} (Nom) \\
\eng{They \textbf{dis}COUNT the price for regulars.} (Verbe) \\
\eng{Fresh \textbf{PRO}duce arrives Monday.} (Nom) \\
\eng{This region \textbf{pro}DUCES great coffee.} (Verbe)
\end{exemplebox}

\courssub{Sentence stress and intonation patterns}

In English, \cle{content words} (nouns, main verbs, adjectives, question words) are stressed. \cle{Function words} (articles, auxiliaries, prepositions, pronouns) are weak (schwa /ə/).

\begin{propriete}[Règles d'intonation pour l'achat-vente]
\begin{itemize}
\item \textbf{Questions WH- (How much, What, Which) :} Intonation \textbf{descendante} (\eng{How MUCH is it?} ↘).
\item \textbf{Questions Oui/Non (Do you, Can I, Is it) :} Intonation \textbf{montante} (\eng{Do you HAVE any?} ↗).
\item \textbf{Demandes polies (Could I, Would you) :} Intonation \textbf{montante-descendante} (courbe polie) (\eng{Could I HAVE some?} ↗↘).
\item \textbf{Affirmations / Prix :} Intonation \textbf{descendante} (\eng{It's 5,000 FCFA.} ↘).
\item \textbf{Listes (énumération) :} Montante sur chaque élément, descendante sur le dernier (\eng{Rice, beans, oil, and bread.} ↗ ↗ ↗ ↘).
\end{itemize}
\end{propriete}

\courssub{Key sounds for shopping vocabulary}

Practice these minimal pairs and difficult sounds. Prononciation notée entre guillemets « … » à la française.

\begin{center}
\begin{tabularx}{\linewidth}{|X|X|X|}
\hline
\rowcolor{popblueL} \textbf{Son} & \textbf{Mots clés (achat-vente)} & \textbf{Prononciation FR} \\
\hline
/θ/ (th non voisé) & \eng{think, three, thousand, nothing} & « sink », « sri », « zaou-zend », « na-zing » (langue entre dents) \\
/ð/ (th voisé) & \eng{this, that, these, those, other} & « zis », « zat », « ziz », « zoz », « o-zè » (langue entre dents + voix) \\
/ʃ/ (ch) & \eng{shop, cash, fish, sugar, dish} & « chop », « kèch », « fich », « chou-gue », « dich » \\
/tʃ/ (tch) & \eng{change, cheap, choose, catch, kitchen} & « tchéindj », « tchip », « tchouze », « kètch », « kit-chène » \\
/dʒ/ (dj) & \eng{just, large, orange, vegetables, budget} & « djeust », « lar-dje », « o-ran-dje », « vè-dje-ta-beuls », « beu-djèt » \\
/ɒ/ (o court) & \eng{shop, cost, lot, bottle, onion} & « chop », « coste », « lot », « bo-tle », « o-gnon » \\
/ɜː/ (eur long) & \eng{work, girl, purchase, merchant, serve} & « weurk », « geurl », « per-tchisse », « meur-tchant », « seurve » \\
\hline
\end{tabularx}
\end{center}

\begin{methode}[Comment travailler la prononciation seul]
\begin{enumerate}
\item \textbf{Écoutez} le modèle (professeur, audio OnBuch+).
\item \textbf{Répétez} en chœur, puis seul (murmure, puis voix haute).
\item \textbf{Enregistrez-vous} sur votre téléphone ; comparez.
\item \textbf{Exagérez} l'intonation (montée/descente) au début.
\item \textbf{Ciblez} 3 sons difficiles par séance (ex. /θ/, /ʃ/, /dʒ/).
\end{enumerate}
\end{methode}

\begin{exercice}[Prononciation — Lecture expressive]
Lisez les phrases suivantes en respectant l'accentuation de phrase (mots en \textbf{gras} = forts) et l'intonation (flèche).
\begin{enumerate}
\item \eng{\textbf{HOW MUCH} is a \textbf{KILO} of \textbf{MEAT}?} ↘
\item \eng{\textbf{DO} you \textbf{HAVE} any \textbf{FRESH} \textbf{TOMATOES}?} ↗
\item \eng{\textbf{I'd LIKE} a \textbf{BAG} of \textbf{RICE}, \textbf{PLEASE}.} ↗↘
\item \eng{\textbf{THAT'S} \textbf{TOO EXPENSIVE}. \textbf{CAN} you \textbf{GIVE} a \textbf{DISCOUNT}?} ↘
\item \eng{\textbf{HERE} is \textbf{10,000} FCFA. \textbf{THANK YOU}.} ↘
\end{enumerate}
\end{exercice}

\begin{corrige}[Exercice Prononciation]
Vérifiez : 1. WH-question descendante. 2. Yes/No montante. 3. Demande polie courbe. 4. Plainte + demande : descendante sur plainte, montante sur demande. 5. Affirmation + politesse : descendante.
\end{corrige}

\coursec{Guided role-plays: let's speak!}

\courssub{Instructions for role-plays}

\begin{methode}[Déroulement d'un jeu de rôle guidé]
\begin{enumerate}
\item \textbf{Préparation (2 min) :} Lisez votre carte de rôle. Ne la montrez pas à votre partenaire. Préparez mentalement le vocabulaire (conteneurs, prix, politesse).
\item \textbf{Jeu (3 min) :} Jouez la scène. Le professeur circule, note les points forts / faiblesses (vocabulaire, grammaire, prononciation, fluidité, politesse).
\item \textbf{Échange des rôles :} Rejouez en inversant Client / Vendeur.
\item \textbf{Retour collectif (5 min) :} Le professeur corrige 2--3 erreurs fréquentes au tableau. Deux binômes volontaires rejouent devant la classe.
\end{enumerate}
\end{methode}

\courssub{Role-play cards (Cartes à photocopier / afficher)}

\begin{experience}[Situation 1 : Ingrédients pour le Ndolé (Marché)]
\begin{tabularx}{\linewidth}{|X|X|}
\hline
\rowcolor{poporangeL} \textbf{Student A : Customer} & \textbf{Student B : Stallholder} \\
\hline
Vous préparez le Ndolé pour dimanche. & Vous vendez : feuilles de ndolé (botte), \\
Achetez : 2 bottes de ndolé, 500g crevettes & crevettes séchées (bol), arachides (bol), \\
séchées, 1 bol d'arachides, 2 cubes Maggi. & cubes Maggi (unité). Prix libres (ex: 500, \\
Négociez le prix total. Payez 5,000 FCFA. & 1000, 300, 50 FCFA). Donnez la monnaie. \\
\hline
\end{tabularx}
\end{experience}

\begin{experience}[Situation 2 : Fournitures scolaires (Kiosque / Boutique)]
\begin{tabularx}{\linewidth}{|X|X|}
\hline
\rowcolor{popgreenL} \textbf{Student A : Customer} & \textbf{Student B : Shopkeeper} \\
\hline
Rentrée des classes. Achetez : & Prix fixes (pas de négociation). \\
3 cahiers (200 p.), 2 stylos (100 p.), & Cahier 200, Stylo 100, Règle 50, \\
1 règle, 1 boîte crayons de couleur. & Crayons 500. Calculez le total. Rendez la \\
Demandez le prix total. Payez 2,000 FCFA. & monnaie. Emballez dans un sachet. \\
\hline
\end{tabularx}
\end{experience}

\begin{experience}[Situation 3 : Crédit téléphone \& Boisson (Kiosque bord de route)]
\begin{tabularx}{\linewidth}{|X|X|}
\hline
\rowcolor{popblueL} \textbf{Student A : Customer (à moto)} & \textbf{Student B : Kiosk keeper} \\
\hline
Pressé. Voulez : 1 grande bouteille Coca & Vendez : Coca 500, Malta 500, Eau 300. \\
(500 FCFA) + Recharge MTN 1000 FCFA. & Recharge MTN/Orange 500, 1000, 2000. \\
Demandez si Coca est frais. Payez 2000 FCFA. & Vérifiez code recharge. Rendez monnaie. \\
\hline
\end{tabularx}
\end{experience}

\begin{experience}[Situation 4 : Réclamation \& Échange (Supermarché / Boutique)]
\begin{tabularx}{\linewidth}{|X|X|}
\hline
\rowcolor{poppurpleL} \textbf{Student A : Customer} & \textbf{Student B : Shopkeeper / Manager} \\
\hline
Acheté hier : 1 sac riz 25kg (25,000 FCFA). & Le client revient avec sac (pas ouvert ?) ou \\
En rentrant : sac déchiré / riz mouillé / & ticket. Politique : échange si défaut, pas \\
cailloux dedans. Revenez pour échange ou & remboursement cash. Proposez échange même \\
remboursement. Restez poli mais ferme. & marque ou marque supérieure + différence. \\
\hline
\end{tabularx}
\end{experience}

\courssub{Evaluation grid (Grille d'auto-évaluation / prof)}

After each role-play, tick the boxes.

\begin{center}
\begin{tabularx}{\linewidth}{|X|c|c|c|c|}
\hline
\rowcolor{popblueL} \textbf{Critère} & \textbf{4 Excellent} & \textbf{3 Bien} & \textbf{2 Passable} & \textbf{1 Insuffisant} \\
\hline
Vocabulaire (articles, mesures, argent) & Précis, varié & Correct, suffisant & Limité, répétitif & Très pauvre, erreurs \\
Grammaire (How much/many, would like, prix) & Sans erreur & 1--2 erreurs mineures & Erreurs fréquentes & Bloquant \\
Prononciation / Intonation & Naturelle, expressive & Compréhensible, quelques fautes & Accent fort, incompréhensible par moments & Incompréhensible \\
Fluidité / Interaction & Spontané, relance bien & Réactif, quelques pauses & Longues pauses, monosyllabes & Ne parle presque pas \\
Politesse / Conventions sociales & Parfait (please, thanks, titles) & Correct & Oubli fréquent (pas de please) & Brutal, impoli \\
\hline
\end{tabularx}
\end{center}

\begin{aretenir}[Synthèse de la leçon — Speaking: Buying and Selling]
\begin{itemize}
\item \textbf{Lieux \& Acteurs} : \eng{store, stall, kiosk, supermarket} ; \eng{customer, stallholder, shopkeeper}.
\item \textbf{Lexique clé} : \eng{foodstuff, provisions, containers (bag, kilo, litre, tin, bunch, cup)}.
\item \textbf{Fonctions} : Demander prix (\eng{How much}), Quantité (\eng{Do you have}), Négocier (\eng{Too expensive, Last price, Discount}), Conclure (\eng{Change, Come again}).
\item \textbf{Grammaire} : \eng{How much} (indénombrable) vs \eng{How many} (pluriel) ; \eng{I'd like} (politesse) ; \eng{This/These} (proximité).
\item \textbf{Prononciation} : Accent de mot (noun/verb), Intonation WH- (↘) vs Yes/No (↗), Sons /θ/, /ð/, /ʃ/, /tʃ/, /dʒ/.
\item \textbf{Culture} : Négociation au marché (pas au supermarché), \eng{dash} (petit cadeau), titres de politesse (\eng{Madam, Sir, My sister/brother}).
\end{itemize}
\end{aretenir}

\begin{exercice}[Écrit : Composition guidée — “My Saturday at the Market”]
Rédigez un paragraphe de 80--100 mots en anglais. Décrivez votre dernier passage au marché (Mokolo, Mfoundi, ou votre marché local). Utilisez au moins : \textbf{5 articles/foodstuff}, \textbf{3 mesures/conteneurs}, \textbf{2 formules de négociation}, \textbf{les mots price, change, discount}.
\textit{Indication :} Utilisez le passé simple (\eng{I went, I bought, I asked, The seller gave…}).
\end{exercice}

\begin{corrige}[Exercice Composition — Modèle de réponse]
\eng{Last Saturday, I went to Mfoundi market to buy foodstuff for the week. I needed a bag of rice, two kilos of tomatoes, a litre of groundnut oil and a bunch of plantains. At the stall, I asked, “How much is the rice?” The seller said 25,000 FCFA. I said, “That's too expensive, can you give me a discount?” He agreed on 24,000 FCFA. The tomatoes were 800 FCFA per kilo. I paid 30,000 FCFA and received my change. The seller was kind and gave me a dash of onions. I said “Thank you, come again!” and went home happy.}
\end{corrige}

\coursec{Functional language: asking for and giving information}

Dans la section précédente, nous avons découvert les lieux et les personnages de notre scène : le \eng{store} (la boutique du quartier) et le \eng{market} (le marché), le \eng{customer} (le client) et le \eng{seller} (le vendeur). Nous savons maintenant \emph{où} nous sommes et \emph{qui} parle. Il reste l'essentiel : \emph{que dire} ? Cette section vous donne toutes les formules fonctionnelles — les « briques de langue » toutes prêtes — pour demander une information et la donner, du premier \eng{Good morning!} jusqu'à la question sur la fraîcheur du poisson. Observez d'abord ce mini-dialogue, puis nous démonterons chaque formule.

\begin{textebox}[At the provision store — a short exchange]
\textbf{Seller:} Good morning, Madam! Can I help you?\\
\textbf{Customer:} Good morning. I'm looking for some rice and palm oil. Do you have any?\\
\textbf{Seller:} Yes, we have plenty. How much rice do you need?\\
\textbf{Customer:} Two kilos, please. How much is it?\\
\textbf{Seller:} It's 600 FCFA a kilo. And the palm oil is 1,000 FCFA a bottle.\\
\textbf{Customer:} Is the oil pure?\\
\textbf{Seller:} Yes, it arrived this morning. It's the best quality in the market.\\
\textbf{Customer:} Fine. I'll take two kilos of rice and one bottle of oil. Thank you!\\
\textbf{Seller:} You're welcome, Madam. Come again!
\end{textebox}

\textbf{Glossaire :} \emph{to look for} (v., chercher) ; \emph{plenty} (adv., beaucoup, en abondance) ; \emph{pure} (adj., pur, non mélangé) ; \emph{to arrive} (v., arriver) ; \emph{quality} (n., qualité) ; \emph{You're welcome} (expr., je vous en prie) ; \emph{Come again!} (expr., revenez nous voir !).

\textbf{Questions de compréhension :} 1. What does the customer want to buy? 2. How much is a kilo of rice? 3. Why is the customer sure the oil is good?

\begin{corrige}[Questions de compréhension]
1. \eng{She wants to buy some rice and palm oil.} « Elle veut acheter du riz et de l'huile de palme. »\\
2. \eng{It is 600 FCFA a kilo.} « C'est 600 FCFA le kilo. »\\
3. \eng{Because the seller says it arrived this morning and that it is the best quality in the market.} « Parce que le vendeur dit qu'elle est arrivée ce matin et que c'est la meilleure qualité du marché. »
\end{corrige}

\courssub{Opening the conversation: the seller greets and offers help}

En anglais, c'est presque toujours le \textbf{vendeur} qui ouvre la conversation. Trois formules dominent :

\begin{exemplebox}[The seller opens the conversation]
\eng{Good morning! Can I help you?} « Bonjour ! Je peux vous aider ? » (formule la plus courante, partout)\\
\eng{What can I do for you?} « Que puis-je faire pour vous ? » (un peu plus formel)\\
\eng{Are you being served?} « Est-ce qu'on s'occupe de vous ? » (très britannique ; littéralement : « êtes-vous en train d'être servi ? »)
\end{exemplebox}

Notez la construction de la dernière formule : \eng{are you being served} est un présent continu \emph{passif} (\eng{being} + participe passé). Inutile de la produire vous-même à l'oral en Première, mais il faut la \emph{comprendre} : elle signifie simplement « quelqu'un vous sert-il déjà ? ».

\begin{attention}[Ne traduisez pas mot à mot]
Un francophone qui entend \eng{Can I help you?} répond parfois \eng{Yes, you can.} — erreur ! La question n'attend pas une réponse grammaticale mais l'expression de votre besoin : répondez directement \eng{Yes, please. I need some sugar.} De même, \eng{What can I do for you?} ne veut pas dire « qu'est-ce que tu peux faire pour moi ? » au sens littéral ; c'est une simple offre de service.
\end{attention}

\courssub{Expressing your needs: the customer speaks}

Pour dire ce que vous cherchez, l'anglais offre cinq formules, de la plus directe à la plus interrogative :

\begin{center}
\begin{tabularx}{\linewidth}{|X|X|X|}
\hline
\rowcolor{popblueL} English & Français & Registre / remarque \\
\hline
\eng{I want to buy some bread.} & Je veux acheter du pain. & direct, un peu sec \\
\hline
\eng{I need some sugar.} & J'ai besoin de sucre. & direct, courant \\
\hline
\eng{I'm looking for dried fish.} & Je cherche du poisson fumé. & poli, très naturel \\
\hline
\eng{Do you have any onions?} & Avez-vous des oignons ? & question, très poli \\
\hline
\eng{Have you got any onions?} & Avez-vous des oignons ? & identique, britannique \\
\hline
\end{tabularx}
\end{center}

\begin{propriete}[Some ou any ?]
\cle{any} s'emploie dans les \textbf{questions} et les \textbf{phrases négatives} : \eng{Do you have any sugar?} / \eng{We don't have any sugar.}\\
\cle{some} s'emploie dans les \textbf{phrases affirmatives} et dans les \textbf{offres} : \eng{We have some sugar.} / \eng{Would you like some?}\\
Exception utile : dans une offre ou une demande où l'on attend « oui », on garde \eng{some} même à la forme interrogative : \eng{Could I have some water, please?}
\end{propriete}

\courssub{Asking for prices and quantities}

\begin{exemplebox}[Asking the price]
\eng{How much is it?} « Combien ça coûte ? »\\
\eng{How much are the tomatoes?} « Combien coûtent les tomates ? »\\
\eng{How much is a kilo of meat?} « Combien coûte le kilo de viande ? »\\
\eng{What's the price of this palm oil?} « Quel est le prix de cette huile de palme ? »
\end{exemplebox}

\begin{propriete}[How much is / How much are / How many]
\eng{How much + is} + nom \textbf{singulier ou indénombrable} : \eng{How much is the oil?} « Combien coûte l'huile ? »\\
\eng{How much + are} + nom \textbf{pluriel} : \eng{How much are the eggs?} « Combien coûtent les œufs ? »\\
\eng{How many} + nom \textbf{pluriel dénombrable} pour les \emph{quantités} : \eng{How many eggs do you want?} « Combien d'œufs voulez-vous ? »\\
Rappel : en anglais, \eng{rice, sugar, oil, meat, bread, money} sont \textbf{indénombrables} → toujours avec \eng{how much} et le verbe au singulier.
\end{propriete}

Pour la quantité, le vendeur demande : \eng{How many do you want?} (dénombrable) ou \eng{How much rice do you need?} (indénombrable). Le client répond simplement : \eng{Two kilos, please.} « Deux kilos, s'il vous plaît. » / \eng{Just a small packet.} « Juste un petit paquet. »

\begin{exemplebox}[Giving information — the seller answers]
\eng{It's 500 FCFA a kilo.} « C'est 500 FCFA le kilo. »\\
\eng{They cost 250 FCFA each.} « Elles coûtent 250 FCFA pièce. »\\
\eng{We sell it at 1,000 FCFA a bottle.} « Nous le vendons à 1 000 FCFA la bouteille. »\\
\eng{I'm sorry, we don't have any left.} « Je suis désolé, il n'en reste plus. »\\
\eng{It's out of stock.} « C'est en rupture de stock. »
\end{exemplebox}

Notez l'article indéfini \eng{a} qui, devant une unité, correspond au français « le/la » : \eng{500 FCFA a kilo} = « 500 FCFA \emph{le} kilo ». Et \eng{each} = « pièce, chacun ».

\courssub{Asking about quality and freshness}

\begin{exemplebox}[Quality and freshness]
\eng{Is it fresh?} « C'est frais ? » — \eng{Yes, it arrived this morning.} « Oui, c'est arrivé ce matin. »\\
\eng{Is this palm oil pure?} « Cette huile de palme est-elle pure ? » — \eng{It's the best quality in the market.} « C'est la meilleure qualité du marché. »\\
\eng{Do you have any dried fish?} « Avez-vous du poisson fumé ? » — \eng{Yes, we have plenty.} « Oui, nous en avons beaucoup. »\\
\eng{How much is a tin of tomatoes?} « Combien coûte une boîte de tomates ? » — \eng{It's 350 FCFA.} « C'est 350 FCFA. »
\end{exemplebox}

\begin{popfigure}
\begin{center}
\begin{tikzpicture}[
  box/.style={rounded corners=3pt, draw=popline, thick, align=center, font=\small, inner sep=5pt, minimum width=5.6cm},
  q/.style={box, fill=popblueL},
  a/.style={box, fill=poporangeL},
  lab/.style={font=\small\bfseries, text=popdark}]
\node[lab, text=popblue] (t1) at (0,0) {The CUSTOMER asks};
\node[lab, text=poporange] (t2) at (7.6,0) {The SELLER answers};
\node[q] (q1) at (0,-1.1) {Do you have any rice?};
\node[a] (a1) at (7.6,-1.1) {Yes, we have plenty.};
\node[q] (q2) at (0,-2.3) {How much is it a kilo?};
\node[a] (a2) at (7.6,-2.3) {It's 600 FCFA a kilo.};
\node[q] (q3) at (0,-3.5) {Is it fresh?};
\node[a] (a3) at (7.6,-3.5) {Yes, it arrived this morning.};
\node[q] (q4) at (0,-4.7) {Do you have any palm oil?};
\node[a] (a4) at (7.6,-4.7) {I'm sorry, it's out of stock.};
\foreach \i in {1,2,3,4}{\draw[-{Stealth[length=3mm]}, thick, poppurple] (q\i.east) -- (a\i.west);}
\end{tikzpicture}
\end{center}
\legende{L'échange question → réponse entre le client et le vendeur.}
\end{popfigure}

\begin{methode}[Demander poliment en anglais]
\begin{enumerate}
\item \textbf{Ouvrez} toujours par \eng{Excuse me} (pour attirer l'attention) ou saluez : \eng{Good morning!}
\item \textbf{Formulez} votre demande avec \eng{please} : \eng{Two kilos of rice, please.}
\item \textbf{Préférez} les formes interrogatives (\eng{Do you have...?} / \eng{Could I have...?}) aux formes directes (\eng{I want...}), plus sèches.
\item \textbf{Terminez} par \eng{Thank you} ou \eng{Thanks a lot}.
\end{enumerate}
La politesse anglaise passe par le \emph{ton} autant que par les mots : une voix montante et souriante vaut tous les \eng{please} du monde.
\end{methode}

\begin{exoresolu}[Put the exchange in order]
Énoncé — Remettez les répliques suivantes dans l'ordre logique d'un dialogue au marché :\\
a) \eng{It's 500 FCFA a kilo.} \quad b) \eng{Good morning! Can I help you?} \quad c) \eng{Two kilos, please.} \quad d) \eng{Yes, I need some meat. How much is it?} \quad e) \eng{How much do you want?} \quad f) \eng{Thank you very much!}
\tcblower
Solution détaillée — Un échange commercial suit toujours la même logique : salutation et offre d'aide → expression du besoin → question sur le prix → réponse du vendeur → question sur la quantité → commande → remerciement. L'ordre correct est donc : \textbf{b → d → a → e → c → f}.\\
\eng{b) Good morning! Can I help you? — d) Yes, I need some meat. How much is it? — a) It's 500 FCFA a kilo. — e) How much do you want? — c) Two kilos, please. — f) Thank you very much!}\\
Remarquez que la réplique (e) ne peut venir qu'après le prix : le vendeur ne demande la quantité qu'une fois le produit choisi.
\end{exoresolu}

\begin{exercice}[Your turn: complete the exchange]
Complétez ce dialogue avec les formules de la section.\\
\textbf{Seller:} Good afternoon! \ldots\ldots\ldots\ldots\ldots\ldots ?\\
\textbf{Customer:} Yes, please. \ldots\ldots\ldots\ldots\ldots\ldots (chercher) some groundnuts. Do you have any?\\
\textbf{Seller:} Yes, we do. \ldots\ldots\ldots\ldots\ldots\ldots do you want?\\
\textbf{Customer:} One kilo. \ldots\ldots\ldots\ldots\ldots\ldots is it?\\
\textbf{Seller:} It's 400 FCFA a kilo.\\
\textbf{Customer:} Fine, I'll take it. \ldots\ldots\ldots\ldots\ldots\ldots !\\
\textbf{Seller:} You're welcome. Come again!
\end{exercice}

\begin{corrige}[Exercice : complete the exchange]
\eng{Good afternoon! Can I help you? — Yes, please. I'm looking for some groundnuts. Do you have any? — Yes, we do. How much do you want? — One kilo. How much is it? — It's 400 FCFA a kilo. — Fine, I'll take it. Thank you! — You're welcome. Come again!}\\
Note : \eng{groundnuts} est indénombrable dans ce contexte (une denrée en vrac), d'où \eng{How much do you want?} et non \eng{How many}.
\end{corrige}

\begin{attention}[Erreurs fréquentes des francophones]
\begin{itemize}
\item \eng{How much are the rice?} est faux : \eng{rice} est indénombrable → \eng{How much \textbf{is} the rice?}
\item Ne dites pas \eng{I want that you give me...} ; dites \eng{Could you give me...?} ou \eng{Can I have...?}
\item \eng{How many costs...?} n'existe pas : pour le prix, c'est toujours \eng{How much}.
\item Faux ami : \eng{actually} ne veut pas dire « actuellement » mais « en fait ». Pour « actuellement », dites \eng{at the moment} : \eng{We don't have any at the moment.}
\item Attention à l'ordre des mots dans la question : \eng{How much is it?} et non \eng{How much it is?} (inversion obligatoire).
\end{itemize}
\end{attention}

\begin{savaistu}[“Are you all right there?”]
En Grande-Bretagne, le vendeur dit souvent \eng{Are you all right there?} ou simplement \eng{You OK?} pour proposer son aide. Ce n'est \emph{pas} une question sur votre santé ! La réponse attendue est \eng{Yes, thanks, I'm just looking} (« oui, merci, je regarde seulement ») si vous ne voulez rien acheter, ou directement votre demande. Au Cameroun anglophone, on entend aussi le chaleureux \eng{Welcome, my customer!} — une formule toute locale qui met tout de suite à l'aise.
\end{savaistu}

\coursec{Setting the scene: at the store and at the market}

Cette leçon vous plonge dans une situation de communication authentique : \textbf{acheter et vendre des articles ou denrées alimentaires dans un magasin ou un marché de proximité}. Au Cameroun, comme dans tout pays anglophone, cette interaction suit un rituel précis : saluer, s'informer (prix, qualité, quantité), négocier (au marché), décider, payer, éventuellement relancer, et se quitter poliment. Maîtriser ce « script » oral est indispensable pour l'autonomie quotidienne et pour l'examen oral du Baccalauréat.

\begin{textebox}[Scène de marché : négociation et conclusion]
Customer: Good morning, Madam. How much is a kilo of tomatoes?
Trader: Good morning, my dear. It is 1,000 FCFA per kilo. Very fresh, brought from the West Region this morning.
Customer: One thousand? That is too expensive! The other seller sells them at 800 FCFA.
Trader: Ah, but Madam, look at the colour! These are \eng{Grade A}. The others are \eng{Grade B}. But because you are a regular customer, I can give you a discount. Let's say 900 FCFA.
Customer: Can you make it 850 FCFA if I take two kilos?
Trader: \eng{Hmm...} All right, two kilos for 1,700 FCFA. That is my last price.
Customer: OK, I'll take them. Here is 2,000 FCFA.
Trader: Thank you. Here is your change: 300 FCFA. Anything else?
Customer: What about some fresh pepper?
Trader: Yes, I have yellow and green. 200 FCFA the heap.
Customer: I'll take one heap of green pepper. Here is the money.
Trader: Here you are. Thank you very much. Come again!
Customer: You're welcome. Goodbye!
\end{textebox}

\paragraph{Glossaire}
\begin{itemize}
    \item \textbf{Grade A / Grade B} (n.) : qualité supérieure / qualité inférieure.
    \item \textbf{regular customer} (n.) : client habituel / fidèle.
    \item \textbf{discount} (n.) : réduction, remise.
    \item \textbf{last price} (n.) : prix final, prix net (non négociable).
    \item \textbf{heap} (n.) : tas, petit tas (pour légumes, fruits).
    \item \textbf{change} (n.) : monnaie (rendue après paiement).
    \item \textbf{Madam / Sir} (n.) : Madame / Monsieur (formules de politesse, capitalisées en vocatif).
\end{itemize}

\coursec{Functional language: asking for and giving information}

Avant de négocier, il faut obtenir l'information. Cette section regroupe les formules fonctionnelles pour \emph{demander} (prix, quantité, qualité, origine) et \emph{donner} ces informations.

\courssub{Demander le prix, la quantité, la qualité}

\begin{definition}[Formules pour s'informer (Asking for information)]
\begin{tabularx}{\linewidth}{|X|X|}\hline
\rowcolor{popblueL} \textbf{Fonction} & \textbf{Formules types (Anglais)} \\ \hline
Prix (général) & \eng{How much is...? / How much does... cost?} \\ \hline
Prix (quantité précise) & \eng{How much is a kilo of...? / What is the price per kilo?} \\ \hline
Qualité / Catégorie & \eng{What quality is it? / Are they Grade A? / Are they fresh?} \\ \hline
Origine / Provenance & \eng{Where do they come from? / Are they local?} \\ \hline
Quantité disponible & \eng{How many do you have? / Do you have enough for...?} \\ \hline
\end{tabularx}
\end{definition}

\courssub{Donner l'information (Giving information)}

\begin{propriete}[Structures de réponse du vendeur]
\begin{itemize}
    \item \textbf{Prix :} \eng{It is [prix] per kilo / each / per item.} (\eng{It's 500 FCFA per kilo.})
    \item \textbf{Qualité :} \eng{They are Grade A / fresh / organic / just arrived.}
    \item \textbf{Origine :} \eng{They come from the West Region / They are local / imported.}
    \item \textbf{Disponibilité :} \eng{I have plenty. / Only a few left. / Sold out.}
\end{itemize}
\textbf{Rappel grammaire :} \eng{It is} $\rightarrow$ \eng{It's} (oral) ; \eng{They are} $\rightarrow$ \eng{They're}. \eng{Per} = « le / par » (distribution). \eng{Each} = « l'unité / pièce ».
\end{propriete}

\begin{exemplebox}[Échange informationnel : exemples traduits]
\begin{tabularx}{\linewidth}{|X|X|}\hline
\rowcolor{popblueL} \textbf{Anglais} & \textbf{Français} \\ \hline
\eng{Customer: How much are the bananas?} & Client : Combien coûtent les bananes ? \\ \hline
\eng{Seller: They are 300 FCFA per kilo.} & Vendeur : Elles sont à 300 FCFA le kilo. \\ \hline
\eng{Customer: Are they ripe?} & Client : Sont-elles mûres ? \\ \hline
\eng{Seller: Yes, very ripe and sweet.} & Vendeur : Oui, très mûres et sucrées. \\ \hline
\eng{Customer: Where do they come from?} & Client : D'où viennent-elles ? \\ \hline
\eng{Seller: From Buea, Madam. Local produce.} & Vendeur : De Buéa, Madame. Produit local. \\ \hline
\end{tabularx}
\end{exemplebox}

\coursec{Bargaining, relaunching and concluding the exchange}

Une fois l'information obtenue, la transaction entre dans sa phase dynamique : la \textbf{négociation} (bargaining), la \emph{relance} pour augmenter le panier d'achat, et enfin la \textbf{conclusion} (paiement, monnaie, au revoir).

\courssub{Relancer la conversation : proposer un article supplémentaire (Relaunching)}

Quand le client hésite ou a fini son achat principal, le vendeur \emph{relance} pour augmenter la vente (\eng{upselling}) ou par politesse commerciale. En anglais, on utilise des structures ouvertes, non agressives.

\begin{definition}[Formules de relance (Relaunching)]
Ces formules invitent le client à continuer ses achats sans pression.
\begin{itemize}
    \item \eng{Anything else?} — Court, standard, très fréquent.
    \item \eng{What else do you need?} — Plus direct, orienté « besoin ».
    \item \eng{Would you like something else?} — Poli, conditionnel de politesse (\eng{would like}).
    \item \eng{What about some fresh eggs? / How about some fresh eggs?} — Proposition ciblée (suggestion d'un produit précis).
    \item \eng{We also have... today.} — Information nouvelle pour déclencher l'achat d'appoint.
\end{itemize}
\end{definition}

\begin{methode}[Comment relancer naturellement (Technique de vente)]
\begin{enumerate}
    \item \textbf{Attendez} que le client ait validé l'article principal (\eng{OK, I'll take it}).
    \item \textbf{Utilisez} \eng{Anything else?} pour une relance générale OU \eng{What about...?} pour suggérer un produit précis (souvent un complément : poivre pour la tomate, pain pour le beurre).
    \item \textbf{Évitez} l'impératif \eng{Buy this!} qui est agressif. Préférez le conditionnel \eng{Would you like...?} ou la question ouverte \eng{What about...?}.
\end{enumerate}
\end{methode}

\begin{exemplebox}[Relance : exemples traduits]
\begin{tabularx}{\linewidth}{|X|X|}\hline
\rowcolor{popblueL} \textbf{Anglais} & \textbf{Français} \\ \hline
\eng{Anything else, Sir?} & Et avec ceci, Monsieur ? / Autre chose ? \\ \hline
\eng{Would you like some bread to go with the butter?} & Voulez-vous du pain pour accompagner le beurre ? \\ \hline
\eng{What about a bottle of juice?} & Et si vous preniez une bouteille de jus ? \\ \hline
\eng{We also have fresh fish today.} & Nous avons aussi du poisson frais aujourd'hui. \\ \hline
\end{tabularx}
\end{exemplebox}

\courssub{Réagir à un prix jugé trop élevé}

Le client exprime son désaccord. En anglais, l'adjectif standard est \cle{expensive}. \cle{Dear} existe en anglais britannique (soutenu), mais \eng{expensive} est universel.

\begin{propriete}[Exprimer « c'est trop cher »]
\begin{itemize}
    \item \eng{That's too expensive!} (Standard, neutre).
    \item \eng{That's too much for me.} (Subjectif : « pour mon budget »).
    \item \eng{The price is too high.} (Factuel, un peu plus formel).
    \item \eng{It costs an arm and a leg.} (Idiome : « ça coûte un bras et une jambe » = très cher).
\end{itemize}
\textbf{Attention :} On ne dit pas \eng{*It is too much expensive} (pléonasme syntaxique). \eng{Too} modifie l'adjectif \eng{expensive}, pas le nom \eng{price} directement dans cette structure.
\end{propriete}

\begin{attention}[Erreur fréquente : \eng{*I will reflect} vs \eng{I'll think about it}]
\begin{itemize}
    \item \textbf{Faux :} \eng{*I will reflect.} (Signifie « je vais méditer/réfléchir profondément », style littéraire).
    \item \textbf{Vrai :} \eng{I'll think about it.} (Décision d'achat différée, courant).
    \item \textbf{Vrai :} \eng{I'll have a think.} (Familier, très britannique).
\end{itemize}
De même, \eng{dear} (cher) s'emploie surtout en UK (\eng{It's a bit dear}), mais \eng{expensive} marche partout.
\end{attention}

\courssub{Négocier poliment (Bargaining / Haggling)}

La négociation suit un code de politesse. Le client propose, le vendeur contre-propose ou refuse. L'usage du \textbf{modal \eng{can} / \eng{could}} est central pour la politesse.

\begin{propriete}[Structures de la négociation (Client)]
\begin{tabularx}{\linewidth}{|X|X|X|}\hline
\rowcolor{popblueL} \textbf{Fonction} & \textbf{Structure} & \textbf{Exemple} \\ \hline
Demander une réduction & \eng{Can / Could you reduce the price?} & \eng{Could you reduce the price a little?} \\ \hline
Demander une remise & \eng{Can you give me a discount?} & \eng{Can you give me a discount for two kilos?} \\ \hline
Demander le prix final & \eng{What is your last price?} & \eng{What's your last price for this bag?} \\ \hline
Proposer un deal (conditionnel) & \eng{If I buy X, how much...?} & \eng{If I buy three, how much will I pay?} \\ \hline
Proposer un prix précis & \eng{I'll give you [prix] for it.} & \eng{I'll give you 500 FCFA for the lot.} \\ \hline
\end{tabularx}
\end{propriete}

\begin{propriete}[Réponses du vendeur (Trader / Shopkeeper)]
\begin{tabularx}{\linewidth}{|X|X|}\hline
\rowcolor{poporangeL} \textbf{Situation} & \textbf{Réponse type} \\ \hline
Prix fixe (supermarché, prix affiché) & \eng{Sorry, the price is fixed.} / \eng{No discount, Madam, fixed price.} \\ \hline
Petite concession & \eng{OK, I can give you a small discount.} \\ \hline
Concession pour fidélité & \eng{Because you are a good customer, take it for 900 FCFA.} \\ \hline
Prix final non négociable & \eng{That's my last price.} / \eng{I can't go lower than that.} \\ \hline
\end{tabularx}
\end{propriete}

\begin{savaistu}[Le \eng{Haggling} au Nigeria et au Ghana]
Dans les grands marchés de Lagos (Nigeria) ou d'Accra (Ghana), comme au Marché Central de Douala ou au Marché de Bamenda, la négociation (\eng{haggling} ou \eng{bargaining}) est un \textbf{art social}. Le premier prix annoncé (\eng{asking price}) est presque \emph{toujours} supérieur de 20 à 50\% au prix final (\eng{final price}). Ne pas marchander peut passer pour de la naïveté ou de l'arrogance. Cependant, dans les supermarchés modernes (Shoprite, Carrefour, Mahima), les prix sont \eng{fixed} (fixes) : on ne négocie pas.
\end{savaistu}

\courssub{Conclure l'achat : décision, paiement et monnaie}

Une fois le prix arrêté, le client prend la décision d'achat. C'est ici qu'intervient une règle grammaticale cruciale : l'usage de \eng{will} pour la décision \emph{immédiate}.

\begin{propriete}[Règle : \eng{I'll take it} vs \eng{*I take it}]
\begin{itemize}
    \item \textbf{Correct :} \eng{I'll take it.} (\eng{I will take it} contracté).
    \item \textbf{Incorrect :} \eng{*I take it.}
    \item \textbf{Explication :} \eng{Will} (futur simple) exprime ici une \textbf{décision spontanée} prise au moment de parler (\eng{instant decision}). Le présent simple (\eng{I take}) exprime une habitude ou une vérité générale, pas une décision d'achat ponctuelle.
    \item \textbf{Variantes :} \eng{All right, I'll take it.} / \eng{OK, give me two kilos, please.} / \eng{I'll have this one.}
\end{itemize}
\end{propriete}

\begin{exemplebox}[Paiement et monnaie : scripts types]
\begin{tabularx}{\linewidth}{|X|X|}\hline
\rowcolor{popgreenL} \textbf{Anglais} & \textbf{Français} \\ \hline
\eng{Here is the money.} & Voici l'argent. \\ \hline
\eng{Here is 5,000 FCFA.} & Voici 5 000 FCFA. \\ \hline
\eng{Here is your change: 500 FCFA.} & Voici votre monnaie : 500 FCFA. \\ \hline
\eng{Keep the change.} & Gardez la monnaie (pourboire). \\ \hline
\eng{Can I have a receipt, please?} & Puis-je avoir un reçu, s'il vous plaît ? \\ \hline
\eng{Do you have change for 10,000?} & Avez-vous de la monnaie pour 10 000 ? \\ \hline
\end{tabularx}
\end{exemplebox}

\courssub{Refuser poliment et prendre congé}

Savoir partir sans acheter est une compétence sociale. Le client garde la face, le vendeur garde le client potentiel.

\begin{definition}[Refus poli et formules de départ]
\begin{tabularx}{\linewidth}{|X|X|}\hline
\rowcolor{poppurpleL} \textbf{Client (Refus)} & \textbf{Vendeur (Au revoir)} \\ \hline
\eng{No, thank you.} & \eng{You're welcome. Come again!} \\ \hline
\eng{I'll think about it.} & \eng{Take your time. Have a nice day!} \\ \hline
\eng{Maybe next time.} & \eng{We'll be here. Goodbye!} \\ \hline
\eng{Not today, thanks.} & \eng{Thank you for coming. Bye!} \\ \hline
\end{tabularx}
\end{definition}

\begin{attention}[Pièges de traduction : « Au revoir »]
\begin{itemize}
    \item \eng{Goodbye} est définitif ou formel.
    \item \eng{Bye} / \eng{Bye-bye} est standard, amical.
    \item \eng{See you later} / \eng{See you soon} implique qu'on se reverra (marché hebdomadaire).
    \item \textbf{Ne dites pas} \eng{*Go well} (traduction littérale de « va bien » / \eng{Allez bien}) pour « au revoir ».
\end{itemize}
\end{attention}

\begin{popfigure}
\begin{tikzpicture}[
    node distance=1.2cm and 0.5cm,
    every node/.style={font=\small, align=center},
    stepnode/.style={rectangle, rounded corners, draw=popblue, thick, fill=popblueL, minimum width=3.2cm, minimum height=1.4cm},
    arrownode/.style={-Stealth, thick, popdark}
]
% Nœuds
\node[stepnode, align=center] (greet){\textbf{1. Greeting}\\\eng{Good morning! / Hello!}};
\node[stepnode, below=of greet, align=center] (ask){\textbf{2. Asking for Info}\\\eng{How much? / What quality?}};
\node[stepnode, below=of ask, align=center] (bargain){\textbf{3. Bargaining}\\\eng{Too expensive! / Discount? / Last price?}};
\node[stepnode, below=of bargain, align=center] (pay){\textbf{4. Paying}\\\eng{I'll take it. / Here is money. / Change.}};
\node[stepnode, below=of pay, align=center] (bye){\textbf{5. Saying Goodbye}\\\eng{Thank you! / Come again!}};

% Flèches principales
\draw[arrownode] (greet) -- (ask);
\draw[arrownode] (ask) -- (bargain);
\draw[arrownode] (bargain) -- (pay);
\draw[arrownode] (pay) -- (bye);

% Boucle optionnelle (relance) - pointillés
\draw[arrownode, dashed, poporange] (pay.east) -- ++(1.5,0) |- (ask.east) node[pos=0.25, right, fill=popcream, rounded corners, inner sep=2pt, font=\tiny, text=popdark] {Relance: \eng{Anything else?}};
\end{tikzpicture}
\legende{Organigramme de la transaction commerciale en 5 étapes (la flèche pointillée orange montre la relance possible après paiement).}
\end{popfigure}

\courssub{Récapitulatif lexical : du bonjour au revoir}

\begin{center}
\begin{tabularx}{\linewidth}{|X|X|X|}\hline
\rowcolor{popblueL} \textbf{Étape} & \textbf{Client (Customer)} & \textbf{Vendeur (Trader / Shopkeeper)} \\ \hline
\textbf{1. Saluer} & \eng{Good morning. / Hello.} & \eng{Good morning, Madam/Sir. Welcome!} \\ \hline
\textbf{2. S'informer} & \eng{How much is...? / Price of...?} & \eng{It's 500 FCFA. / 1,000 per kilo.} \\ \hline
\textbf{3. Réagir / Négocier} & \eng{Too expensive! / Discount? / Last price?} & \eng{Fixed price. / OK, 900 for you. / Last price.} \\ \hline
\textbf{4. Décider / Payer} & \eng{I'll take it. / Here is 2,000.} & \eng{Here is your change. / Receipt?} \\ \hline
\textbf{5. Relance (optionnel)} & \eng{No thanks. / What about...?} & \eng{Anything else? / What about...?} \\ \hline
\textbf{6. Quitter} & \eng{Thank you. Bye!} & \eng{Come again! / Have a nice day!} \\ \hline
\end{tabularx}
\end{center}

\coursec{Pronunciation, stress and intonation}

La prononciation, l'accentuation (stress) et l'intonation sont essentielles pour être compris et paraître poli. En anglais, ce ne sont pas les syllabes qui comptent (comme en français) mais les \textbf{mots porteurs de sens} (mots de contenu) qui sont accentués et plus longs.

\courssub{Accent de mot (Word Stress) — Vocabulaire clé}

Les mots ci-dessous sont fréquents dans la transaction. L'accent tonique est indiqué en \textbf{gras} et la prononciation approximative en lettres françaises entre guillemets.

\begin{center}
\begin{tabularx}{\linewidth}{|X|X|X|}\hline
\rowcolor{popblueL} \textbf{Mot} & \textbf{Accentuation (IPA simplifiée)} & \textbf{Prononciation (lettres fr.)} \\ \hline
expensive & \eng{ex\textbf{PEN}sive} /ɪkˈspensɪv/ & « ik-\textbf{pen}-siv » \\ \hline
discount (n.) & \eng{\textbf{DIS}count} /ˈdɪskaʊnt/ & « \textbf{dis}-kaount » \\ \hline
discount (v.) & \eng{dis\textbf{COUNT}} /dɪsˈkaʊnt/ & « dis-\textbf{kaount} » \\ \hline
tomato & \eng{to\textbf{MA}to} /təˈmɑːtəʊ/ & « to-\textbf{ma}-to » \\ \hline
banana & \eng{ba\textbf{NA}na} /bəˈnɑːnə/ & « ba-\textbf{na}-na » \\ \hline
kilo & \eng{\textbf{KI}lo} /ˈkiːləʊ/ & « \textbf{ki}-lo » \\ \hline
receipt & \eng{re\textbf{CEIPT}} /rɪˈsiːt/ & « ri-\textbf{sit} » (le \eng{p} est muet) \\ \hline
change & \eng{\textbf{CHANGE}} /tʃeɪndʒ/ & « \textbf{tchéindj} » \\ \hline
customer & \eng{\textbf{CUS}tomer} /ˈkʌstəmə/ & « \textbf{cus}-to-meu » \\ \hline
bargain & \eng{\textbf{BAR}gain} /ˈbɑːɡɪn/ & « \textbf{ba}-guine » \\ \hline
\end{tabularx}
\end{center}

\courssub{Accent de phrase (Sentence Stress) et Intonation}

Dans une phrase, on accentue les \textbf{mots de contenu} (noms, verbes pleins, adjectifs, adverbes, mots interrogatifs) et on « avale » les mots grammaticaux (articles, auxiliaires, prépositions, pronoms).

\begin{propriete}[Règles d'intonation pour le dialogue d'achat]
\begin{itemize}
    \item \textbf{Questions fermées (Yes/No) :} Intonation \textbf{montante} final (\eng{/ \nearrow/}).
    \exemple{ \eng{Can you reduce the price?} « \textbf{Can you re-DUCE the PRICE} \nearrow » }
    \item \textbf{Questions ouvertes (Wh-) :} Intonation \textbf{descendante} final (\eng{/ \searrow/}).
    \exemple{ \eng{How much is it?} « \textbf{HOW MUCH is IT} \searrow » }
    \item \textbf{Affirmations / Décisions :} Intonation \textbf{descendante}.
    \exemple{ \eng{I'll take it.} « \textbf{I'll TAKE IT} \searrow » }
    \item \textbf{Politesse (Would you like... / Could you...) :} Intonation \textbf{montante-descendante} (courbe en cloche) ou haute montante pour marquer la déférence.
    \exemple{ \eng{Would you like a bag?} « \textbf{WOULD you LIKE a BAG} \nearrow\searrow » }
\end{itemize}
\end{propriete}

\begin{methode}[Entraînement à la prosodie (3 étapes)]
\begin{enumerate}
    \item \textbf{Identifiez} les mots de contenu (surlignez-les).
    \item \textbf{Tapez} du pied ou frappez dans les mains \emph{uniquement} sur les mots de contenu en lisant la phrase.
    \item \textbf{Dites} la phrase en allongeant les voyelles des mots accentués et en respectant la courbe d'intonation.
\end{enumerate}
\end{methode}

\begin{exemplebox}[Exemples d'accentuation de phrase (mots en \textbf{gras} = accentués)]
\begin{tabularx}{\linewidth}{|X|X|}\hline
\rowcolor{popgreenL} \textbf{Phrase} & \textbf{Intonation} \\ \hline
\eng{\textbf{HOW} \textbf{MUCH} \textbf{IS} \textbf{IT}?} & Descendante (Wh-) \\ \hline
\eng{\textbf{CAN} you \textbf{GIVE} me a \textbf{DISCOUNT}?} & Montante (Yes/No) \\ \hline
\eng{\textbf{THAT'S} \textbf{TOO} \textbf{EXPENSIVE}!} & Descendante (Exclamation) \\ \hline
\eng{\textbf{I'LL} \textbf{TAKE} \textbf{IT}.} & Descendante (Décision) \\ \hline
\eng{\textbf{WOULD} you \textbf{LIKE} \textbf{ANYTHING} \textbf{ELSE}?} & Montante (Politesse/Offre) \\ \hline
\end{tabularx}
\end{exemplebox}

\coursec{Model dialogues}

Deux dialogues modèles complets illustrant la différence entre \textbf{supermarché (prix fixes)} et \textbf{marché (négociation)}.

\begin{textebox}[Modèle 1 : Supermarché — Prix fixes (Fixed prices)]
\eng{Shop Assistant: Good morning, Madam. Can I help you?}
\eng{Customer: Good morning. Yes, how much are the apples?}
\eng{Shop Assistant: The imported apples are 2,500 FCFA per kilo. The local bananas are 300 FCFA per kilo.}
\eng{Customer: 2,500? That's too expensive! Can you give me a discount?}
\eng{Shop Assistant: I'm sorry, Madam, the price is fixed in the supermarket. But the bananas are on special offer today.}
\eng{Customer: All right. I'll take two kilos of bananas. Here is 1,000 FCFA.}
\eng{Shop Assistant: Here is your change: 400 FCFA. Anything else?}
\eng{Customer: No, thank you. Goodbye.}
\eng{Shop Assistant: You're welcome. Have a nice day!}
\end{textebox}

\begin{textebox}[Modèle 2 : Marché Central — Négociation autorisée (Bargaining)]
\eng{Trader: Welcome, my brother! Beautiful tomatoes today, fresh fish from Kribi!}
\eng{Customer: Good morning. What is the price for the tomatoes and that big fish?}
\eng{Trader: Tomatoes 1,200 FCFA per kilo. Fish 3,000 FCFA each.}
\eng{Customer: Eh! Too high! The man next door sells tomatoes at 1,000. I want 5 kilos and two fish. What is your last price?}
\eng{Trader: My brother, my tomatoes are Grade A! But for you... 5,000 FCFA for the tomatoes and 5,000 for the two fish. Total 10,000 FCFA. Last price.}
\eng{Customer: OK, I'll take them. Here is 10,000 FCFA.}
\eng{Trader: Thank you. Here is your receipt. Anything else? Mangoes are sweet today.}
\eng{Customer: I'll think about the mangoes. Thank you. Goodbye!}
\eng{Trader: Come again! Have a nice day!}
\end{textebox}

\coursec{Guided role-plays: let's speak!}

\courssub{Exercice guidé 1 : Reconstitution de dialogue (Examen type)}

\begin{exoresolu}[Mettre en ordre et compléter le dialogue]
\textbf{Énoncé :} Remettez les répliques dans l'ordre logique (1 à 8). Complétez les trous avec \eng{will / would / can / too / fixed / change}.

\begin{enumerate}
    \item Trader: Good morning! \_\_\_\_\_\_ I help you?
    \item Customer: How much \_\_\_\_\_\_\_ this shirt?
    \item Trader: It is 5,000 FCFA.
    \item Customer: 5,000! That's \_\_\_\_\_\_ expensive. \_\_\_\_\_\_ you reduce the price?
    \item Trader: Sorry, the price is \_\_\_\_\_\_.
    \item Customer: OK, I \_\_\_\_\_\_ take it then. Here is 5,000 FCFA.
    \item Trader: Thank you. Here is your \_\_\_\_\_\_. Anything else?
    \item Customer: No, thanks. Goodbye!
\end{enumerate}

\tcblower
\textbf{Solution détaillée :}
\begin{enumerate}
    \item \textbf{Can} (ou \textbf{Could} plus poli) : Offre d'aide standard. \eng{Can I help you?}
    \item \textbf{is} (verbe \eng{be}, pas un modal) : \eng{How much is this shirt?}
    \item (Donnée)
    \item \textbf{too} (adverbe de degré devant adjectif) / \textbf{Can} (ou \textbf{Could}) : \eng{That's too expensive. Can you reduce the price?}
    \item \textbf{fixed} (adjectif) : \eng{The price is fixed.}
    \item \textbf{will} (décision immédiate) : \eng{I'll take it.} (\eng{I will} $\rightarrow$ \eng{I'll}).
    \item \textbf{change} (nom, monnaie) : \eng{Here is your change.}
    \item (Donnée)
\end{enumerate}
\textbf{Dialogue reconstitué :}
\eng{Trader: Good morning! Can I help you?}
\eng{Customer: How much is this shirt?}
\eng{Trader: It is 5,000 FCFA.}
\eng{Customer: 5,000! That's too expensive. Can you reduce the price?}
\eng{Trader: Sorry, the price is fixed.}
\eng{Customer: OK, I'll take it then. Here is 5,000 FCFA.}
\eng{Trader: Thank you. Here is your change. Anything else?}
\eng{Customer: No, thanks. Goodbye!}
\end{exoresolu}

\courssub{Exercice guidé 2 : Scénarios de jeu de rôle (Évaluation orale)}

\begin{exercice}[Entraînement oral : Scénarios de marché]
\textbf{Consigne :} En binôme, jouez les deux scénarios ci-dessous. Alternez les rôles (Client / Vendeur). Utilisez le vocabulaire et les structures de cette leçon. Le professeur écoutera la prononciation (section 5) et l'intonation (montante pour questions fermées, descendante pour questions ouvertes et affirmations).

\textbf{Scénario A : Au rayon fruits (Supermarché — Prix fixes).}
\begin{itemize}
    \item Client : Demande le prix des pommes (importées) et des bananes (locales).
    \item Vendeur : Annonce les prix (ex. 2 500 FCFA/kg pommes, 300 FCFA/kg bananes).
    \item Client : Trouve les pommes chères, demande une remise.
    \item Vendeur : Refuse poliment (prix fixe), propose les bananes en promo.
    \item Client : Achète 2 kg de bananes, paie 1 000 FCFA, récupère monnaie.
    \item Fin : Formules de politesse.
\end{itemize}

\textbf{Scénario B : Au marché Central (Négociation autorisée).}
\begin{itemize}
    \item Vendeur : Accueille, propose « beautiful tomatoes, fresh fish ».
    \item Client : Demande prix tomates (1 200 FCFA/kg) et poisson (3 000 FCFA/pièce).
    \item Client : Négocie fermement (compare avec concurrent, achat gros volume : 5 kg tomates + 2 poissons).
    \item Vendeur : Résiste, puis cède « last price » global.
    \item Client : Accepte (\eng{I'll take them}), paie gros billet (10 000 FCFA), vérifie monnaie.
    \item Vendeur : Relance (\eng{Anything else?}), client refuse (\eng{I'll think about it for the mangoes}).
    \item Fin : Au revoir chaleureux.
\end{itemize}
\end{exercice}

\begin{corrige}[Exercice : Scénarios de marché — Productions attendues types]
\textbf{Scénario A (Extrait type) :}
\eng{Customer: Good morning. How much are the apples?}
\eng{Seller: The apples are 2,500 FCFA per kilo, Madam. The bananas are 300 FCFA.}
\eng{Customer: 2,500? That's too expensive! Can you give me a discount?}
\eng{Seller: I'm sorry, the price is fixed in the supermarket. But the bananas are on special offer today.}
\eng{Customer: All right. I'll take two kilos of bananas. Here is 1,000 FCFA.}
\eng{Seller: Here is your change: 400 FCFA. Anything else?}
\eng{Customer: No, thank you. Goodbye.}
\eng{Seller: You're welcome. Have a nice day!}

\textbf{Scénario B (Extrait type) :}
\eng{Trader: Welcome, my brother! Beautiful tomatoes today, fresh fish from Kribi!}
\eng{Customer: Good morning. Price for tomatoes and that big fish?}
\eng{Trader: Tomatoes 1,200 FCFA per kilo. Fish 3,000 FCFA each.}
\eng{Customer: Eh! Too high! The man next door sells tomatoes at 1,000. I want 5 kilos and two fish. What is your last price?}
\eng{Trader: My brother, my tomatoes are Grade A! But for you... 5,000 FCFA for the tomatoes and 5,000 for the two fish. Total 10,000 FCFA. Last price.}
\eng{Customer: OK, I'll take them. Here is 10,000 FCFA.}
\eng{Trader: Thank you. Here is your receipt. Anything else? Mangoes are sweet today.}
\eng{Customer: I'll think about the mangoes. Thank you. Goodbye!}
\eng{Trader: Come again! Have a nice day!}
\end{corrige}

\coursec{Model dialogues}

Dans la section précédente, nous avons appris à \cle{négocier} (\eng{bargain}), à \cle{relancer} le vendeur (\eng{What is your last price?}) et à \cle{conclure} l'échange (\eng{I'll take it.}). Il est maintenant temps de rassembler toutes ces formules dans des \textbf{dialogues complets}, du premier mot au dernier. C'est en mémorisant et en jouant ces dialogues modèles que tu deviendras capable de mener toi-même une conversation réelle au marché de Mokolo ou dans une boutique de ton quartier.

\courssub{Dialogue 1: At the market — buying foodstuff}

\begin{textebox}[Dialogue 1 — At the market stall]
\textbf{Customer:} Good morning, Madam!\\
\textbf{Trader:} Good morning! What can I do for you?\\
\textbf{Customer:} I would like to buy some tomatoes and onions. How much are the tomatoes?\\
\textbf{Trader:} They are 500 FCFA a kilo.\\
\textbf{Customer:} Are they fresh?\\
\textbf{Trader:} Yes, they arrived this morning.\\
\textbf{Customer:} Right, give me two kilos of tomatoes and one kilo of onions, please.\\
\textbf{Trader:} Anything else?\\
\textbf{Customer:} No, thank you. How much is everything?\\
\textbf{Trader:} 1,300 FCFA, please.\\
\textbf{Customer:} Here is 1,500 FCFA.\\
\textbf{Trader:} Here is your change: 200 FCFA. Thank you, come again!
\end{textebox}

\textbf{Consigne de lecture :} lis ce dialogue à voix haute. Souligne en \textbf{bleu} les formules de \emph{demande d'information} (\eng{How much are the tomatoes? Are they fresh?}) et en \textbf{rouge} les formules de \emph{conclusion} (\eng{Here is 1,500 FCFA. Thank you, come again!}).

\begin{exemplebox}[Traduction du dialogue 1]
— Bonjour, Madame ! — Bonjour ! Que puis-je faire pour vous ? — Je voudrais acheter des tomates et des oignons. Combien coûtent les tomates ? — Elles sont à 500 FCFA le kilo. — Sont-elles fraîches ? — Oui, elles sont arrivées ce matin. — D'accord, donnez-moi deux kilos de tomates et un kilo d'oignons, s'il vous plaît. — Autre chose ? — Non, merci. Combien coûte le tout ? — 1 300 FCFA, s'il vous plaît. — Voici 1 500 FCFA. — Voici votre monnaie : 200 FCFA. Merci, revenez nous voir !
\end{exemplebox}

\begin{propriete}[Les formules-clés du dialogue 1]
\begin{itemize}
\item \eng{What can I do for you?} = « Que puis-je faire pour vous ? » — formule d'accueil du vendeur.
\item \eng{I would like to buy...} / \eng{I'd like to buy...} = « Je voudrais acheter... » — plus poli que \eng{I want to buy}.
\item \eng{How much are the tomatoes?} = « Combien coûtent les tomates ? » — \eng{how much} + verbe \eng{be} au pluriel car \eng{tomatoes} est pluriel.
\item \eng{Anything else?} = « Autre chose ? » — question raccourcie de \eng{Do you want anything else?}
\item \eng{Here is your change.} = « Voici votre monnaie. » — \eng{change} (la monnaie rendue) est invariable et singulier.
\end{itemize}
\end{propriete}

\courssub{Dialogue 2: At the shop — buying articles}

\begin{textebox}[Dialogue 2 — At the provision store]
\textbf{Shopkeeper:} Good afternoon! Can I help you?\\
\textbf{Customer:} Yes, please. Do you have any batteries and candles?\\
\textbf{Shopkeeper:} Yes, we do. The batteries are 250 FCFA each and the candles 100 FCFA each.\\
\textbf{Customer:} I'd like four batteries and a packet of candles, please.\\
\textbf{Shopkeeper:} Here you are. Anything else?\\
\textbf{Customer:} How much is a bar of soap?\\
\textbf{Shopkeeper:} It's 400 FCFA.\\
\textbf{Customer:} That's a bit expensive. Can you reduce the price?\\
\textbf{Shopkeeper:} Sorry, our prices are fixed.\\
\textbf{Customer:} All right, I'll take one. How much do I owe you?\\
\textbf{Shopkeeper:} That makes 2,000 FCFA, please.\\
\textbf{Customer:} Here you are. Thank you!\\
\textbf{Shopkeeper:} You're welcome. Have a nice day!
\end{textebox}

\begin{exemplebox}[Formules de paiement traduites]
\eng{How much do I owe you?} = « Combien vous dois-je ? »\\
\eng{Here you are.} = « Voilà. » (en tendant l'argent ou l'article)\\
\eng{That makes 2,000 FCFA.} = « Cela fait 2 000 FCFA. »\\
\eng{Our prices are fixed.} = « Nos prix sont fixes. » (pas de négociation possible)\\
\eng{I'll take one.} = « J'en prends un. » (décision d'achat)
\end{exemplebox}

\begin{attention}[Here you are ou Here we are ?]
\eng{Here you are} signifie « voilà » quand on \textbf{donne} quelque chose à quelqu'un : \eng{Here you are, your batteries.} (Voilà, vos piles.)\\
\eng{Here we are} signifie « nous voilà (arrivés) » : \eng{Here we are, at the market!} (Nous voilà au marché !)\\
Ne confonds jamais les deux : dire \eng{Here we are} en tendant de l'argent au vendeur serait comique !
\end{attention}

\courssub{Dialogue 3: Bargaining at the market}

\begin{textebox}[Dialogue 3 — Bargaining for dried fish]
\textbf{Customer:} How much do you sell this dried fish \textbf{for}?\\
\textbf{Trader:} 2,000 FCFA.\\
\textbf{Customer:} Oh! That's too much. What is your last price?\\
\textbf{Trader:} For you, Madam, 1,800 FCFA.\\
\textbf{Customer:} If I buy two, how much will I pay?\\
\textbf{Trader:} Two for 3,500 FCFA.\\
\textbf{Customer:} OK, I'll take two.
\end{textebox}

\begin{exemplebox}[Traduction du dialogue 3]
— Combien vendez-vous ce poisson fumé ? — 2 000 FCFA. — Oh ! C'est trop cher. Quel est votre dernier prix ? — Pour vous, Madame, 1 800 FCFA. — Si j'en achète deux, combien vais-je payer ? — Deux pour 3 500 FCFA. — D'accord, j'en prends deux.
\end{exemplebox}

\begin{propriete}[La négociation en trois coups]
\begin{enumerate}
\item \textbf{Protester poliment :} \eng{That's too much!} / \eng{That's a bit expensive.} (C'est trop cher !)
\item \textbf{Demander un effort :} \eng{What is your last price?} / \eng{Can you reduce the price?} (Quel est votre dernier prix ?)
\item \textbf{Négocier la quantité :} \eng{If I buy two, how much will I pay?} (Si j'en prends deux, combien je paie ?) — note la structure : \eng{if} + présent, \eng{will} dans l'autre partie de la phrase (conditionnelle de type 1).
\end{enumerate}
\end{propriete}

\courssub{Guided analysis: the five steps of a shopping dialogue}

Observe les trois dialogues : ils suivent tous la \textbf{même architecture en cinq étapes}. C'est le « squelette » de tout échange commercial en anglais.

\begin{center}
\begin{tabularx}{\linewidth}{|l|X|X|}
\hline
\rowcolor{popblueL} \textbf{Étape} & \textbf{Fonction} & \textbf{Formules modèles} \\
\hline
1. Greeting & Saluer et accueillir & \eng{Good morning! Can I help you? / What can I do for you?} \\
\hline
2. Asking for information & Demander prix et qualité & \eng{How much are the tomatoes? Are they fresh? Do you have any batteries?} \\
\hline
3. Bargaining & Négocier le prix & \eng{That's too much. What is your last price? Can you reduce the price?} \\
\hline
4. Paying & Commander et payer & \eng{I'd like two kilos, please. How much do I owe you? Here you are. Here is your change.} \\
\hline
5. Leave-taking & Prendre congé & \eng{Thank you! You're welcome. Have a nice day! Come again!} \\
\hline
\end{tabularx}
\end{center}

\begin{popfigure}
\begin{tikzpicture}[node distance=0.9cm]
\node[draw=popblue, fill=popblueL, rounded corners, minimum width=2.1cm, minimum height=1cm, align=center, font=\small\bfseries] (s1) {1. Hello!};
\node[draw=poporange, fill=poporangeL, rounded corners, minimum width=2.1cm, minimum height=1cm, align=center, font=\small\bfseries, right=of s1] (s2) {2. How\\much...?};
\node[draw=poppurple, fill=poppurpleL, rounded corners, minimum width=2.1cm, minimum height=1cm, align=center, font=\small\bfseries, right=of s2] (s3) {3. Bargain};
\node[draw=popgreen, fill=popgreenL, rounded corners, minimum width=2.1cm, minimum height=1cm, align=center, font=\small\bfseries, right=of s3] (s4) {4. Pay};
\node[draw=poppink, fill=poppinkL, rounded corners, minimum width=2.1cm, minimum height=1cm, align=center, font=\small\bfseries, right=of s4] (s5) {5. Goodbye!};
\draw[-{Stealth}, thick, popmuted] (s1) -- (s2);
\draw[-{Stealth}, thick, popmuted] (s2) -- (s3);
\draw[-{Stealth}, thick, popmuted] (s3) -- (s4);
\draw[-{Stealth}, thick, popmuted] (s4) -- (s5);
\node[below=0.15cm of s1, font=\scriptsize, align=center, text=popdark] {\eng{Good morning!}\\ \eng{Can I help you?}};
\node[below=0.15cm of s2, font=\scriptsize, align=center, text=popdark] {\eng{How much}\\ \eng{is / are...?}};
\node[below=0.15cm of s3, font=\scriptsize, align=center, text=popdark] {\eng{Your last}\\ \eng{price?}};
\node[below=0.15cm of s4, font=\scriptsize, align=center, text=popdark] {\eng{Here you are.}\\ \eng{Your change.}};
\node[below=0.15cm of s5, font=\scriptsize, align=center, text=popdark] {\eng{Thank you!}\\ \eng{Come again!}};
\end{tikzpicture}
\legende{Le squelette du dialogue commercial : cinq bulles à compléter avec tes propres formules.}
\end{popfigure}

\begin{methode}[Comment apprendre un dialogue en cinq étapes]
\begin{enumerate}
\item \textbf{Écoute} le professeur lire le dialogue, livre fermé : concentre-toi sur la mélodie des phrases.
\item \textbf{Répète en chœur} avec toute la classe, phrase par phrase, en imitant l'intonation.
\item \textbf{Lis à deux} : un élève joue le client, l'autre le vendeur, puis échangez les rôles.
\item \textbf{Remplace} les mots soulignés (articles, prix, quantités) par d'autres mots : c'est la \emph{substitution}.
\item \textbf{Joue sans regarder le texte} : tu parles maintenant vraiment anglais !
\end{enumerate}
\end{methode}

\begin{exoresolu}[Retrouver les cinq étapes]
\textbf{Énoncé.} Voici des répliques mélangées d'un dialogue au marché. Remets-les dans l'ordre des cinq étapes (greeting, information, bargaining, paying, leave-taking) :\\
a) \eng{Here is your change: 300 FCFA.} — b) \eng{Good morning, Sir!} — c) \eng{What is your last price?} — d) \eng{How much is a cup of rice?} — e) \eng{Thank you, come again!}
\tcblower
\textbf{Solution détaillée.} On suit le squelette du dialogue :\\
1. \textbf{Greeting :} b) \eng{Good morning, Sir!} — on commence toujours par saluer.\\
2. \textbf{Information :} d) \eng{How much is a cup of rice?} — on demande le prix.\\
3. \textbf{Bargaining :} c) \eng{What is your last price?} — on négocie.\\
4. \textbf{Paying :} a) \eng{Here is your change: 300 FCFA.} — le vendeur rend la monnaie.\\
5. \textbf{Leave-taking :} e) \eng{Thank you, come again!} — on prend congé poliment.
\end{exoresolu}

\courssub{Practice: read aloud, then substitute}

\begin{experience}[Substitution drill — change the dialogue!]
Travaille en binôme. Reprends le \textbf{dialogue 1} et remplace les éléments soulignés par ceux de la liste ci-dessous. Chaque binôme joue sa nouvelle version devant la classe.
\begin{itemize}
\item \textbf{Articles :} tomatoes and onions $\rightarrow$ carrots and cabbage / oranges and bananas / beans and groundnuts.
\item \textbf{Prix :} 500 FCFA a kilo $\rightarrow$ 300 FCFA / 750 FCFA / 1,000 FCFA a kilo.
\item \textbf{Quantités :} two kilos $\rightarrow$ three kilos / half a kilo / five kilos.
\item \textbf{Monnaie :} 1,500 FCFA $\rightarrow$ 2,000 FCFA / 5,000 FCFA.
\end{itemize}
Exemple de début transformé : \eng{Good morning, Madam! — Good morning! What can I do for you? — I'd like to buy some carrots and cabbage. How much are the carrots?}
\end{experience}

\begin{attention}[Erreurs fréquentes des francophones dans ces dialogues]
\begin{itemize}
\item Ne dis pas \eng{How much are the tomatoes cost?} — avec \eng{how much}, on utilise soit \eng{be} (\eng{How much are the tomatoes?}), soit \eng{cost} (\eng{How much do the tomatoes cost?}), jamais les deux ensemble.
\item \eng{Change} (monnaie) est singulier : \eng{Here \textbf{is} your change}, pas \eng{here are}.
\item \eng{Each} signifie « chacun, pièce » : \eng{The batteries are 250 FCFA each} = « les piles coûtent 250 FCFA pièce ».
\item Pour commander, dis \eng{I'd like two kilos, please} ou \eng{Give me two kilos, please} ; le mot \eng{please} est indispensable pour rester poli.
\item \eng{Sell} est transitif indirect pour le prix : \eng{How much do you sell it \textbf{for}?} (pas \eng{How much do you sell it?}).
\end{itemize}
\end{attention}

\begin{savaistu}[Fixed prices or bargaining?]
Au Royaume-Uni et aux États-Unis, on ne négocie presque \textbf{jamais} dans les magasins : les prix sont fixes (\eng{fixed prices}), comme dans le dialogue 2. En revanche, on négocie volontiers dans les marchés (\eng{markets}), les brocantes (\eng{car boot sales}) et lors de l'achat d'une voiture d'occasion (\eng{a second-hand car}) ! Au Cameroun, c'est l'inverse : la négociation fait partie de la vie quotidienne au marché, et dire \eng{What is your last price?} est presque un signe de politesse.
\end{savaistu}

\begin{exercice}[Your turn to speak!]
En binôme, écris puis joue un mini-dialogue complet (au moins huit répliques) qui respecte les cinq étapes du squelette. Situation au choix : acheter du poisson fumé au marché de Buea, ou acheter du savon et des allumettes dans une boutique de Bamenda. Utilise au moins deux formules de négociation et une formule de conclusion.
\end{exercice}

\begin{corrige}[Exercice — Your turn to speak!]
Exemple de production attendue :\\
— \eng{Good afternoon, Madam!} — \eng{Good afternoon! What can I do for you?} — \eng{How much do you sell this soap for?} — \eng{It is 400 FCFA a bar.} — \eng{Oh, that's a bit expensive. Can you reduce the price?} — \eng{For you, 350 FCFA.} — \eng{If I buy three bars, how much will I pay?} — \eng{Three for 1,000 FCFA.} — \eng{OK, I'll take three. Here is 1,000 FCFA.} — \eng{Thank you very much. Have a nice day!} — \eng{Goodbye!}\\
Barème indicatif : salutation (1 pt), question de prix correcte (1 pt), négociation (2 pts), paiement (1 pt), congé (1 pt), prononciation et fluidité (4 pts).
\end{corrige}

\coursec{Pronunciation, stress and intonation}

Vous venez de travailler les dialogues modèles : vous savez maintenant quoi dire au marché et à la boutique. Mais entre lire un dialogue sur le papier et le dire à voix haute, il y a un pas décisif : la \cle{prononciation}. Un prix mal prononcé peut coûter cher — imaginez que vous confondiez \eng{thirteen} (13) et \eng{thirty} (30) en pleine négociation ! Cette section vous apprend à prononcer les prix, à accentuer les mots importants et à utiliser la bonne \cle{intonation}, pour que votre anglais soit compris immédiatement par n'importe quel vendeur ou client anglophone.

\courssub{Saying prices and numbers aloud}

Commençons par observer. Écoutez votre professeur lire ces prix à voix haute, puis répétez en chœur :

\begin{exemplebox}[Observe and repeat: prices at the market]
\begin{itemize}
\item \eng{One hundred FCFA} — prononcez « ouane \textbf{han}-drèd èf-si-èf-é ».
\item \eng{Two fifty} (= 250 FCFA) — prononcez « tou \textbf{fif}-ti ».
\item \eng{Five hundred FCFA} — prononcez « faïve \textbf{han}-drèd ».
\item \eng{One thousand three hundred FCFA} — prononcez « ouane \textbf{zaou}-zande tri \textbf{han}-drèd ».
\item \eng{Two thousand five hundred FCFA} — prononcez « tou \textbf{zaou}-zande faïve \textbf{han}-drèd ».
\end{itemize}
\end{exemplebox}

\begin{propriete}[Règle : comment dire les prix en anglais]
\begin{enumerate}
\item On dit d'abord le \cle{nombre}, puis éventuellement la monnaie : \eng{five hundred FCFA} (« cinq cents francs »).
\item \cle{hundred} et \cle{thousand} restent \textbf{invariables} : on dit \eng{two hundred}, \eng{five thousand}, jamais \eng{two hundreds}.
\item Pour les prix entre 100 et 999, les anglophones utilisent souvent la forme courte : \eng{two fifty} = 250, \eng{three twenty-five} = 325. La monnaie est sous-entendue.
\item Le mot \eng{FCFA} s'épelle lettre par lettre : « èf-si-èf-é ».
\end{enumerate}
\end{propriete}

\begin{methode}[Bien prononcer un prix en trois étapes]
\begin{enumerate}
\item Dites d'abord le nombre \textbf{lentement}, syllabe par syllabe : \eng{one — thou — sand}.
\item Ajoutez la monnaie : \eng{one thousand FCFA} (« ouane zaou-zande èf-si-èf-é »).
\item Accélérez progressivement jusqu'au débit naturel, en gardant l'accent sur le mot fort : \eng{one THOUsand FCFA}.
\end{enumerate}
Répétez chaque prix trois fois : lent, moyen, rapide. C'est la méthode des sportifs : la lenteur construit la précision, la vitesse vient ensuite.
\end{methode}

\courssub{The -teen / -ty trap: thirteen or thirty?}

Voici l'une des distinctions les plus importantes de toute la prononciation anglaise, surtout au marché. Observez :

\begin{center}
\begin{tabularx}{\linewidth}{|X|X|X|}
\hline
\rowcolor{popblueL} Nombre & Prononciation & Accent \\
\hline
\eng{thirteen} (13) & « seur-\textbf{tine} » & accent sur la \textbf{2\textsuperscript{e}} syllabe \\
\hline
\eng{thirty} (30) & « \textbf{seur}-ti » & accent sur la \textbf{1\textsuperscript{re}} syllabe \\
\hline
\eng{fifteen} (15) & « fif-\textbf{tine} » & accent sur la \textbf{2\textsuperscript{e}} syllabe \\
\hline
\eng{fifty} (50) & « \textbf{fif}-ti » & accent sur la \textbf{1\textsuperscript{re}} syllabe \\
\hline
\eng{fourteen} (14) & « for-\textbf{tine} » & accent sur la \textbf{2\textsuperscript{e}} syllabe \\
\hline
\eng{forty} (40) & « \textbf{for}-ti » & accent sur la \textbf{1\textsuperscript{re}} syllabe \\
\hline
\end{tabularx}
\end{center}

\begin{propriete}[Règle : -teen contre -ty]
\begin{itemize}
\item Les nombres en \cle{-teen} (13 à 19) portent l'accent sur la \textbf{dernière syllabe} et le « n » final se prononce clairement : thir\textbf{TEEN}, fif\textbf{TEEN}.
\item Les nombres en \cle{-ty} (20, 30, 40…) portent l'accent sur la \textbf{première syllabe} : \textbf{THIR}ty, \textbf{FIF}ty.
\item Le « t » de \eng{-ty} est souvent doux, presque un « d » : \eng{twenty} se dit « touén-di ».
\end{itemize}
\end{propriete}

\begin{attention}[Piège : 1 500 ou 50 000 FCFA ?]
Au marché de Mokolo ou de Buea, confondre \eng{fifteen hundred} et \eng{fifty thousand} peut vous coûter très cher ! Pour éviter tout malentendu :
\begin{itemize}
\item insistez bien sur la syllabe accentuée : fif\textbf{TEEN} hundred (1 500) / \textbf{FIF}ty thousand (50 000) ;
\item en cas de doute, demandez au vendeur de répéter ou d'écrire le chiffre : \eng{Could you write it down, please?} (« Pourriez-vous l'écrire, s'il vous plaît ? »).
\end{itemize}
\end{attention}

\begin{exoresolu}[Exercice résolu : quel prix entendez-vous ?]
Énoncé. Votre professeur prononce chaque prix deux fois. Entourez le prix que vous entendez :
\begin{enumerate}
\item \eng{thirty FCFA} ou \eng{thirteen FCFA} ?
\item \eng{fifty FCFA} ou \eng{fifteen FCFA} ?
\item \eng{fourteen hundred} ou \eng{four thousand} ?
\end{enumerate}
\tcblower
Solution détaillée.
\begin{enumerate}
\item Si vous entendez « \textbf{seur}-ti » avec l'accent au début, c'est \eng{thirty FCFA} (30 F). Si vous entendez « seur-\textbf{tine} » avec l'accent à la fin et un « n » net, c'est \eng{thirteen FCFA} (13 F).
\item « \textbf{fif}-ti » = \eng{fifty} (50 F) ; « fif-\textbf{tine} » = \eng{fifteen} (15 F).
\item « for-\textbf{tine} han-drèd » = \eng{fourteen hundred} (1 400 F) ; « \textbf{fo} zaou-zande » = \eng{four thousand} (4 000 F). Remarquez que l'on ne dit jamais \eng{forty hundred} en anglais correct : au-delà de 1 900, on utilise \eng{thousand}.
\end{enumerate}
Astuce de vérification : fermez les yeux et repérez uniquement la syllabe la plus forte. Si elle est à la fin, c'est -teen ; si elle est au début, c'est -ty.
\end{exoresolu}

\courssub{Word stress: the music of market words}

En anglais, chaque mot de plusieurs syllabes possède une \cle{syllabe accentuée}, plus forte et plus longue que les autres. Si vous accentuez la mauvaise syllabe, votre interlocuteur peut ne pas comprendre. Voici les mots essentiels du marché :

\begin{center}
\begin{tabularx}{\linewidth}{|X|X|X|}
\hline
\rowcolor{popblueL} English & Prononciation & Français \\
\hline
\eng{toMAtoes} & « to-\textbf{ma}-toz » & tomates \\
\hline
\eng{ONions} & « \textbf{a}-nionz » & oignons \\
\hline
\eng{CABbage} & « \textbf{ka}-bidj » & chou \\
\hline
\eng{BARgain} & « \textbf{bar}-guine » & marchander / bonne affaire \\
\hline
\eng{DIScount} & « \textbf{dis}-kaounte » & remise, rabais \\
\hline
\eng{change} & « \textbf{tchéin}-dj » & monnaie (à rendre) \\
\hline
\eng{reCEIPT} & « ri-\textbf{site} » & reçu \\
\hline
\eng{CUStomer} & « \textbf{keu}-sto-meur » & client \\
\hline
\end{tabularx}
\end{center}

\begin{attention}[Erreurs fréquentes des francophones]
\begin{itemize}
\item Ne prononcez \textbf{jamais} le « h » de \eng{how} à la française : dites « \textbf{haou} » (comme un souffle suivi de « aou »), pas « o ».
\item \eng{much} ne se dit pas « moutch » mais « \textbf{meutch} » : \eng{How much is it?} = « haou meutch iz ite ».
\item \eng{the} se prononce « ze » (langue entre les dents), pas « dé » : \eng{the market} = « ze markète ».
\item Dans \eng{change}, le « ch » anglais se prononce « tch » : « tchéindj », jamais « chandj » à la française.
\item Attention à l'orthographe de \eng{receipt} : le « p » ne se prononce pas (« ri-site »), mais il s'écrit !
\end{itemize}
\end{attention}

\courssub{Intonation: going up and going down}

L'\cle{intonation} est la mélodie de la phrase : la voix monte ou descend à la fin. En français, notre mélodie est relativement plate ; en anglais, elle est très marquée, et elle \textbf{change le sens} de la phrase. Observez ces deux questions du dialogue précédent :

\begin{itemize}
\item \eng{Do you have any RICE? ↗} — la voix \textbf{monte} : le vendeur attend \eng{yes} ou \eng{no}.
\item \eng{HOW much IS it? ↘} — la voix \textbf{descend} : le client attend une information précise (un prix).
\end{itemize}

\begin{popfigure}
\begin{tikzpicture}[font=\small]
% Flèche montante
\node[align=center] at (3,0) {\eng{Do you have any RICE?}};
\draw[-{Stealth[length=3mm]}, line width=1.5pt, popgreen] (0.8,0.5) .. controls (2.2,0.6) and (3.6,0.9) .. (5.2,1.8);
\node[popgreen, align=center] at (3,2.2) {question yes/no : la voix \textbf{monte} $\nearrow$};
% Flèche descendante
\node[align=center] at (11,0) {\eng{HOW much IS it?}};
\draw[-{Stealth[length=3mm]}, line width=1.5pt, poporange] (8.8,1.8) .. controls (10,1.4) and (11.6,0.9) .. (13.2,0.5);
\node[poporange, align=center] at (11,2.2) {question wh- : la voix \textbf{descend} $\searrow$};
\end{tikzpicture}
\legende{Les deux mélodies de base de l'anglais : intonation montante pour les questions fermées, descendante pour les questions en wh- et les affirmations.}
\end{popfigure}

\begin{propriete}[Règle d'intonation]
\begin{enumerate}
\item Les \cle{questions fermées} (réponse \eng{yes} ou \eng{no}, commençant par \eng{do}, \eng{is}, \eng{can}, \eng{have}…) se terminent par une intonation \textbf{montante} ↗ : \eng{Is it FRESH? ↗} (« Est-ce frais ? » — on attend yes/no).
\item Les \cle{questions ouvertes} (commençant par un mot en \eng{wh-} : \eng{what}, \eng{how much}, \eng{where}, \eng{when}) se terminent par une intonation \textbf{descendante} ↘ : \eng{HOW much ARE they? ↘} (« Combien coûtent-ils ? » — on attend un prix).
\item Les \cle{affirmations} et les réponses descendent également ↘ : \eng{It's FIVE hundred FCFA. ↘} (« Cela coûte cinq cents francs. »).
\item \textbf{Exception} : pour exprimer la surprise ou demander de répéter, une question en wh- peut monter : \eng{How MUCH? ↗} (« Combien, dites-vous ?! »).
\end{enumerate}
\end{propriete}

\begin{exemplebox}[Exemples traduits : monte ou descend ?]
\begin{itemize}
\item \eng{Is it FRESH? ↗} « Est-ce frais ? » (attend : yes/no)
\item \eng{Do you SELL plantains? ↗} « Vendez-vous des plantains ? » (attend : yes/no)
\item \eng{HOW much ARE they? ↘} « Combien coûtent-ils ? » (attend : un prix)
\item \eng{WHERE is the market? ↘} « Où est le marché ? » (attend : un lieu)
\item \eng{That's TWO fifty, please. ↘} « Cela fera deux cent cinquante, s'il vous plaît. » (affirmation)
\end{itemize}
\end{exemplebox}

\courssub{Politeness through tone of voice}

En anglais, la politesse ne tient pas seulement aux mots (\eng{please}, \eng{thank you}) : elle tient au \cle{ton}. Un ton chaleureux, qui monte légèrement, adoucit la demande ; un ton plat ou descendant peut sembler sec, voire impoli.

\begin{exemplebox}[The polite trio]
\begin{itemize}
\item \eng{Excuse ME ↗…} « Excusez-moi… » — la voix monte pour attirer l'attention sans agresser.
\item \eng{…please ↗} « …s'il vous plaît » — dit avec un ton montant, la demande devient aimable.
\item \eng{Thank YOU ↘} « Merci » — la voix descend, ferme et chaleureuse, pour conclure.
\end{itemize}
Comparez : \eng{Thank YOU ↘} (sincère) et \eng{thank you ↗} (distrait, automatique). Au marché, préférez toujours le ton descendant et appuyé sur \eng{YOU}.
\end{exemplebox}

\begin{savaistu}[Les prix sans le mot de la monnaie]
En Grande-Bretagne, les vendeurs disent très souvent le prix \textbf{sans nommer la monnaie} : \eng{That'll be two fifty} signifie « Cela fera 2,50 £ ». Le contexte suffit ! De même, au Cameroun anglophone, un vendeur de Bamenda peut dire simplement \eng{Five hundred} pour 500 FCFA. Et le mot \eng{change} (« monnaie à rendre ») vient du vieux français \emph{change} — les langues se sont prêté des mots pendant des siècles !
\end{savaistu}

\courssub{Practice time: choral drills and mimed mini-dialogues}

\begin{experience}[Choral drill: the price ladder]
Répétition en chœur, menée par le professeur. Montez l'échelle des prix, d'abord lentement, puis de plus en plus vite :
\begin{enumerate}
\item \eng{One hundred FCFA.} (« ouane han-drèd »)
\item \eng{Two fifty.} (« tou fif-ti »)
\item \eng{Five hundred FCFA.} (« faïve han-drèd »)
\item \eng{One thousand FCFA.} (« ouane zaou-zande »)
\item \eng{Two thousand five hundred FCFA.} (« tou zaou-zande faïve han-drèd »)
\end{enumerate}
Variante : le professeur montre un prix écrit au tableau ; la classe le dit à voix haute, puis un élève seul le répète. Attention à l'accent : \eng{one THOUsand}, pas « one thou-SAND ».
\end{experience}

\begin{experience}[Mimed mini-dialogue: say it with your hands!]
Par binômes, jouez ce mini-dialogue en mimant l'intonation avec la main : la main monte ↗ pour les questions yes/no, descend ↘ pour les questions en wh- et les réponses.
\begin{itemize}
\item A : \eng{Excuse ME ↗, do you have any ONions? ↗}
\item B : \eng{Yes, we DO. ↘ Fresh ONions! ↘}
\item A : \eng{HOW much ARE they? ↘}
\item B : \eng{Three hundred FCFA a KIlo. ↘}
\item A : \eng{Can you give me a DIScount? ↗}
\item B : \eng{Two FIFty, my friend. ↘}
\item A : \eng{Thank YOU! ↘}
\end{itemize}
Puis inversez les rôles et remplacez \eng{onions} par \eng{tomatoes}, \eng{cabbage} ou \eng{rice}, avec un nouveau prix.
\end{experience}

\begin{exercice}[Pour vous entraîner]
\begin{enumerate}
\item Écrivez en toutes lettres, puis prononcez à voix haute : 150 FCFA ; 1 300 FCFA ; 2 750 FCFA.
\item Indiquez par une flèche (↗ ou ↘) l'intonation correcte : \eng{Is the rice fresh?} / \eng{How much is a bag of rice?} / \eng{It costs five thousand FCFA.} / \eng{Can I pay now?}
\item Soulignez la syllabe accentuée : \eng{tomatoes}, \eng{discount}, \eng{cabbage}, \eng{bargain}, \eng{thirteen}, \eng{thirty}.
\end{enumerate}
\end{exercice}

\begin{corrige}[Exercice « Pour vous entraîner »]
\begin{enumerate}
\item \eng{One hundred and fifty FCFA} (« ouane han-drèd and fif-ti ») ; \eng{One thousand three hundred FCFA} (« ouane zaou-zande tri han-drèd ») ; \eng{Two thousand seven hundred and fifty FCFA} (« tou zaou-zande sé-ven han-drèd and fif-ti »).
\item \eng{Is the rice fresh? ↗} (question yes/no) ; \eng{How much is a bag of rice? ↘} (question wh-) ; \eng{It costs five thousand FCFA. ↘} (affirmation) ; \eng{Can I pay now? ↗} (question yes/no).
\item \eng{to\textbf{MA}toes}, \eng{\textbf{DIS}count}, \eng{\textbf{CAB}bage}, \eng{\textbf{BAR}gain}, \eng{thir\textbf{TEEN}}, \eng{\textbf{THIR}ty}.
\end{enumerate}
\end{corrige}

\begin{aretenir}[L'essentiel de la section]
\begin{itemize}
\item Les prix : nombre + monnaie ; \eng{hundred} et \eng{thousand} invariables ; forme courte \eng{two fifty} = 250.
\item \textbf{-TEEN} accentué à la fin (13, 15) contre \textbf{-TY} accentué au début (30, 50) : la distinction vitale du marché.
\item Questions yes/no : intonation \textbf{montante} ↗ ; questions en wh- et affirmations : intonation \textbf{descendante} ↘.
\item La politesse passe par le ton : \eng{Excuse me ↗, please ↗, thank YOU ↘}.
\item Prononciations à soigner : \eng{how} « haou », \eng{much} « meutch », \eng{change} « tchéindj ».
\end{itemize}
\end{aretenir}

\coursec{Guided role-plays: let's speak!}

Après avoir travaillé la prononciation, l'accentuation et l'intonation des formules clés, nous passons maintenant à la mise en situation réelle. C'est le moment de transformer la classe en un espace de communication authentique où l'anglais devient l'outil indispensable pour \eng{get things done} (« accomplir les choses »). Les trois jeux de rôle ci-dessous sont conçus selon une progression pédagogique stricte : du très guidé (sécurité maximale) au libre (autonomie totale), conformément aux exigences du programme MINESEC pour la compétence \eng{Speaking}.

\courssub{Role-play 1: Highly guided practice — Swapping items and prices}

\begin{textebox}{Role-play 1 Script (Base Dialogue)}
\textbf{Customer:} “Good morning. How much are the \textbf{beans}, please?”
\textbf{Seller:} “Good morning. They are \textbf{800 FCFA a kilo}.”
\textbf{Customer:} “That's a bit expensive. Can you give me a discount?”
\textbf{Seller:} “Alright. For you, 750 FCFA a kilo.”
\textbf{Customer:} “OK. I'll take two kilos, please.”
\textbf{Seller:} “Here you are. That makes 1,500 FCFA.”
\textbf{Customer:} “Here is 2,000 FCFA.”
\textbf{Seller:} “Thank you. Here is your change: 500 FCFA. Have a nice day!”
\textbf{Customer:} “Thank you. Goodbye!”
\end{textebox}

\begin{methode}{Comment réussir le Jeu de rôle 1 (Très guidé)}
\begin{enumerate}
    \item \textbf{Formez des paires :} Un élève A (Customer), un élève B (Seller).
    \item \textbf{Distribuez les cartes « Articles/Prix » :} Chaque paire reçoit 3 à 4 cartes (ex. \eng{eggs — 100 FCFA each} ; \eng{tomatoes — 600 FCFA a pile} ; \eng{bread — 300 FCFA a loaf}).
    \item \textbf{Répétez le dialogue de base} en remplaçant \textbf{uniquement} les mots en gras (article, prix unitaire, quantité, total, monnaie).
    \item \textbf{Changez de rôle} après chaque carte pour que chacun pratique les deux parties.
    \item \textbf{Objectif :} Fluidité et prononciation correcte des nombres et des formules de politesse.
\end{enumerate}
\end{methode}

\begin{exemplebox}{Cartes d'entraînement (Item Cards)}
\begin{center}
\begin{tabular}{|l|l|l|}
\hline
\rowcolor{popblueL} \textbf{Article} & \textbf{Unit Price} & \textbf{Quantity bought} \\
\hline
Beans & 800 FCFA a kilo & 2 kilos \\
Eggs & 100 FCFA each & 10 eggs \\
Tomatoes & 600 FCFA a pile & 3 piles \\
Bread & 300 FCFA a loaf & 2 loaves \\
Milk (powder) & 500 FCFA a sachet & 4 sachets \\
\hline
\end{tabular}
\end{center}
\emph{Note :} Pour \eng{eggs} et \eng{bread}, on utilise \eng{each} (pièce) et \eng{a loaf} (pain). Pour \eng{tomatoes} au marché camerounais, \eng{a pile} (un tas) est l'unité courante.
\end{exemplebox}

\begin{exoresolu}{Calcul du total et de la monnaie (Role-play 1)}
\textbf{Énoncé :} Le client achète 3 kilos de \eng{beans} à 750 FCFA le kilo (prix négocié). Il donne 3 000 FCFA. Écrivez la réplique du vendeur pour annoncer le total et rendre la monnaie.
\\
\textbf{Solution détaillée :}
\begin{enumerate}
    \item Calcul du total : 3 × 750 = 2 250 FCFA.
    \item Calcul de la monnaie : 3 000 - 2 250 = 750 FCFA.
    \item Réplique type : \eng{“That makes 2,250 FCFA. Here is your change: 750 FCFA.”}
\end{enumerate}
\emph{Astuce :} Entraînez-vous à dire les nombres en anglais : « two thousand two hundred and fifty » / « seven hundred and fifty ».
\end{exoresolu}

\courssub{Role-play 2: Semi-guided practice — Shopping list vs Price list}

Ici, le dialogue n'est plus écrit mot pour mot. Vous devez \textbf{construire} l'échange en utilisant les formules fonctionnelles apprises (Sections 2 et 3). C'est l'étape cruciale pour passer de la reproduction à la production langagière.

\begin{experience}{Fiches de rôle (Role Cards)}
\begin{center}
\begin{tabularx}{\linewidth}{|X|X|}
\hline
\rowcolor{popgreenL} \textbf{CUSTOMER CARD} & \textbf{SELLER CARD} \\
\hline
\textbf{You need to buy:} & \textbf{Your prices:} \\
\begin{itemize}
    \item 1 bag of rice
    \item 2 litres of oil
    \item 1 packet of sugar
\end{itemize}
\textbf{Your budget:} 10 000 FCFA. &
\begin{itemize}
    \item Rice: 6 000 FCFA a bag
    \item Oil: 1 500 FCFA a litre
    \item Sugar: 900 FCFA a packet
\end{itemize}
\textbf{Strategy:} You can give a small discount (max 5\%) if the customer buys everything. \\
\hline
\end{tabularx}
\end{center}
\end{experience}

\begin{methode}{Déroulement du Jeu de rôle 2 (Semi-guidé)}
\begin{enumerate}
    \item \textbf{Préparation silencieuse (2 min) :} Le client calcule le coût total estimé (6 000 + 3 000 + 900 = 9 900 FCFA). Le vendeur prépare ses arguments de vente.
    \item \textbf{Interaction orale (3-4 min) :} Pas de texte écrit. Utilisez les formules :
    \begin{itemize}
        \item \eng{“Good morning. I'd like to buy... / How much is...?”}
        \item \eng{“That's too much. Can you lower the price?”}
        \item \eng{“I can offer you... / Let's say...”}
        \item \eng{“Deal. / No deal, thank you.”}
    \end{itemize}
    \item \textbf{Conclusion obligatoire :} Le client doit annoncer s'il achète ou non, et le vendeur doit calculer le total exact et la monnaie si achat il y a.
    \item \textbf{Échange des rôles} et recommencement avec une nouvelle liste (ex. \eng{1 kg meat, 5 kg plantains, 1 bottle soap}).
\end{enumerate}
\end{methode}

\begin{exemplebox}{Exemple d'échange semi-libre (Transcription annotée)}
\textbf{Customer:} \eng{“Good morning. How much is a bag of rice?”} (Formule standard)
\textbf{Seller:} \eng{“Good morning. It's 6,000 FCFA.”} (Réponse directe)
\textbf{Customer:} \eng{“And the oil?”} (Ellipse courante : \eng{“How much is the oil?”})
\textbf{Seller:} \eng{“1,500 FCFA a litre.”}
\textbf{Customer:} \eng{“I need two litres. That's 3,000. And sugar?”} (Calcul mental verbalisé)
\textbf{Seller:} \eng{“900 FCFA a packet.”}
\textbf{Customer:} \eng{“Total is 9,900. I have 10,000. Can you make it 9,500 for all?”} (Négociation globale)
\textbf{Seller:} \eng{“Hmm... OK. 9,500 for everything. Good price!”} (Acceptation + commentaire)
\textbf{Customer:} \eng{“Deal. Here is 10,000.”} (Conclusion)
\textbf{Seller:} \eng{“Thank you. Here is 500 FCFA change. Have a good day!”} (Monnaie + formule de clôture)
\end{exemplebox}

\courssub{Role-play 3: Free practice — The Classroom Market Simulation}

C'est l'activité sommative de la leçon. La classe devient un véritable \eng{market} (marché). L'objectif est l'interaction spontanée, la mobilisation du lexique (quantités, contenants, articles) et la gestion de l'imprévu (rupture de stock, budget insuffisant, négociation acharnée).

\begin{popfigure}
\centering
\begin{tikzpicture}[
    stall/.style={draw=popdark, thick, fill=popcream, rounded corners=8pt, minimum width=3.5cm, minimum height=2.5cm, align=center, font=\small},
    arrow/.style={->, >=Stealth, thick, popblue},
    customer/.style={circle, draw=poporange, fill=poporangeL, minimum size=8pt, inner sep=1pt},
    label/.style={font=\tiny, text=popmuted, anchor=north}
]
\node[stall, align=center] (S1) at (-4, 3){STALL 1\\[2pt] \textbf{Mama Ngo}\\[2pt] Rice, Beans, Oil\\[2pt] \emph{Fixed prices}};
\node[stall, align=center] (S2) at (4, 3){STALL 2\\[2pt] \textbf{Papa Jean}\\[2pt] Meat, Fish, Eggs\\[2pt] \emph{Negotiable}};
\node[stall, align=center] (S3) at (-4, -3){STALL 3\\[2pt] \textbf{Tante Marie}\\[2pt] Tomatoes, Onions,\\[2pt] Peppers, Plantains\\[2pt] \emph{Per pile/heap}};
\node[stall, align=center] (S4) at (4, -3){STALL 4\\[2pt] \textbf{Frère Paul}\\[2pt] Bread, Milk, Sugar,\\[2pt] Soap, Maggi\\[2pt] \emph{Promo today!}};

\node[customer] (C1) at (-1.5, 0) {};
\node[customer] (C2) at (1.5, 0) {};
\node[customer] (C3) at (0, 1.5) {};
\node[customer] (C4) at (0, -1.5) {};

\draw[arrow, dashed] (C1) -- (S1.west);
\draw[arrow, dashed] (C1) -- (S3.west);
\draw[arrow, dashed] (C2) -- (S2.east);
\draw[arrow, dashed] (C2) -- (S4.east);
\draw[arrow, dashed] (C3) -- (S1.north);
\draw[arrow, dashed] (C3) -- (S2.north);
\draw[arrow, dashed] (C4) -- (S3.south);
\draw[arrow, dashed] (C4) -- (S4.south);

\node[label] at (-4, 4.5) {Fixed Price Zone};
\node[label] at (4, 4.5) {Bargaining Zone};
\node[label] at (-4, -4.5) {Fresh Produce Zone};
\node[label] at (4, -4.5) {General Store Zone};
\node[label, align=center] at (0, 0){START:\\[-2pt] 5 000 FCFA Budget};
\end{tikzpicture}
\legende{Disposition de la classe pour le « Classroom Market » : 4 étals spécialisés, clients au centre avec budget de 5 000 FCFA, déplacements libres (flèches pointillées).}
\end{popfigure}

\begin{definition}{Règles du « Classroom Market » (Role-play 3)}
\begin{enumerate}
    \item \textbf{Rôles :} 8 à 10 \eng{Sellers} (derrière leur table/étal avec affichettes de prix), le reste = \eng{Customers}.
    \item \textbf{Budget Client :} 5 000 FCFA (billets/monnaie factices fournis par le professeur).
    \item \textbf{Objectif Client :} Acheter \textbf{exactement 3 articles différents} chez \textbf{au moins 2 vendeurs différents} en négociant.
    \item \textbf{Objectif Vendeur :} Vendre son stock, gagner de l'argent, pratiquer l'anglais commercial.
    \item \textbf{Contrainte linguistique :} \textbf{Zéro français toléré.} Si bloqué $\rightarrow$ \eng{“Sorry, can you repeat? / Speak slowly, please.”}
    \item \textbf{Minimum fonctionnel :} Utiliser au moins \textbf{4 formules distinctes} de la leçon (salutation, prix, négociation, paiement, remerciement).
    \item \textbf{Durée :} 15 minutes intenses. Rotation possible (Clients $\leftrightarrow$ Vendeurs) à mi-temps.
\end{enumerate}
\end{definition}

\begin{aretenir}{Check-liste « Minimum Syndical » pour le Jeu de rôle 3}
Avant de quitter un étal, le client doit avoir dit :
\begin{itemize}
    \item[\checkmark] \eng{“Good morning / Hello.”} (Salutation)
    \item[\checkmark] \eng{“How much is...? / What's the price of...?”} (Question prix)
    \item[\checkmark] \eng{“Can you reduce...? / Any discount?”} (Négociation)
    \item[\checkmark] \eng{“I'll take it. / Here is the money.”} (Paiement)
    \item[\checkmark] \eng{“Thank you. Goodbye / Have a nice day.”} (Clôture)
\end{itemize}
Le vendeur doit avoir répondu à chaque étape.
\end{aretenir}

\courssub{Peer Observation Grid — Grille d'observation par les pairs}

Pendant que deux élèves jouent, deux autres observent avec cette grille. Cela maintient toute la classe active et développe l'écoute critique (\eng{critical listening}).

\begin{center}
\begin{tabularx}{\linewidth}{|>{\hsize=0.8\hsize}X|>{\hsize=1.2\hsize}X|>{\hsize=1.2\hsize}X|>{\hsize=0.8\hsize}X|}
\hline
\rowcolor{poppurpleL} \textbf{Critère d'observation} & \textbf{Formules attendues (Exemples)} & \textbf{Observé ? (✔/✘)} & \textbf{Commentaire / Correction} \\
\hline
\textbf{1. Salutation} & \eng{“Good morning. / Hello. / Hi.”} & & \\
\hline
\textbf{2. Question sur le prix / article} & \eng{“How much is...? / How much for...? / I'd like...”} & & \\
\hline
\textbf{3. Négociation / Relance} & \eng{“Too expensive. / Can you lower...? / Last price?”} & & \\
\hline
\textbf{4. Paiement \& Monnaie} & \eng{“Here is... FCFA. / That makes... / Your change.”} & & \\
\hline
\textbf{5. Remerciements / Clôture} & \eng{“Thank you. / Thanks a lot. / Goodbye. / See you.”} & & \\
\hline
\textbf{6. Prononciation claire} & Nombres (\eng{hundred, thousand}), \eng{FCFA}, \eng{discount}, \eng{receipt}. & & \\
\hline
\end{tabularx}
\end{center}

\begin{methode}{Comment utiliser la grille (Pour les observateurs)}
\begin{enumerate}
    \item Cochez ✔ si la formule est utilisée \textbf{correctement} (forme + intonation).
    \item Cochez ✘ si absent ou erreur majeure (ex. \eng{*How much cost?} au lieu de \eng{How much is it?}).
    \item Notez un exemple exact entendu dans la colonne « Commentaire » (bon ou à corriger).
    \item À la fin, donnez la grille aux acteurs pour auto-correction immédiate.
\end{enumerate}
\end{methode}

\courssub{Collective Debriefing — Feedback collectif structuré}

Le débriefing n'est pas une option, c'est \textbf{là que l'apprentissage se consolide}. Le professeur mène la séance au tableau en 3 temps forts.

\begin{experience}{Déroulement du Debriefing (10-15 min)}
\textbf{Temps 1 : Les « Golden Formulas » (Les meilleures formules entendues)}
\begin{itemize}
    \item Le professeur note au tableau 4-5 formules \textbf{parfaites} ou \textbf{créatives} entendues pendant l'activité.
    \item Ex. : \eng{“If I take three, do I get a better price?”} (Négociation conditionnelle spontanée).
    \item Ex. : \eng{“Keep the change.”} (Pourboire / Arrondi — culture pourboire).
    \item La classe répète en chœur (drill) pour ancrer la prosodie.
\end{itemize}

\textbf{Temps 2 : « Error Clinic » (Correction collective bienveillante)}
\begin{itemize}
    \item Le professeur note 3-4 erreurs \textbf{récurrentes} (anonymes) relevées sur les grilles d'observation.
    \item Ex. erreur type : \eng{* “How much it cost?”} $\rightarrow$ Correction : \eng{“How much does it cost?”} ou \eng{“How much is it?”}.
    \item Ex. erreur type : \eng{* “I want 2 kilos of rices.”} $\rightarrow$ Correction : \eng{“Rice” est indénombrable (uncountable)}.
    \item Ex. erreur type : \eng{* “Give me my rest.”} (Traduction littérale de « Donne-moi mon reste ») $\rightarrow$ Correction : \eng{“Here is your change.”}.
    \item La classe propose la correction ; le professeur valide et explique \textbf{pourquoi} (grammaire, lexique, faux ami).
\end{itemize}

\textbf{Temps 3 : « Market Champions » (Bilan chiffré et ludique)}
\begin{itemize}
    \item \textbf{Meilleur négociateur (Client) :} Celui qui a acheté ses 3 articles pour le montant le plus bas (ex. 3 200 FCFA sur 5 000 budget).
    \item \textbf{Meilleur vendeur :} Celui qui a engrangé le plus de chiffre d'affaires (somme des ventes).
    \item \textbf{Meilleur communicant :} Celui qui a utilisé le plus grand nombre de formules différentes (vérifié sur grille pair).
    \item Remise de « prix » symboliques (bonbons, points bonus, applaudissements).
\end{itemize}
\end{experience}

\begin{propriete}[Règle d'or de l'oral : Fluency before accuracy]
\textbf{« Fluency before accuracy »} (\eng{La fluidité avant la justesse}).
\begin{itemize}
    \item \textbf{Pendant le jeu de rôle :} L'objectif est de \textbf{communiquer}, de faire passer le message, de conclure la transaction. Ne vous arrêtez pas pour chercher le mot parfait ou le temps parfait. Utilisez des gestes, reformulez, demandez de l'aide en anglais (\eng{“What do you call...?”}).
    \item \textbf{Pendant le débriefing :} C'est le moment de la \textbf{justesse} (accuracy). On corrige, on explique, on répète les formes correctes.
    \item \textbf{Pourquoi ?} Si vous avez peur de faire une faute, vous ne parlez pas. Si vous ne parlez pas, vous ne progressez pas. L'oral au Bac évalue votre capacité à \textbf{tenir une interaction}, pas à produire des phrases parfaites de grammairien.
\end{itemize}
\end{propriete}

\begin{attention}[Piège majeur : La traduction mot à mot (Calque du français)]
\begin{itemize}
    \item \textbf{Interdit :} Réfléchir en français, traduire mot à mot, puis parler. Cela tue la fluidité et produit de l'anglais « francisé » incompréhensible.
    \item \textbf{Exemples de calques à bannir :}
    \begin{center}
    \begin{tabular}{|l|l|l|}
    \hline
    \rowcolor{poporangeL} \textbf{Français (Pensée)} & \textbf{Calque incorrect (À éviter)} & \textbf{Anglais correct (Formule mémorisée)} \\
    \hline
    Combien ça coûte ? & *How much it costs? & \eng{How much is it? / How much does it cost?} \\
    C'est trop cher. & *It is too much expensive. & \eng{It's too expensive. / That's a bit pricey.} \\
    Donne-moi mon reste. & *Give me my rest. & \eng{Here is your change. / Can I have my change?} \\
    Je veux ça. & *I want that. (Trop direct/impoli) & \eng{I'd like that, please. / I'll take this one.} \\
    C'est bon / D'accord. & *It is good. & \eng{Deal. / Agreed. / That works for me.} \\
    \hline
    \end{tabular}
    \end{center}
    \item \textbf{Solution :} Mémorisez les \textbf{blocs entiers} (chunks) de la leçon comme des « briques Lego » que vous assemblez sans les découper.
\end{itemize}
\end{attention}

\begin{savaistu}[Le Bac Oral et les « Interaction Skills »]
À l'épreuve orale du Baccalauréat (série A), l'examinateur utilise une grille d'évaluation où le critère \textbf{INTERACTION} pèse aussi lourd que \textbf{GRAMMAR} ou \textbf{VOCABULARY}.
\begin{itemize}
    \item \textbf{Interaction} = Savoir \eng{initier} (démarrer), \eng{maintain} (relancer : \eng{And you? / What about...?}), \eng{negotiate meaning} (réparer : \eng{Sorry, I mean... / Let me rephrase}), et \eng{conclude} (clôturer poliment).
    \item Les formules de cette leçon (\eng{How much...?}, \eng{Can you...?}, \eng{Deal}, \eng{Thank you, goodbye}) sont des \textbf{marqueurs d'interaction} standards que l'examinateur attend d'entendre dans une situation de \eng{role-play} ou de \eng{simulated dialogue}.
    \item Un candidat qui dit \eng{“Price?”} (mot unique) perd des points en Interaction. Celui qui dit \eng{“Excuse me, how much is this please?”} gagne des points (politesse + structure complète + intonation).
    \item \textbf{Conseil de pro :} Entraînez-vous à \textbf{ne jamais rester silencieux} plus de 3 secondes. Utilisez des « fillers » anglais : \eng{Well... / Let me see... / Hmm...} au lieu de « Euh... ».
\end{itemize}
\end{savaistu}

\begin{exercice}{Entraînement individuel : « Mon dialogue idéal »}
Rédigez, chez vous, un dialogue complet (10-12 répliques) entre un client et un vendeur au marché de Buea. Le client achète : \eng{2 kg of plantains, 1 kg of tomatoes, 1 bottle of palm oil}. Le vendeur propose une réduction de 10\% pour l'ensemble.
\begin{enumerate}
    \item Utilisez au moins 6 formules fonctionnelles différentes de la leçon.
    \item Indiquez l'intonation (flèche montante $\uparrow$ pour question ouverte, descendante $\downarrow$ pour affirmation/ordre) sur 4 répliques clés.
    \item Traduisez le dialogue en français en regard.
\end{enumerate}
\emph{Cet exercice écrit prépare l'oral : on écrit pour mieux parler ensuite.}
\end{exercice}

\begin{corrige}{Exercice : Mon dialogue idéal}
\textbf{Dialogue type (Customer C / Seller S)} \\
\textbf{C:} \eng{“Good morning.” $\downarrow$} \textbf{S:} \eng{“Good morning, my sister. What can I do for you?” $\downarrow$} \\
\textbf{C:} \eng{“How much is a kilo of plantains?” $\uparrow$} \textbf{S:} \eng{“500 FCFA a kilo. Very ripe today.” $\downarrow$} \\
\textbf{C:} \eng{“And the tomatoes?” $\uparrow$} \textbf{S:} \eng{“600 FCFA a pile.” $\downarrow$} \\
\textbf{C:} \eng{“I'll take 2 kilos of plantains and 2 piles of tomatoes.” $\downarrow$} \textbf{S:} \eng{“That makes 2,200 FCFA.” $\downarrow$} \\
\textbf{C:} \eng{“I also need a bottle of palm oil. How much?” $\uparrow$} \textbf{S:} \eng{“1,200 FCFA. Total 3,400 FCFA.” $\downarrow$} \\
\textbf{C:} \eng{“Can you give me a discount for all?” $\uparrow$} \textbf{S:} \eng{“OK, 10\% off. That's 3,060 FCFA. Say 3,000 even.” $\downarrow$} \\
\textbf{C:} \eng{“Deal. Here is 5,000 FCFA.” $\downarrow$} \textbf{S:} \eng{“Thank you. Here is your change: 2,000 FCFA. Have a nice day!” $\downarrow$} \\
\textbf{C:} \eng{“Thank you very much. Goodbye!” $\downarrow$}
\end{corrige}

\newpage

\coursec{Activité d'intégration}

\courssub{Objectifs de l'activité}
\noindent Cette activité de synthèse vise à mobiliser l'ensemble des savoir-faire travaillés durant la leçon dans une situation de communication authentique, ludique et ancrée dans le quotidien camerounais. À l'issue de cette séance, l'élève doit être capable de :
\begin{itemize}
    \item Saluer poliment et engager la conversation avec un commerçant (\eng{Good morning / Good afternoon}).
    \item Demander des informations précises sur la disponibilité, la qualité, l'origine et le prix de produits alimentaires ou d'articles courants (\eng{How much is...? Do you have...? Is it fresh?}).
    \item Négocier le prix (\emph{bargaining}) en utilisant des formules de relance et de concession (\eng{Can you lower the price? That's too expensive. I'll give you...}).
    \item Conclure la transaction : payer, vérifier la monnaie (\eng{change}) et remercier (\eng{Here is your change. Thank you, have a nice day!}).
    \item Rendre compte par écrit de ses achats en remplissant une fiche de compte rendu structurée.
    \item Respecter la prononciation, l'accentuation des mots-clés (\cle{price}, \cle{change}, \cle{discount}) et l'intonation montante ou descendante des questions et des affirmations.
\end{itemize}

\courssub{Situation}
\noindent Bienvenue au \textbf{Mokolo Market Day} ! La classe est transformée en une réplique animée du grand marché de Mokolo à Yaoundé. L'ambiance est sonore : bruits de \emph{bendskins} qui passent, cris des vendeurs, odeurs d'épices et de poisson braisé.

\begin{textebox}{Mokolo Market Day — Scenario and Shopping Lists}
\textbf{Setting:} The classroom is transformed into Mokolo Market (Yaoundé). Stalls are set up with price tags in English and FCFA. Fake FCFA notes (500, 1,000, 2,000, 5,000) are used.

\textbf{Stalls and Price Tags (Seller's View):}

1. \textbf{Mama Ngo's Foodstuff Stall:} Rice (500 FCFA/kg), Beans (600 FCFA/kg), Groundnuts (800 FCFA/kg), Red Oil (1,500 FCFA/litre), Maggi Cubes (50 FCFA/cube).

2. \textbf{Papa Paul's Provision Store:} Sardines (400 FCFA/tin), Milk (600 FCFA/tin), Sugar (650 FCFA/kg), Flour (550 FCFA/kg), Soap (350 FCFA/bar).

3. \textbf{Brother John's Fish Stall:} Fresh Fish (2,000 FCFA/kg), Smoked Fish (1,500 FCFA/kg), Crayfish (3,000 FCFA/kg).

4. \textbf{Sister Mary's Bread Kiosk:} Bread Loaf (250 FCFA), Meat Pie (300 FCFA), Puff-Puff (50 FCFA/piece), Bottled Water (200 FCFA).

5. \textbf{Uncle Sam's Fruit Corner:} Plantains (100 FCFA/finger), Mangoes (50 FCFA/piece), Oranges (50 FCFA/piece), Pineapple (500 FCFA/whole).

6. \textbf{Auntie Rose's Vegetable Stand:} Tomatoes (200 FCFA/pile), Onions (150 FCFA/pile), Green Vegetables — \emph{Ndolé}/Huckleberry (300 FCFA/bunch), Pepper (100 FCFA/pile).

\textbf{Customer Profiles and Shopping Lists (Buyer's View):}
Each customer receives a list of 3 or 4 items and a budget of 5,000 FCFA.

\textbf{List A:} 1 kg Rice, 1 tin Sardines, 1 loaf Bread, 1 bunch Ndolé. (Approx. 1,450 FCFA)

\textbf{List B:} 1 kg Beans, 1 tin Milk, 500 g Smoked Fish, 5 fingers Plantains. (Approx. 3,150 FCFA)

\textbf{List C:} 1 litre Red Oil, 1 kg Sugar, 1 kg Fresh Fish, 1 Pineapple. (Approx. 4,500 FCFA — tight budget!)

\textbf{List D:} 1 kg Flour, 10 Maggi Cubes, 1 bar Soap, 10 Puff-Puffs. (Approx. 1,750 FCFA)

\textbf{List E:} 1 kg Groundnuts, a small quantity of Crayfish, 10 Oranges, 1 Bottle of Water. (Approx. 2,500 FCFA)
\end{textebox}

\begin{definition}[Glossaire de la situation]
\begin{itemize}
    \item \cle{stall} (n.) : étal, stand de marché. \cle{stall holder} : commerçant qui tient un étal.
    \item \cle{foodstuff} (n. inv.) : denrées alimentaires. Attention, ce mot ne prend pas de \emph{s}.
    \item \cle{provision store} (n.) : boutique d'alimentation générale.
    \item \cle{price tag} (n.) : étiquette de prix.
    \item \cle{tin} (n.) : boîte de conserve (anglais britannique ; \eng{can} en anglais américain).
    \item \cle{loaf} (n.) : pain, miche ; pluriel irrégulier \eng{loaves}.
    \item \cle{finger of plantain} (n.) : doigt de plantain (unité de vente courante au Cameroun).
    \item \cle{pile} (n.) : tas, petit tas (unité de vente pour les tomates, oignons, piments).
    \item \cle{bunch} (n.) : botte (de légumes-feuilles), régime.
    \item \cle{tight budget} (n.) : budget serré.
\end{itemize}
\end{definition}

\courssub{Consignes}
\noindent L'activité se déroule en trois phases successives. Suivez les étapes ci-dessous avec rigueur.

\begin{enumerate}
    \item \textbf{Phase 1 — Mise en place et appropriation (15 min)} \\
    Le professeur désigne 6 \emph{stall holders} (commerçants) qui s'installent avec leurs étiquettes de prix et leur caisse (fausse monnaie). Les autres élèves sont des \emph{customers} (clients) et reçoivent une liste de courses (List A à E) et un portefeuille de 5,000 FCFA (billets découpés). Chaque élève prend 3 minutes pour lire sa liste, repérer les étals concernés et préparer mentalement ses questions (vocabulaire, structures de négociation).

    \item \textbf{Phase 2 — Simulation d'achats : « Let's go shopping! » (45 min)}
    \begin{enumerate}
        \item Les clients circulent librement. Chaque client doit effectuer \textbf{au moins 3 transactions complètes} (dialogue + paiement) auprès de 3 étals différents.
        \item \textbf{Obligation de négociation :} sur au moins \textbf{un} des articles achetés, le client doit tenter d'obtenir une réduction (\eng{discount}) ou un « petit cadeau » (\eng{throw-in}), même symbolique (par exemple un cube Maggi offert, 50 FCFA de moins).
        \item \textbf{Vérification de la monnaie :} à chaque paiement, le client annonce le montant donné, le vendeur rend la monnaie à voix haute (\eng{Here is 500 FCFA change}). Le client vérifie et confirme (\eng{That's correct, thank you}).
        \item \textbf{Fiche de compte rendu :} immédiatement après chaque achat, le client note sur sa fiche \emph{Shopping Report} : Article / Quantity / Asking Price / Final Price / Discount? (Yes/No) / Change Received / Stall Name.
    \end{enumerate}

    \item \textbf{Phase 3 — Inversion des rôles et « Best Scene » (30 min)}
    \begin{enumerate}
        \item Les commerçants deviennent clients (ils reçoivent une nouvelle liste et un budget) et les clients tiennent les étals (ils gèrent le stock et la caisse). Nouvelle session de 20 minutes.
        \item \textbf{Présentation finale :} un binôme volontaire (ou désigné) rejoue sa meilleure scène de négociation devant la classe (\emph{Role-play Performance}).
        \item \textbf{Évaluation par les pairs :} la classe observe avec la \emph{Peer Observation Grid} (grille d'observation) : \textit{Greeting and Politeness / Vocabulary Accuracy / Bargaining Strategies / Pronunciation and Intonation / Transaction Completion / Natural Flow}. Notation sur 6 critères (de 1 à 3 étoiles). Retour constructif oral du professeur.
    \end{enumerate}
\end{enumerate}

\courssub{Ressources à mobiliser}
\noindent Voici un rappel synthétique des outils linguistiques à réactiver pendant l'activité. Gardez cette page accessible.

\begin{propriete}[Functional Language Toolkit — Buying and Selling]
\textbf{1. Greeting and Opening} \\
\eng{Good morning / Good afternoon.} \quad \eng{How can I help you?} (Seller) \\
\eng{Good morning. I'm looking for...} \quad \eng{Do you have...?} (Buyer) \\[4pt]

\textbf{2. Asking for Information (Price, Quality, Quantity)}
\begin{center}
\begin{tabularx}{\linewidth}{|X|X|}\hline
\rowcolor{popblueL} \textbf{Buyer} & \textbf{Seller} \\ \hline
\eng{How much is this? / How much are these?} & \eng{It's 500 francs per kilo / per tin.} \\ \hline
\eng{What is the price of...?} & \eng{They are 200 francs each.} \\ \hline
\eng{Is it fresh? / Is it good quality?} & \eng{Very fresh, it just arrived this morning.} \\ \hline
\eng{Where does it come from?} & \eng{From the West Region / From Douala / It's local.} \\ \hline
\eng{Can I see it? / Can I taste it?} & \eng{Sure, go ahead.} \\ \hline
\end{tabularx}
\end{center}

\textbf{3. Bargaining — Négocier le prix (stratégies progressives)}
\begin{center}
\begin{tabular}{|l|l|}\hline
\rowcolor{poporangeL} \textbf{Étape} & \textbf{Formules types} \\ \hline
\textit{Surprise / Refus} & \eng{Oh! That's too expensive! / That's a bit steep.} \\ \hline
\textit{Contre-proposition} & \eng{Can you lower the price? / Can you do 400? / I'll give you 450.} \\ \hline
\textit{Argumentation} & \eng{I'm a regular customer. / I'm buying in quantity.} \\ \hline
 & \eng{The other stall sells it cheaper.} \\ \hline
\textit{Concession vendeur} & \eng{Okay, last price: 480. / Alright, for you: 450.} \\ \hline
 & \eng{I can't go lower.} \\ \hline
\textit{Clôture} & \eng{Deal! / Agreed. / No, thank you, I'll leave it.} \\ \hline
\end{tabular}
\end{center}

\textbf{4. Quantity and Packaging (unités de mesure courantes au Cameroun)} \\
\eng{A kilo of... / Half a kilo of... / A pile of... (tomatoes, onions) / A bunch of... (ndolé, huckleberry) / A tin of... / A bottle of... / A loaf of... / A piece of... / A finger of plantain.} \\[4pt]

\textbf{5. Payment and Closing}
\begin{center}
\begin{tabularx}{\linewidth}{|X|X|}\hline
\rowcolor{popgreenL} \textbf{Buyer} & \textbf{Seller} \\ \hline
\eng{Here is 1,000 francs.} & \eng{That makes 850 francs. Here is your change: 150 francs.} \\ \hline
\eng{Keep the change. / Is that correct?} & \eng{Thank you! Have a nice day! / Come again!} \\ \hline
\end{tabularx}
\end{center}
\end{propriete}

\begin{attention}[Pièges fréquents pour francophones — à surveiller pendant l'activité]
\begin{itemize}
    \item \textbf{Demander le prix :} on dit \eng{How much is it?} ou \eng{How much does it cost?}, jamais \eng{« It costs how much? »}. Le verbe \eng{cost} s'emploie surtout dans les affirmations avec un sujet inanimé : \eng{It costs 500 francs.} (« Cela coûte 500 francs. »)
    \item \textbf{Monnaie :} \cle{change} est \textbf{invariable} (pas de \emph{s}). On dit \eng{Here is your change} (singulier). Ne pas confondre \eng{change} (la monnaie rendue) et \eng{coins} (les pièces de monnaie).
    \item \textbf{Quantifieurs :} \eng{How much} pour l'indénombrable (\eng{rice, oil, sugar}) ; \eng{How many} pour le dénombrable (\eng{tins, loaves, mangoes}). Exemples : \eng{How much rice do you want?} (« Combien de riz voulez-vous ? ») ; \eng{How many tins do you need?} (« Combien de boîtes vous faut-il ? »).
    \item \textbf{Articles partitifs :} \eng{a kilo \textbf{of} rice}, \eng{a tin \textbf{of} milk}, \eng{a loaf \textbf{of} bread}. Ne jamais oublier la préposition \eng{of}.
    \item \textbf{Prononciation :} \eng{price} se prononce « praïss » (pas « prise ») ; \eng{change} se prononce « tchéïndj » ; \eng{bargain} se prononce « ba-gueune » avec l'accent sur la première syllabe ; \eng{receipt} se prononce « ri-ssite » (le \eng{p} est muet).
    \item \textbf{Politesse :} toujours ouvrir par \eng{Good morning} ou \eng{Good afternoon} et fermer par \eng{Thank you / Have a nice day}. L'omission de ces formules est perçue comme impolie en contexte anglophone.
\end{itemize}
\end{attention}

\begin{savaistu}[Le saviez-vous ? — Le « bargaining » au Cameroun]
Au marché Mokolo (Yaoundé) ou au Marché Central (Douala), la négociation (\emph{bargaining} ou \emph{haggling}) est une norme culturelle, pas une impolitesse. Le prix affiché (\eng{asking price} ou \eng{marked price}) est souvent un prix de départ (\eng{starting price}) : le client s'attend à obtenir une petite réduction.
\begin{itemize}
    \item \textbf{Expressions locales :} \eng{« Give me my price »} (« fais-moi mon prix ») ou \eng{« Do me well »} (« traite-moi bien, fais-moi un bon prix »).
    \item \textbf{Le « dash » :} petit cadeau ajouté par le vendeur pour fidéliser le client (par exemple un oignon en plus, un cube Maggi). En anglais camerounais : \eng{« Put something small for me »} ou \eng{« Dash me one »}.
    \item \textbf{Argent :} on parle de \eng{FCFA} (\eng{francs}). Les billets courants sont de 500, 1,000, 2,000, 5,000 et 10,000 francs. Vérifier sa monnaie (\eng{count your change}) est une habitude universelle.
\end{itemize}
\end{savaistu}

\courssub{Proposition de production et corrigé commenté}
\noindent Voici un modèle de dialogue complet et réaliste (niveau A2+/B1) illustrant une transaction réussie avec négociation, basé sur la \textbf{Liste C} (client) et l'étal \textbf{Mama Ngo's Foodstuff Stall} (vendeuse). Ce dialogue sert de référence pour l'évaluation finale (\emph{Best Scene}).

\begin{exoresolu}{Modèle de dialogue : achat d'huile rouge et de riz — négociation réussie}
\textbf{Contexte :} le client (Liste C : 1 litre de red oil, 1 kg de rice, 1 kg de fresh fish, 1 pineapple ; budget : 5,000 FCFA) se trouve à l'étal de Mama Ngo pour l'huile et le riz. Prix affichés : red oil 1,500 FCFA/litre, rice 500 FCFA/kg. Total affiché : 2,000 FCFA.

\textbf{Dialogue :}

\textbf{Seller (Mama Ngo):} \eng{Good afternoon, my sister! Welcome to my stall. How can I help you today?}

\textbf{Buyer:} \eng{Good afternoon, Mama. I am looking for red oil and rice. Do you have good quality red oil?}

\textbf{Seller:} \eng{Yes, my son. This one is from the West Region, very red and thick. Best quality in Mokolo! It is 1,500 francs per litre. The rice is 500 francs per kilo. Local rice, very clean.}

\textbf{Buyer:} \eng{One litre of red oil and one kilo of rice, please. How much is that altogether?}

\textbf{Seller:} \eng{That makes 2,000 francs exactly.}

\textbf{Buyer:} \eng{Oh! 2,000 francs? That is a bit steep, Mama. Can you lower the price? I am a regular customer. I'll give you 1,800 francs for both.}

\textbf{Seller:} \eng{Eh, my son, you know the market is hard today. Transport has increased. Last price: 1,900 francs. I cannot go lower.}

\textbf{Buyer:} \eng{Okay, 1,900 francs. Deal. Here is 2,000 francs.}

\textbf{Seller:} \eng{Here is your change: 100 francs. Thank you very much! Come again tomorrow!}

\textbf{Buyer:} \eng{Thank you, Mama. Have a nice afternoon! Bye!}

\tcblower

\textbf{Commentaire détaillé des choix de langue (corrigé commenté) :}

\begin{enumerate}
    \item \textbf{Accroche culturelle (\eng{my sister / my son}) :} très courante au Cameroun anglophone et francophone pour créer du lien social (\emph{rapport}). Elle remplace \eng{Sir/Madam}, trop formel. À reproduire.
    \item \textbf{Demande de prix global (\eng{How much is that altogether?}) :} structure clé ; \eng{altogether} signifie « en tout, au total ». Elle évite de devoir additionner mentalement en anglais.
    \item \textbf{Expression de surprise polie (\eng{That is a bit steep}) :} \cle{steep} (adj.) signifie « élevé, excessif » pour un prix. C'est plus naturel que \eng{too expensive} seul ; \eng{a bit} adoucit la critique.
    \item \textbf{Stratégie \eng{regular customer} :} argument standard pour justifier une réduction. \eng{Regular} signifie « habituel, fidèle » (client régulier), et non « ordinaire ».
    \item \textbf{Contre-proposition directe (\eng{I'll give you 1,800}) :} la structure \eng{I'll give you + montant} correspond à « je vous en donne... ». Très idiomatique pour le marchandage.
    \item \textbf{Réplique du vendeur (\eng{Last price}) :} formule figée signifiant « prix final, non négociable ». \eng{I cannot go lower} signifie « je ne peux pas descendre plus bas ».
    \item \textbf{Clôture (\eng{Deal}) :} court, ferme, scelle l'accord oral. Équivalent de « C'est entendu / Marché conclu ».
    \item \textbf{Paiement (\eng{Here is 2,000 francs}) :} présentatif \eng{here is/are} + montant donné. \textbf{Piège évité :} ne pas dire \eng{« I give you 2,000 »}.
    \item \textbf{Monnaie (\eng{Here is your change: 100 francs}) :} \cle{change} reste au singulier. Le vendeur annonce le montant rendu \textbf{avant} de le tendre : c'est la bonne pratique.
    \item \textbf{Formules de clôture symétriques :} \eng{Thank you / Have a nice afternoon / Come again}. Rituel de politesse indispensable.
\end{enumerate}

\textbf{Fiche de compte rendu associée (Shopping Report) :}
\begin{center}
\begin{tabular}{|l|c|c|c|c|c|c|}\hline
\rowcolor{popblueL} \textbf{Article} & \textbf{Qty} & \textbf{Asking Price} & \textbf{Final Price} & \textbf{Discount?} & \textbf{Change} & \textbf{Stall} \\ \hline
Red Oil & 1 L & 1,500 & 1,425* & Yes & 100 FCFA & Mama Ngo \\ \hline
Rice & 1 kg & 500 & 475* & Yes & --- & Mama Ngo \\ \hline
\end{tabular}
\end{center}
\eng{*Negotiated global price of 1,900 FCFA, split proportionally for the report.} (« Prix global négocié de 1 900 FCFA, réparti proportionnellement pour le compte rendu. »)
\end{exoresolu}

\courssub{Bilan de l'activité}
\noindent Cette activité d'intégration \textbf{Mokolo Market Day} a permis de valider les compétences orales ciblées dans un environnement simulé mais réaliste, ancré dans le contexte sociolinguistique camerounais.

\begin{aretenir}[Bilan des acquis — grille d'auto-évaluation élève]
À la fin de la séance, cochez ce que vous savez faire \textbf{seul(e)} en anglais :
\begin{itemize}
    \item[$\square$] Saluer et prendre congé poliment (\eng{Good morning / Have a nice day}).
    \item[$\square$] Demander le prix (\eng{How much is...?}) et la qualité (\eng{Is it fresh?}).
    \item[$\square$] Comprendre le prix annoncé par le vendeur (écoute de nombres en FCFA).
    \item[$\square$] Négocier : dire « c'est cher », proposer un prix, argumenter (\eng{regular customer}).
    \item[$\square$] Accepter (\eng{Deal}) ou refuser (\eng{No, thank you}) l'offre finale.
    \item[$\square$] Payer, vérifier la monnaie (\eng{change}) et signaler une erreur le cas échéant.
    \item[$\square$] Remplir une fiche de compte rendu en anglais (vocabulaire des articles, nombres, \eng{Yes/No}).
    \item[$\square$] Prononcer correctement : \eng{price, change, bargain, receipt, thousand, hundred}.
\end{itemize}
\end{aretenir}

\begin{methode}[Pour progresser — pistes de travail personnel]
\begin{enumerate}
    \item \textbf{Entraînement aux nombres :} dictez-vous des prix au hasard (par exemple 1,750 ; 3,200 ; 10,500) et écrivez-les en toutes lettres (\eng{one thousand seven hundred and fifty}). Révisez la prononciation du \eng{th} (\eng{thousand, three}) et la liaison dans \eng{hundred and}.
    \item \textbf{Flashcards « Marché » :} faites des cartes recto-verso : image de l'article (tomate, savon, poisson) au recto, expression au verso : \eng{a pile of tomatoes / a bar of soap / a kilo of fish}. Révisez les partitifs avec \eng{of}.
    \item \textbf{Jeu de l'ombre (\emph{shadowing}) :} réécoutez le dialogue modèle (ou enregistrez-vous en binôme) et répétez en même temps que l'audio en imitant l'intonation : montante pour \eng{Can you lower the price?}, descendante pour \eng{Deal.}
    \item \textbf{Sortie réelle :} lors de votre prochain passage au marché (Mokolo, Nkolndongo, Mfoundi...), essayez de formuler mentalement la transaction en anglais. Notez trois nouvelles expressions entendues en anglais local et traduisez-les en français.
\end{enumerate}
\end{methode}

\begin{experience}[Extension possible — « Market Survey Project » (projet sur deux semaines)]
En groupes de 4, menez une mini-enquête hors classe : \eng{Price Comparison: Supermarket vs. Open Market}.
\begin{itemize}
    \item Choisissez 5 produits identiques (par exemple 1 kg de rice, 1 litre d'oil, 1 tin de milk, 1 loaf de bread, 1 kg de sugar).
    \item Relevez les prix dans un supermarché (par exemple Carrefour, Mahima, Spar) et au marché (Mokolo, Marché Central).
    \item Présentez un poster bilingue (anglais/français) avec un tableau comparatif, un graphique simple et une conclusion orale de 2 minutes en anglais : \eng{« We found that rice is cheaper at the market because... The supermarket is better for... »}.
    \item Évaluation : vocabulaire spécifique, comparatifs (\eng{cheaper than, more expensive than, the best value for money}), prononciation, travail d'équipe.
\end{itemize}
\end{experience}

\newpage

\coursec{Méthodes et exercices résolus}

\coursec{Lesson 13 — Speaking: Exchanges information about buying and selling articles/foodstuff at a nearby store/market}

\courssub{Échanger des informations dans une situation d'achat et de vente}

\begin{methode}[Mener un échange complet au magasin ou au marché]
Pour réussir un dialogue d'achat-vente à l'oral, suis toujours la même architecture en cinq étapes :
\begin{enumerate}
\item \textbf{Saluer et ouvrir l'échange} : \eng{“Good morning! Can I help you?”} « Bonjour ! Puis-je vous aider ? » / \eng{“Good morning! I'm looking for some tomatoes.”} « Bonjour ! Je cherche des tomates. » Ne commence jamais par la demande brute : la salutation est obligatoire dans la culture anglophone.
\item \textbf{Demander l'information} : utilise les structures interrogatives correctes : \eng{“How much is this bag of rice?”} « Combien coûte ce sac de riz ? », \eng{“Do you have fresh fish?”} « Avez-vous du poisson frais ? », \eng{“How many mangoes do you want?”} « Combien de mangues voulez-vous ? » Attention à l'auxiliaire \cle{do/does} et à l'ordre sujet-verbe.
\item \textbf{Donner l'information} : réponds de façon complète : \eng{“It costs 500 francs a kilo.”} « Ça coûte 500 francs le kilo. », \eng{“We have some, but they are not very fresh today.”} « Nous en avons, mais ils ne sont pas très frais aujourd'hui. »
\item \textbf{Négocier et relancer} : \eng{“That's too expensive! Can you reduce the price?”} « C'est trop cher ! Pouvez-vous baisser le prix ? », \eng{“What about 400 francs?”} « Et 400 francs ? », \eng{“If you buy two kilos, I'll give you a discount.”} « Si vous achetez deux kilos, je vous ferai une remise. »
\item \textbf{Conclure poliment} : \eng{“Fine, I'll take it. Here is the money.”} « D'accord, je le prends. Voici l'argent. », \eng{“Thank you very much. Have a nice day!”} « Merci beaucoup. Bonne journée ! », \eng{“Goodbye, madam!”} « Au revoir, madame ! »
\end{enumerate}
Réflexe : prépare toujours mentalement ces cinq blocs avant de parler ; un dialogue sans salutation ni conclusion est pénalisé à l'examen.
\end{methode}

\begin{exoresolu}[Compléter un dialogue au marché de Mokolo]
Énoncé — Complete the dialogue between a customer and a trader at Mokolo market (Yaoundé) with the missing sentences. Choose from: \eng{“How much do they cost?”} / \eng{“That's too expensive!”} / \eng{“Good morning, madam!”} / \eng{“I'll take two kilos then.”} / \eng{“Can I help you?”}

\eng{Trader: (1) \ldots}
\eng{Customer: Good morning! I need some onions.}
\eng{Trader: Of course. These are very fresh.}
\eng{Customer: (2) \ldots}
\eng{Trader: 800 francs a kilo.}
\eng{Customer: (3) \ldots In Douala I paid 600 francs.}
\eng{Trader: This is Yaoundé, my dear! But for you, 700 francs.}
\eng{Customer: (4) \ldots}
\eng{Trader: Here you are. Anything else?}
\eng{Customer: No, thank you. Goodbye!}
\tcblower
Solution détaillée :
\begin{enumerate}
\item \eng{“Can I help you?”} — C'est la formule d'accueil type du vendeur. Le client répond par sa demande (\eng{“I need some onions”}), donc le trader doit ouvrir l'échange. \eng{“Good morning, madam!”} serait possible mais la case (1) exige la relance du vendeur, et la salutation convient mieux en ouverture globale.
\item \eng{“How much do they cost?”} — La réponse du trader est un prix (\eng{“800 francs a kilo”}), donc la question porte sur le prix. \eng{Onions} est pluriel, d'où \cle{they} et non \emph{it}.
\item \eng{“That's too expensive!”} — Le client compare avec Douala et proteste : c'est la phase de négociation (bargaining).
\item \eng{“I'll take two kilos then.”} — Après la réduction à 700 francs, le client conclut l'achat. \eng{I'll take} = futur de décision spontanée (\cle{will} + base verbale).
\end{enumerate}
Il reste \eng{“Good morning, madam!”} qui pourrait ouvrir le dialogue ; la formulation la plus naturelle est : Trader: \eng{“Good morning, madam! Can I help you?”} — mais comme une seule proposition par case est demandée, la réponse attendue en (1) est \eng{“Can I help you?”}, la salutation étant déjà implicite dans la réponse du client. Conclusion : toujours vérifier que chaque réplique choisie est justifiée par la réplique qui suit (prix $\rightarrow$ question sur le prix ; réduction $\rightarrow$ décision d'achat).
\end{exoresolu}

\courssub{Poser des questions et répondre de façon adaptée}

\begin{methode}[Questionner et répondre correctement dans un échange commercial]
\begin{enumerate}
\item \textbf{Identifier le type de question} : question fermée (réponse yes/no) avec inversion de l'auxiliaire : \eng{“Do you sell groundnuts?”} « Vendez-vous des arachides ? » ; question ouverte avec mot interrogatif en tête : \eng{“How much is a bottle of palm oil?”} « Combien coûte une bouteille d'huile de palme ? », \eng{“Where can I find fresh pepper?”} « Où puis-je trouver du piment frais ? »
\item \textbf{Choisir le bon mot interrogatif} : \cle{how much} + nom indénombrable (\eng{rice, oil, sugar}) ; \cle{how many} + nom dénombrable pluriel (\eng{eggs, oranges, yams}) ; \cle{which} pour un choix limité (\eng{“Which one do you prefer?”} « Lequel préférez-vous ? »).
\item \textbf{Construire la question} : mot interrogatif + auxiliaire (\cle{do/does/is/are/can}) + sujet + verbe. Jamais d'inversion sans auxiliaire : \eng{“How much costs this?”} est FAUX ; dire \eng{“How much does this cost?”}
\item \textbf{Répondre de façon adaptée} : à une question fermée, réponds d'abord \eng{“Yes, we do. / No, we don't.”} puis ajoute une information ; à une question ouverte, donne directement l'information demandée : \eng{“It's 1,500 francs.”}
\item \textbf{Relancer si la réponse est incomplète} : \eng{“Sorry, how much did you say?”} « Désolé, combien avez-vous dit ? », \eng{“And what about the plantains?”} « Et pour les plantains ? »
\end{enumerate}
\end{methode}

\begin{exoresolu}[Transformer des demandes en questions correctes]
Énoncé — You are at a store near your school. Turn these ideas into correct English questions, then give a possible short answer.
\begin{enumerate}
\item (demander le prix du savon)
\item (demander s'ils vendent du lait en poudre)
\item (demander combien d'œufs il reste)
\item (demander où sont les boîtes de sardines)
\end{enumerate}
\tcblower
Solution détaillée :
\begin{enumerate}
\item \eng{“How much is the soap?”} / \eng{“How much does this soap cost?”} — Réponse possible : \eng{“It's 250 francs.”} Justification : prix $\rightarrow$ \cle{how much} ; avec le verbe \emph{to be}, simple inversion ; avec \emph{cost}, auxiliaire \cle{does} obligatoire (sujet singulier).
\item \eng{“Do you sell powdered milk?”} — Réponse : \eng{“Yes, we do. It's on the second shelf.”} Question fermée $\rightarrow$ \cle{do} + sujet + base verbale. Ne pas dire \eng{“Sell you...?”} (calque du français « vendez-vous... ? »).
\item \eng{“How many eggs are left?”} / \eng{“How many eggs do you have left?”} — Réponse : \eng{“There are ten eggs left.”} \emph{Eggs} est dénombrable pluriel $\rightarrow$ \cle{how many}, jamais \emph{how much}.
\item \eng{“Where are the tins of sardines?”} / \eng{“Where can I find the tinned sardines?”} — Réponse : \eng{“They are at the back of the store, next to the rice.”} \emph{Tins} est pluriel $\rightarrow$ \cle{are}.
\end{enumerate}
Conclusion : la bonne question dépend de trois choix successifs : mot interrogatif (prix, quantité, lieu), auxiliaire adapté au sujet, et ordre des mots interrogatif. S'entraîner à verbaliser ces trois choix avant de parler.
\end{exoresolu}

\courssub{Négocier, relancer et conclure : le bargaining}

\begin{methode}[Négocier un prix en anglais correct]
\begin{enumerate}
\item \textbf{Exprimer que le prix est élevé} sans agressivité : \eng{“That's a bit expensive.”} « C'est un peu cher. », \eng{“Is that your last price?”} (formule idiomatique du marché : « C'est votre dernier prix ? »).
\item \textbf{Faire une contre-offre} avec \cle{what about} / \cle{how about} + nom ou \emph{-ing} : \eng{“What about 300 francs?”} « Et 300 francs ? », \eng{“How about reducing it to 300?”} « Et si on le réduisait à 300 ? »
\item \textbf{Argumenter} : \eng{“If I buy three, will you give me a discount?”} « Si j'en achète trois, me ferez-vous une remise ? », \eng{“The tomatoes are not very fresh.”} « Les tomates ne sont pas très fraîches. »
\item \textbf{Accepter ou refuser} : \eng{“OK, that's fine. I'll take it.”} « D'accord, c'est bon. Je le prends. » / \eng{“No, thank you. That's still too much. Maybe next time.”} « Non, merci. C'est encore trop. Peut-être la prochaine fois. »
\item \textbf{Conclure la transaction} : payer (\eng{“Here you are.”} / \eng{“Keep the change.”} « Gardez la monnaie. »), remercier, saluer : \eng{“Thanks a lot. See you next week!”} « Merci beaucoup. À la semaine prochaine ! »
\end{enumerate}
Réflexe de prononciation : la négociation se dit avec une intonation montante sur la contre-offre (\eng{“What about three hundred?”}) et descendante sur la décision finale (\eng{“I'll take it.”}).
\end{methode}

\begin{exoresolu}[Rédiger la phase de négociation d'un dialogue]
Énoncé — At Buea market, a customer wants to buy a bunch of plantains. The trader asks for 2,000 francs; the customer finds it expensive and offers 1,500 francs; they finally agree on 1,700 francs. Write the bargaining part of the dialogue (at least 6 lines).
\tcblower
Solution détaillée — Raisonnement : il faut respecter la progression protestation $\rightarrow$ contre-offre $\rightarrow$ argument $\rightarrow$ compromis $\rightarrow$ conclusion, avec des formules variées et un anglais correct.

\eng{Customer: “Two thousand francs? That's too expensive for a bunch of plantains!”}
\eng{Trader: “But look at the size, my friend. They are very fresh.”}
\eng{Customer: “What about one thousand five hundred?”}
\eng{Trader: “No, that's too low. I bought them at a high price myself.”}
\eng{Customer: “Come on! If you accept one thousand seven hundred, I'll take the bunch right now.”}
\eng{Trader: “OK, deal! One thousand seven hundred it is.”}
\eng{Customer: “Fine. Here is the money. Thank you very much!”}
\eng{Trader: “You're welcome. Come again!”}

Justifications grammaticales : \eng{“That's too expensive”} (\cle{too} + adjectif = « trop ») ; \eng{“What about”} + complément nominal pour la contre-offre ; \eng{“If you accept..., I'll take...”} = premier conditionnel (\cle{if} + présent, \cle{will} + base verbale) pour la promesse d'achat ; \eng{“Here is the money”} : \emph{money} est indénombrable, donc singulier. Conclusion : une négociation réussie alterne répliques courtes et naturelles ; éviter les phrases longues et littéraires.
\end{exoresolu}

\courssub{Mener un court dialogue en anglais : le jeu de rôle}

\begin{methode}[Préparer et jouer un rôle (customer ou trader)]
\begin{enumerate}
\item \textbf{Lire la consigne et repérer ton rôle} : qui es-tu (acheteur/vendeur), que veux-tu (quel article, quel budget), où es-tu (store, market) ?
\item \textbf{Préparer 5 à 8 répliques} en suivant l'architecture : salutation $\rightarrow$ demande $\rightarrow$ information/prix $\rightarrow$ négociation $\rightarrow$ conclusion.
\item \textbf{Mémoriser les formules-clés, pas des phrases entières} : \eng{“Can I help you?”} / \eng{“How much...?”} / \eng{“That's too expensive.”} / \eng{“What about...?”} / \eng{“I'll take it.”} / \eng{“Keep the change.”}
\item \textbf{Soigner la prononciation} : accentue les mots porteurs de sens (\eng{“HOW much is THIS bag?”}), monte le ton dans les questions fermées (\eng{“Do you have FISH?”}), descends-le dans les questions en \emph{wh-} (\eng{“HOW much is it?”}).
\item \textbf{Improviser si le partenaire sort du script} : relance avec \eng{“Sorry, can you repeat that?”}, \eng{“Pardon?”}, \eng{“Do you mean...?”}
\end{enumerate}
\end{methode}

\begin{exoresolu}[Jeu de rôle noté : au magasin du quartier]
Énoncé — Role-play: Student A is a shopkeeper at a provision store in Bamenda. Student B is a customer who wants to buy a loaf of bread, a tin of milk and half a kilo of sugar, with only 1,000 francs. Perform the dialogue. (Assessment: greeting, correct questions, bargaining, conclusion.)
\tcblower
Solution détaillée — Modèle de performance attendue :

\eng{A: “Good afternoon! Can I help you?”}
\eng{B: “Good afternoon. How much is a loaf of bread?”}
\eng{A: “It's 300 francs.”}
\eng{B: “And a tin of milk?”}
\eng{A: “450 francs.”}
\eng{B: “Oh, and half a kilo of sugar?”}
\eng{A: “That's 350 francs. So, 1,100 francs altogether.”}
\eng{B: “Oh no! I only have 1,000 francs. Can you reduce the price a little?”}
\eng{A: “Hmm... OK, because you are a regular customer: 1,000 francs.”}
\eng{B: “Thank you so much! Here you are.”}
\eng{A: “You're welcome. Have a nice evening!”}

Analyse : la salutation ouvre et clôt l'échange ; les questions de prix utilisent \cle{how much} + inversion (\eng{“How much is...?”}) ; le problème du budget crée une négociation naturelle (\eng{“Can you reduce the price?”}) ; le vendeur conclut avec une justification (\eng{“because you are a regular customer”}). Prononciation : \eng{bread} « brède », \eng{sugar} « chou-gar », \eng{altogether} « ol-to-gué-zar ». Conclusion : pour être bien noté, il faut les quatre étapes + une prononciation claire, pas un dialogue long.
\end{exoresolu}

\courssub{Prononciation, accentuation et intonation dans le dialogue commercial}

\begin{methode}[Soigner l'intonation et les sons difficiles]
\begin{enumerate}
\item \textbf{Questions fermées (yes/no)} : intonation montante à la fin : \eng{“Do you sell YAMS?”} (le ton monte sur « yams »).
\item \textbf{Questions ouvertes (wh-)} : intonation descendante : \eng{“HOW much IS it?”} (le ton descend sur « it »).
\item \textbf{Accent de phrase} : accentuer les mots porteurs de sens (noms, nombres, adjectifs), pas les petits mots : \eng{“I'll take TWO KILOS of RICE.”}
\item \textbf{Sons-pièges pour les francophones} : \emph{th} de \eng{three, thousand} (langue entre les dents, « sss » soufflé, jamais « tri ») ; \emph{h} aspiré de \eng{how, hundred} (« haou, ran-dred ») ; \eng{cheap} « tchipe » (ne pas confondre avec \eng{sheep} « chipe ») ; \eng{change} « tchèndj ».
\item \textbf{Nombres} : bien distinguer \eng{thirteen} « seur-tine » (accent final) et \eng{thirty} « seur-ti » (accent initial) — confusion fréquente dans les prix !
\end{enumerate}
\end{methode}

\begin{exoresolu}[Repérer et corriger des erreurs de prononciation et d'intonation]
Énoncé — A student says: « Tri sousand frans? Zat is too expensive! » and asks \eng{“How much is it?”} with a rising intonation. Identify and correct the mistakes, then transcribe the correct pronunciation in French letters.
\tcblower
Solution détaillée :
\begin{enumerate}
\item « Tri » $\rightarrow$ \eng{three} se prononce « ssri » : le \emph{th} est une fricative dentale sourde (langue entre les dents), pas un \emph{t} français.
\item « sousand » $\rightarrow$ \eng{thousand} : « ssaou-zand » ; le \emph{th} initial subit la même correction, et \emph{ou} se dit « aou ».
\item « frans » $\rightarrow$ \eng{francs} : « fran-kss » ; le \emph{c} final se prononce.
\item « Zat » $\rightarrow$ \eng{that} : « zatte » avec \emph{th} sonore (vibration), pas un simple « z ».
\item Intonation : \eng{“How much is it?”} est une question ouverte (mot interrogatif \emph{how much}), donc l'intonation doit être \textbf{descendante} sur « it ». L'intonation montante est réservée aux questions fermées : \eng{“Is it fresh?”} (ton montant).
\end{enumerate}
Phrase corrigée : \eng{“Three thousand francs? That is too expensive!”} — « ssri ssaou-zand fran-kss ? zatte iz tou èk-spè-n-siv ! » Conclusion : à l'oral, une mauvaise intonation peut transformer une question en affirmation ; s'entraîner à enregistrer ses dialogues pour vérifier montées et descentes du ton.
\end{exoresolu}

\begin{attention}[Les 5 erreurs qui coûtent des points au Bac]
\begin{enumerate}
\item \textbf{Question sans auxiliaire} : \eng{“How much costs this?”} est faux. Dire \eng{“How much does this cost?”} ou \eng{“How much is this?”} Toute question avec un verbe lexical exige \cle{do/does/did}.
\item \textbf{Confusion how much / how many} : \eng{“How many rice?”} est faux : \emph{rice, sugar, oil, money} sont indénombrables $\rightarrow$ \cle{how much}. Réserve \cle{how many} aux dénombrables pluriels (\eng{“how many eggs”}).
\item \textbf{Calque du français} : \eng{“It is how much?”} (calque de « c'est combien ? ») et \eng{“Give me reduction”} (calque de « faites-moi une réduction »). Dire \eng{“How much is it?”} et \eng{“Can you give me a discount?”} / \eng{“Can you reduce the price?”}
\item \textbf{Faux amis et vocabulaire} : \emph{actually} ne veut pas dire « actuellement » (\eng{currently}) ; \eng{“I want to buy some devoirs”} — attention à ne pas franciser : « provisions » se dit \eng{foodstuff / groceries} ; « un magasin » = \eng{a store/shop}, pas \emph{a magazine} (qui signifie « une revue »).
\item \textbf{Pluriels et indénombrables} : \eng{“some breads”}, \eng{“many moneys”}, \eng{“informations”} sont faux. Dire \eng{“some bread / some loaves of bread”}, \eng{“much money”}, \eng{“some information”}. Et dans les prix : \eng{“It costs 500 francs”} — n'oublie pas le \cle{-s} de la troisième personne au présent simple.
\end{enumerate}
\end{attention}

\newpage

\coursec{Exercices}

\courssub{Je vérifie mes connaissances}

\begin{exercice}[QCM : les bonnes formules]
Pour chaque situation, choisis la formule anglaise correcte (a, b ou c).

\begin{enumerate}
\item Tu entres dans une boutique et tu veux attirer poliment l'attention du vendeur :
\begin{enumerate}
\item \eng{Hey you! Come here!}
\item \eng{Excuse me, could you help me, please?}
\item \eng{Give me your attention now!}
\end{enumerate}
\item Tu veux connaître le prix des tomates :
\begin{enumerate}
\item \eng{How many are the tomatoes?}
\item \eng{How much are the tomatoes?}
\item \eng{What is the tomatoes' cost price?}
\end{enumerate}
\item Le vendeur te demande si tu veux autre chose :
\begin{enumerate}
\item \eng{Anything else?}
\item \eng{Something more again?}
\item \eng{Do you want else?}
\end{enumerate}
\item Tu trouves le prix trop élevé et tu veux négocier :
\begin{enumerate}
\item \eng{It is too much expensive, I refuse!}
\item \eng{That's a bit expensive. Can you give me a discount?}
\item \eng{Your price is a robbery!}
\end{enumerate}
\end{enumerate}
\end{exercice}

\begin{exercice}[Vrai ou faux ? Justifie]
Réponds par \eng{True} ou \eng{False} et justifie ta réponse en une phrase en anglais.

\begin{enumerate}
\item \eng{We use “How much...?” with uncountable nouns like rice or oil.}
\item \eng{“How many...?” is used with uncountable nouns like money.}
\item \eng{“Here is your change” is said by the customer.}
\item \eng{“I'm just looking, thank you” is a polite way to refuse help in a shop.}
\item \eng{In a Cameroonian market, bargaining is considered rude and must be avoided.}
\end{enumerate}
\end{exercice}

\begin{exercice}[Vocabulaire : au magasin et au marché]
Complète chaque phrase avec le mot anglais qui convient, choisi dans la liste : \eng{stall — change — receipt — discount — customer — trader — bag — price}.

\begin{enumerate}
\item The \ldots\ldots\ldots\ldots sells fruit and vegetables at the market.
\item I paid 5,000 FCFA and the shopkeeper gave me 1,500 FCFA \ldots\ldots\ldots\ldots .
\item Can I have a \ldots\ldots\ldots\ldots to carry my tomatoes, please?
\item If you buy three kilos, I will give you a \ldots\ldots\ldots\ldots of 200 FCFA.
\item She sells second-hand clothes at a small \ldots\ldots\ldots\ldots near the bus station.
\item Always ask for a \ldots\ldots\ldots\ldots when you buy something in a supermarket.
\end{enumerate}
\end{exercice}

\begin{exercice}[Grammaire : much, many, some, any]
Complète les questions et réponses avec \eng{much}, \eng{many}, \eng{some} ou \eng{any}.

\begin{enumerate}
\item How \ldots\ldots\ldots\ldots is this bag of rice? — It is 6,000 FCFA.
\item How \ldots\ldots\ldots\ldots oranges do you want? — Five, please.
\item Do you have \ldots\ldots\ldots\ldots groundnut oil? — Yes, we have \ldots\ldots\ldots\ldots .
\item I'm sorry, we don't have \ldots\ldots\ldots\ldots fresh fish today.
\item Could I have \ldots\ldots\ldots\ldots bread, please?
\end{enumerate}
\end{exercice}

\courssub{Je m'entraîne}

\begin{exercice}[Traduction : au marché de Mokolo]
Traduis les phrases suivantes en anglais.

\begin{enumerate}
\item Combien coûte un kilo de tomates ?
\item C'est trop cher. Pouvez-vous baisser le prix ?
\item Je prends deux kilos de riz et un litre d'huile.
\item Voulez-vous autre chose, madame ?
\item Voici votre monnaie. Merci et à bientôt !
\end{enumerate}
\end{exercice}

\begin{exercice}[Réagir à des situations]
Écris en anglais la phrase adaptée à chaque situation, en utilisant les formules fonctionnelles de la leçon.

\begin{enumerate}
\item Tu demandes le prix d'un ananas.
\item Tu dis poliment que le prix est trop élevé.
\item Tu demandes au vendeur de te rendre la monnaie sur 10,000 FCFA.
\item Tu refuses poliment l'aide d'un vendeur car tu regardes seulement.
\item Tu demandes si le magasin vend du savon.
\end{enumerate}
\end{exercice}

\begin{textebox}[At the provision store — dialogue à compléter]
Seller: \ldots\ldots\ldots\ldots\ldots\ldots\ldots\ldots\ldots\ldots\ldots\ldots\ldots\ldots ?

Customer: Good morning. I would like some tomatoes, please.

Seller: How many kilos do you want?

Customer: Two kilos. \ldots\ldots\ldots\ldots\ldots\ldots\ldots\ldots\ldots\ldots\ldots\ldots\ldots\ldots ?

Seller: They are 800 FCFA a kilo.

Customer: That is expensive! \ldots\ldots\ldots\ldots\ldots\ldots\ldots\ldots\ldots\ldots\ldots\ldots\ldots\ldots ?

Seller: All right. For you, 1,400 FCFA for the two kilos.

Customer: Fine. Here are 2,000 FCFA.

Seller: \ldots\ldots\ldots\ldots\ldots\ldots\ldots\ldots\ldots\ldots\ldots\ldots\ldots\ldots ?

Customer: No, thank you. That's all.

Seller: \ldots\ldots\ldots\ldots\ldots\ldots\ldots\ldots\ldots\ldots\ldots\ldots\ldots\ldots . Goodbye!

Customer: Thank you. Goodbye!
\end{textebox}

\begin{exercice}[Dialogue à trous]
Complète le dialogue ci-dessus entre un client et un vendeur avec les formules suivantes : \eng{Anything else? — How much are they? — Here is your change. — Can you give me a discount? — Good morning, can I help you?}
\end{exercice}

\begin{exercice}[Remettre le dialogue dans l'ordre]
Les répliques de ce dialogue au marché sont mélangées. Numérote-les de 1 à 8 pour reconstituer la conversation, puis joue-la avec un camarade.

\begin{enumerate}
\item[\ldots] \eng{That's too much! What about 1,200 FCFA?}
\item[\ldots] \eng{Good morning, madam. Fresh plantains! Very sweet!}
\item[\ldots] \eng{Deal! Here is the money.}
\item[\ldots] \eng{Good morning. How much is this bunch of plantains?}
\item[\ldots] \eng{All right, madam. 1,300 FCFA, last price.}
\item[\ldots] \eng{It is 1,500 FCFA, madam.}
\item[\ldots] \eng{Thank you, madam. Come back next week!}
\item[\ldots] \eng{1,300 FCFA? OK, I'll take it.}
\end{enumerate}
\end{exercice}

\courssub{Je me prépare au Bac}

\begin{textebox}[Market day in Bafoussam]
Every Saturday, Mama Ngono wakes up at five o'clock and walks to the big market in Bafoussam with her two daughters. She sells avocados, bananas and vegetables from her farm. Her stall is small, but it is always full of customers because her prices are fair and her products are fresh.

Last Saturday, a young woman stopped at the stall. “How much are these avocados?” she asked. “Five hundred francs for three,” Mama Ngono answered with a smile. The customer thought it was a little expensive, so they negotiated for a minute. Finally, Mama Ngono agreed to sell four avocados for five hundred francs. The woman paid with a two-thousand-franc note, and Mama Ngono gave her the change. “Anything else, my dear?” she asked. The woman also bought a bunch of bananas before leaving, happy with her purchases.

For Mama Ngono, the market is not only a place to earn money. It is also a place where she meets friends, listens to the news of the town and teaches her daughters how to talk to customers politely. “A good trader,” she often says, “must be honest, patient and friendly.”
\end{textebox}

\begin{exercice}[Compréhension et langue — type Bac]
Lis le texte ci-dessus, puis réponds aux questions en anglais, avec tes propres mots.

\textbf{A. Comprehension questions}

\begin{enumerate}
\item Where does Mama Ngono sell her products?
\item Why is her stall always full of customers? Give two reasons.
\item How many avocados did the young woman finally buy, and for how much?
\item How did the customer pay, and what did Mama Ngono give her back?
\item According to Mama Ngono, what are the three qualities of a good trader?
\end{enumerate}

\textbf{B. Language work}

\begin{enumerate}
\item Find in the text the English for : \emph{une échoppe ; négocier ; la monnaie ; un régime de bananes}.
\item Complete with \eng{much} or \eng{many} : (a) How \ldots\ldots money did she earn? (b) How \ldots\ldots bananas did the woman buy?
\item Put this sentence into reported speech : \eng{“A good trader must be honest,” she often says.} $\rightarrow$ \eng{She often says that \ldots}
\end{enumerate}
\end{exercice}

\begin{textebox}[Texte à traduire]
Hier, ma mère m'a envoyé au marché pour acheter des provisions. J'ai demandé le prix des tomates à une vendeuse. Elles coûtaient mille francs le kilo, alors j'ai négocié et elle m'a fait un rabais. J'ai aussi acheté de l'huile et du sel. La vendeuse m'a rendu la monnaie et m'a dit : « Merci, à la prochaine fois ! » J'étais très content de mes achats.
\end{textebox}

\begin{exercice}[Traduction — type Bac]
Traduis le passage ci-dessus en anglais (environ 60 mots).
\end{exercice}

\begin{exercice}[Expression écrite et jeu de rôle — type Bac]
\textbf{Part A — Guided dialogue (about 80 words).}

You are at the Mokolo market in Yaoundé. You have 5,000 FCFA and you must buy : two kilos of rice, one litre of oil and some onions. Write the complete dialogue (8 to 10 exchanges) between you and the trader. Use the functional language of the lesson : greeting, asking prices, bargaining, paying, asking for the change, concluding politely.

\textbf{Part B — Reported paragraph (about 60 words).}

Now transform your dialogue into a short paragraph in the past tense, beginning like this :

\eng{Yesterday, I went to the Mokolo market. I bought \ldots}
\end{exercice}



\newpage

\coursec{Corrigés des exercices}

\begin{corrige}[Exercice 1 — QCM : les bonnes formules]
\begin{enumerate}
\item \textbf{b} — \eng{Excuse me, could you help me, please?} est la formule polie standard pour attirer l'attention d'un vendeur. Les formules a et c sont des ordres directs, très impolis en anglais.
\item \textbf{b} — Pour demander un prix, on utilise \eng{How much...?}, même avec un nom pluriel : \eng{How much are the tomatoes?}. \eng{How many} sert à demander une quantité (\eng{How many tomatoes do you want?}), pas un prix. La formule c est un calque du français.
\item \textbf{a} — \eng{Anything else?} est la formule idiomatique du vendeur pour proposer autre chose. Les formules b et c sont des traductions mot à mot du français et n'existent pas en anglais.
\item \textbf{b} — \eng{That's a bit expensive. Can you give me a discount?} est poli et adapté à la négociation. Attention : on dit \eng{too expensive} ou \eng{much too expensive}, jamais \eng{too much expensive}. La formule c est insultante.
\end{enumerate}
\end{corrige}

\begin{corrige}[Exercice 2 — Vrai ou faux ? Justifie]
\begin{enumerate}
\item \eng{True. “How much...?” is used with uncountable nouns such as rice, oil or money.}
\item \eng{False. “How many...?” is used with countable nouns; with “money” we say “How much money...?”.}
\item \eng{False. “Here is your change” is said by the seller when he gives the money back to the customer.}
\item \eng{True. “I'm just looking, thank you” politely tells the shop assistant that you do not need help.}
\item \eng{False. Bargaining is a normal and friendly practice in Cameroonian markets.}
\end{enumerate}
\end{corrige}

\begin{corrige}[Exercice 3 — Vocabulaire : au magasin et au marché]
\begin{enumerate}
\item \eng{trader} — \eng{The trader sells fruit and vegetables at the market.} (le commerçant, la vendeuse)
\item \eng{change} — \eng{...the shopkeeper gave me 1,500 FCFA change.} (la monnaie rendue)
\item \eng{bag} — \eng{Can I have a bag to carry my tomatoes, please?} (un sac, un sachet)
\item \eng{discount} — \eng{...I will give you a discount of 200 FCFA.} (un rabais, une remise)
\item \eng{stall} — \eng{...at a small stall near the bus station.} (une échoppe, un étal)
\item \eng{receipt} — \eng{Always ask for a receipt when you buy something in a supermarket.} (un reçu, un ticket de caisse)
\end{enumerate}
Les mots \eng{customer} (le client) et \eng{price} (le prix) n'étaient pas utilisés : c'étaient des distracteurs.
\end{corrige}

\begin{corrige}[Exercice 4 — Grammaire : much, many, some, any]
\begin{enumerate}
\item \eng{much} — \eng{How much is this bag of rice?} : on demande un prix, donc \eng{how much}.
\item \eng{many} — \eng{oranges} est un nom comptable au pluriel : \eng{How many oranges do you want?}.
\item \eng{any} ; \eng{some} — \eng{any} dans la question (\eng{Do you have any groundnut oil?}) ; \eng{some} dans la réponse affirmative (\eng{Yes, we have some}).
\item \eng{any} — phrase négative : \eng{We don't have any fresh fish today.}
\item \eng{some} — dans une demande polie (\eng{Could I have some bread, please?}), on utilise \eng{some} et non \eng{any}, car on attend une réponse positive.
\end{enumerate}
\textbf{Rappel} : \eng{much} + indénombrable ; \eng{many} + comptable pluriel ; \eng{some} en affirmatif et dans les offres/demandes polies ; \eng{any} en négatif et dans les questions neutres.
\end{corrige}

\begin{corrige}[Exercice 5 — Traduction : au marché de Mokolo]
\begin{enumerate}
\item \eng{How much is a kilo of tomatoes?} (ou \eng{How much does a kilo of tomatoes cost?})
\item \eng{That's too expensive. Can you reduce the price, please?} (ou \eng{Can you lower the price?})
\item \eng{I'll take two kilos of rice and one litre of oil.}
\item \eng{Would you like anything else, madam?} (ou \eng{Anything else, madam?})
\item \eng{Here is your change. Thank you and see you soon!}
\end{enumerate}
\textbf{Points de vigilance} : « combien coûte » se traduit par \eng{how much is...?} ou \eng{how much does... cost?}, jamais par \eng{how many} ; « je prends » (décision d'achat) se dit \eng{I'll take} ; « autre chose » se dit \eng{anything else}, pas \eng{other thing}.
\end{corrige}

\begin{corrige}[Exercice 6 — Réagir à des situations]
Réponses possibles (toute formule polie et correcte est acceptée) :
\begin{enumerate}
\item \eng{How much is this pineapple, please?}
\item \eng{I'm sorry, that's a bit expensive. Can you give me a discount?}
\item \eng{Here are 10,000 FCFA. Could I have my change, please?}
\item \eng{No, thank you. I'm just looking.}
\item \eng{Excuse me, do you sell soap?}
\end{enumerate}
\textbf{Points de vigilance} : \eng{pineapple} = ananas (faux ami : ce n'est ni une pomme ni un pin) ; \eng{I'm just looking} est la formule idiomatique pour « je regarde seulement » ; on dit \eng{do you sell...?} et non \eng{do you sale...?} (\eng{sale} est un nom, \eng{sell} est le verbe).
\end{corrige}

\begin{corrige}[Exercice 7 — Dialogue à trous]
\begin{enumerate}
\item \eng{Good morning, can I help you?} — le vendeur accueille le client ; le client répond \eng{Good morning. I would like some tomatoes, please.}
\item \eng{How much are they?} — le client demande le prix ; \eng{they} reprend \eng{tomatoes}, qui est pluriel (d'où \eng{are}).
\item \eng{Can you give me a discount?} — le client négocie après avoir trouvé le prix élevé (\eng{That is expensive!}).
\item \eng{Anything else?} — le vendeur propose autre chose ; le client répond \eng{No, thank you. That's all.}
\item \eng{Here is your change.} — le vendeur rend 600 FCFA sur les 2,000 FCFA et conclut la vente.
\end{enumerate}
Vérification du calcul : 2 kilos à 1,400 FCFA, paiement de 2,000 FCFA, monnaie = 600 FCFA. Le dialogue est ainsi cohérent du début à la fin.
\end{corrige}

\begin{corrige}[Exercice 8 — Remettre le dialogue dans l'ordre]
Ordre correct :
\begin{enumerate}
\item[1] \eng{Good morning, madam. Fresh plantains! Very sweet!}
\item[2] \eng{Good morning. How much is this bunch of plantains?}
\item[3] \eng{It is 1,500 FCFA, madam.}
\item[4] \eng{That's too much! What about 1,200 FCFA?}
\item[5] \eng{All right, madam. 1,300 FCFA, last price.}
\item[6] \eng{1,300 FCFA? OK, I'll take it.}
\item[7] \eng{Deal! Here is the money.}
\item[8] \eng{Thank you, madam. Come back next week!}
\end{enumerate}
Logique de la conversation : salutation et appel du vendeur $\rightarrow$ question sur le prix $\rightarrow$ annonce du prix $\rightarrow$ contre-proposition $\rightarrow$ dernier prix $\rightarrow$ acceptation $\rightarrow$ paiement $\rightarrow$ remerciement et conclusion. Les indices sont les connecteurs : \eng{It is 1,500 FCFA} répond à \eng{How much...?} ; \eng{1,300 FCFA? OK} reprend le \eng{last price} du vendeur ; \eng{Deal!} conclut la négociation.
\end{corrige}

\begin{corrige}[Exercice 9 — Compréhension et langue — type Bac]
\textbf{Barème indicatif (sur 20)} : partie A, 5 questions $\times$ 2 points = 10 points (1 point pour le contenu exact, 1 point pour la correction de la langue) ; partie B, 10 points (4 points pour le vocabulaire, 2 points pour \eng{much/many}, 4 points pour le discours rapporté).

\textbf{A. Comprehension questions} (réponses possibles, avec tes propres mots)
\begin{enumerate}
\item \eng{She sells her products at the big market in Bafoussam, at her small stall.}
\item \eng{Her stall is always full because her prices are fair and her products are fresh.}
\item \eng{She finally bought four avocados for five hundred francs.}
\item \eng{She paid with a two-thousand-franc note, and Mama Ngono gave her the change back.}
\item \eng{According to her, a good trader must be honest, patient and friendly.}
\end{enumerate}

\textbf{B. Language work}
\begin{enumerate}
\item \emph{une échoppe} = \eng{a stall} ; \emph{négocier} = \eng{to negotiate} ; \emph{la monnaie} = \eng{the change} ; \emph{un régime de bananes} = \eng{a bunch of bananas}.
\item (a) \eng{much} — \eng{money} est indénombrable ; (b) \eng{many} — \eng{bananas} est comptable au pluriel.
\item \eng{She often says that a good trader must be honest.} — le temps ne change pas car le verbe introducteur \eng{says} est au présent ; on supprime les guillemets et on introduit \eng{that}.
\end{enumerate}
\textbf{Points de vigilance} : répondre par des phrases complètes, pas par un seul mot ; ne pas recopier tout le texte mais reformuler ; \eng{a two-thousand-franc note} = un billet de deux mille francs.
\end{corrige}

\begin{corrige}[Exercice 10 — Traduction — type Bac]
\textbf{Barème indicatif (sur 10)} : respect du sens 4 points ; grammaire et temps 3 points ; vocabulaire 2 points ; orthographe et ponctuation 1 point.

Traduction possible :

\eng{Yesterday, my mother sent me to the market to buy some provisions. I asked a trader the price of the tomatoes. They cost one thousand francs a kilo, so I bargained and she gave me a discount. I also bought some oil and some salt. The trader gave me my change back and said: “Thank you, see you next time!” I was very happy with my purchases.}

\textbf{Points de vigilance} :
\begin{itemize}
\item \eng{sent} est le prétérit irrégulier de \eng{send} (\eng{send — sent — sent}) ;
\item tout le passage est au prétérit : \eng{asked}, \eng{cost} (prétérit identique à la base verbale), \eng{bargained}, \eng{gave}, \eng{said}, \eng{was} ;
\item « négocier » = \eng{to bargain} ou \eng{to negotiate} ; « faire un rabais » = \eng{to give a discount} ;
\item « rendre la monnaie » = \eng{to give somebody their change back} ;
\item « content de mes achats » = \eng{happy with my purchases} (préposition \eng{with}, pas \eng{of}) ;
\item « provisions » = \eng{provisions} ou \eng{foodstuff}, jamais \eng{provisions} orthographié autrement ni confondu avec \eng{providence}.
\end{itemize}
\end{corrige}

\begin{corrige}[Exercice 11 — Expression écrite et jeu de rôle — type Bac]
\textbf{Barème indicatif (sur 20)} : partie A sur 12 (respect de la tâche et des 8 à 10 échanges 4 points ; formules fonctionnelles de la leçon 4 points ; correction grammaticale et vocabulaire 4 points) ; partie B sur 8 (transformation au passé 4 points ; cohérence et langue 4 points).

\textbf{Part A — Modèle de dialogue}

\eng{Trader: Good morning, madam! Can I help you?}

\eng{Customer: Good morning. I would like two kilos of rice, please.}

\eng{Trader: Here you are. Anything else?}

\eng{Customer: Yes, one litre of oil and some onions. How much is everything?}

\eng{Trader: The rice is 2,000 FCFA, the oil is 1,500 FCFA and the onions are 500 FCFA. That is 4,000 FCFA.}

\eng{Customer: That's a bit expensive. Can you give me a discount?}

\eng{Trader: All right, for you: 3,500 FCFA, last price.}

\eng{Customer: Fine, I'll take it. Here are 5,000 FCFA.}

\eng{Trader: Thank you. Here is your change: 1,500 FCFA.}

\eng{Customer: Thank you very much. Goodbye!}

\eng{Trader: Goodbye, madam. Come back next week!}

\textbf{Part B — Modèle de paragraphe rapporté}

\eng{Yesterday, I went to the Mokolo market. I bought two kilos of rice, one litre of oil and some onions. The trader first asked for 4,000 FCFA, but I found it expensive, so I bargained and she gave me a discount. I finally paid 3,500 FCFA. I gave her a five-thousand-franc note and she gave me 1,500 FCFA change. I thanked her politely and went home, very happy with my purchases.}

\textbf{Points de vigilance} :
\begin{itemize}
\item dans le dialogue, respecter l'ordre logique : salutation $\rightarrow$ demande $\rightarrow$ prix $\rightarrow$ négociation $\rightarrow$ paiement $\rightarrow$ monnaie $\rightarrow$ conclusion polie ;
\item vérifier la cohérence des calculs (ici 3,500 FCFA payés sur 5,000 FCFA donnés = 1,500 FCFA de monnaie) ;
\item dans la partie B, passer tous les verbes au prétérit : \eng{went} (\eng{go}), \eng{bought} (\eng{buy}), \eng{asked}, \eng{found} (\eng{find}), \eng{paid} (\eng{pay}), \eng{gave} (\eng{give}) ;
\item \eng{I'll take it} devient au passé \eng{I took it} ou se reformule \eng{I finally paid...} ;
\item ne pas oublier la conclusion polie (\eng{Thank you}, \eng{Goodbye}) : elle fait partie des formules fonctionnelles évaluées.
\end{itemize}
\end{corrige}

\newpage

\coursec{Fiche bilan}

\begin{popfigure}
\begin{tikzpicture}[mindmap, grow cyclic, every node/.style={concept, font=\small},
root concept/.append style={concept color=popblue, font=\small\bfseries},
level 1/.append style={level distance=3.4cm, sibling angle=51.4},
level 2/.append style={level distance=2.2cm, font=\scriptsize}]
\node[root concept]{Buying and selling at the store / market}
  child[concept color=poporange]{ node{Setting the scene}
    child{ node{store, stall, customer, seller} }
    child{ node{foodstuff: rice, oil, tomatoes} }
  }
  child[concept color=popgreen]{ node{Asking for information}
    child{ node{How much is\ldots?} }
    child{ node{Do you have\ldots?} }
  }
  child[concept color=popgold]{ node{Giving information}
    child{ node{It costs\ldots} }
    child{ node{We sell it at\ldots} }
  }
  child[concept color=poppink]{ node{Bargaining}
    child{ node{That's too expensive!} }
    child{ node{Can you reduce the price?} }
  }
  child[concept color=poppurple]{ node{Relaunching \& concluding}
    child{ node{Anything else?} }
    child{ node{Here is your change.} }
  }
  child[concept color=popgreen!70!black]{ node{Pronunciation}
    child{ node{rising intonation: yes/no questions} }
    child{ node{stress on key words} }
  }
  child[concept color=popblue!60]{ node{Role-plays}
    child{ node{pair work, guided dialogues} }
  };
\end{tikzpicture}
\legende{Carte mentale de la leçon : échanger des informations pour acheter et vendre au magasin ou au marché.}
\end{popfigure}

\begin{bilan}
\textbf{1. Lexique clé (à connaître par cœur)}
\begin{itemize}
\item \cle{People and places} : a customer « un client », a seller / a shopkeeper « un vendeur / un commerçant », a stall « un étal », a nearby store « un magasin proche », the market « le marché ».
\item \cle{Foodstuff and articles} : rice « riz », palm oil « huile de palme », tomatoes « tomates », plantains « plantains », dried fish « poisson séché », soap « savon », a bag / a tin / a bunch « un sac / une boîte / un régime ».
\item \cle{Money words} : the price « le prix », change « la monnaie », a discount « une remise », cheap / expensive « bon marché / cher », to bargain « marchander », to cost « coûter ».
\end{itemize}

\textbf{2. Formules fonctionnelles et grammaire}
\begin{center}
\begin{tabularx}{\linewidth}{|X|X|}
\hline
\rowcolor{popblueL} Fonction & Formules en anglais \\
\hline
Demander une information & How much is this? / How much are the tomatoes? / Do you have any rice? \\
\hline
Donner une information & It costs 500 francs. / They are 200 francs a pile. / Yes, we have some. \\
\hline
Marchander & That's too expensive! / Can you reduce the price? / My last price is 300 francs. \\
\hline
Relancer & Anything else? / What else do you need? / Would you like a bag? \\
\hline
Conclure & I'll take it. / Here is your change. / Thank you, come again! \\
\hline
\end{tabularx}
\end{center}
\begin{itemize}
\item \cle{How much} + verbe \emph{to be} : singulier avec les noms indénombrables (\eng{How much is the oil?}), pluriel avec les noms dénombrables (\eng{How much are the plantains?}).
\item \cle{Some / any} : \emph{some} dans les phrases affirmatives et les offres, \emph{any} dans les questions et les négations (\eng{Do you have any sugar? --- No, we don't have any.}).
\item \cle{Present simple} pour les prix et les habitudes du marché : \eng{It costs 1\,000 francs.} « Cela coûte 1\,000 francs. »
\end{itemize}

\textbf{3. Méthodes express}
\begin{itemize}
\item Avant de parler : identifier son rôle (customer ou seller), lister 3 articles et leurs prix, préparer une question d'ouverture.
\item Pendant le dialogue : saluer, demander le prix, réagir (\eng{Really? That's too much!}), marchander, conclure poliment.
\item Prononciation : intonation montante pour les questions fermées (\eng{Do you have any rice?} « monte à la fin »), descendante pour les questions en \emph{wh-} (\eng{How much is it?} « descend à la fin ») ; accentuer les mots porteurs de sens (noms, chiffres).
\end{itemize}
\end{bilan}

\begin{aretenir}[Checklist avant le Bac]
\begin{itemize}
\item[$\square$] Je connais le vocabulaire des articles et des denrées alimentaires (foodstuff).
\item[$\square$] Je sais demander un prix avec \eng{How much is\ldots?} et \eng{How much are\ldots?}.
\item[$\square$] J'accorde correctement le verbe \emph{to be} (is / are) selon le nom.
\item[$\square$] J'utilise \emph{some} et \emph{any} sans erreur.
\item[$\square$] Je sais donner un prix : \eng{It costs\ldots} / \eng{They are\ldots a pile.}
\item[$\square$] Je maîtrise au moins trois formules pour marchander.
\item[$\square$] Je sais relancer la conversation : \eng{Anything else?}
\item[$\square$] Je sais conclure poliment un achat : \eng{Here is your change. Thank you!}
\item[$\square$] Je fais monter l'intonation dans les questions fermées et descendre dans les questions en \emph{wh-}.
\item[$\square$] J'accentue les mots importants (articles, prix) et non les petits mots.
\item[$\square$] Je suis capable de jouer un dialogue complet vendeur--client de 8 répliques sans notes.
\end{itemize}
\end{aretenir}

\findecours

\end{document}
